% Paper 1 of the axiom-omega series: pure mathematics.  Draft (2026-11-19).  Not for circulation.
% Bibliographic items that could not be checked against a source carry the comment % UNVERIFIED.
% AI-use statement: see the unnumbered section before the bibliography.
\documentclass[11pt,reqno]{amsart}
\usepackage[a4paper,margin=1.1in]{geometry}
\usepackage{amsmath,amssymb,amsthm,mathtools}
\usepackage{booktabs}
\usepackage{enumitem}
\usepackage[colorlinks=true,linkcolor=blue,citecolor=blue,urlcolor=blue]{hyperref}

\newtheorem{theorem}{Theorem}[section]
\newtheorem{proposition}[theorem]{Proposition}
\newtheorem{lemma}[theorem]{Lemma}
\newtheorem{corollary}[theorem]{Corollary}
\theoremstyle{definition}
\newtheorem{definition}[theorem]{Definition}
\theoremstyle{remark}
\newtheorem{remark}[theorem]{Remark}
\newtheorem{example}[theorem]{Example}
\numberwithin{equation}{section}

\newcommand{\R}{\mathbb{R}}
\newcommand{\C}{\mathbb{C}}
\newcommand{\HH}{\mathbb{H}}
\newcommand{\Q}{\mathbb{Q}}
\newcommand{\Z}{\mathbb{Z}}
\newcommand{\g}{\mathfrak{g}}
\newcommand{\z}{\mathfrak{z}}
\newcommand{\kk}{\mathfrak{k}}
\newcommand{\su}{\mathfrak{su}}
\newcommand{\cL}{\mathcal{L}}
\newcommand{\cR}{\mathcal{R}}
\newcommand{\Aut}{\operatorname{Aut}}
\newcommand{\Ad}{\operatorname{Ad}}
\newcommand{\ad}{\operatorname{ad}}
\newcommand{\re}{\operatorname{Re}}
\newcommand{\im}{\operatorname{Im}}
\newcommand{\rk}{\operatorname{rank}}
\newcommand{\ip}[2]{\langle #1,#2\rangle}
\newcommand{\hp}[3]{[#1,#2,#3]}
\newcommand{\SU}{\mathrm{SU}}
\newcommand{\SO}{\mathrm{SO}}
\newcommand{\OO}{\mathrm{O}}
\newcommand{\SL}{\mathrm{SL}}
\newcommand{\Cay}{\mathsf{C}}
\newcommand{\hol}{\operatorname{hol}}
% Parts.  Parts are numbered automatically, in order of appearance, by \paperpart; a part that has not yet
% been included simply does not appear, and every reference "Part~\ref{part:...}" follows.  When the algebraic
% Part 0 (expansion module m1) is included, put \setcounter{ppart}{-1} before its \paperpart so that it is Part 0.
\newcounter{ppart}
\renewcommand{\theppart}{\ifnum\value{ppart}=0 0\else\Roman{ppart}\fi}
\newcommand{\paperpart}[2]{\refstepcounter{ppart}\part*{Part \theppart. #2}\label{#1}}

\title[Ternions and minimal arcs]{Ternions and minimal arcs:\\ an origin-free calculus on the three-sphere}
% AUTHORS: to be supplied by the owner
\author{[byline to be confirmed]}
\date{19 November 2026}

\begin{document}
\raggedbottom

\begin{abstract}
A \emph{comparison structure} is a connected manifold with a ternary transport $\hp xyz$ obeying the
heap identities (a group with its identity forgotten) and a comparison $\kappa\colon M\times M\to[0,1]$ that is continuous,
symmetric, invariant under transports, equal to $1$ exactly on the diagonal, has exactly one zero partner for each point, and
satisfies the round-trip law $\kappa(w,\hp xwx)=(2\kappa(w,x)-1)^2$.  We prove that there are exactly two: the unit circle and the
unit quaternions $S^3$, with $\hp xyz=xy^{-1}z$ and $\kappa(x,y)=\lVert(x+y)/2\rVert^2$.  Beyond continuity of the heap operation and of
$\kappa$ nothing is assumed: dimension, compactness, smoothness and positive semidefiniteness are derived.  Associativity is not needed for the
geometry: if the heap identities are replaced by a binary operation with a two-sided unit at one point whose left and right
translations are continuous and preserve $\kappa$, with the round-trip law at that point, the space is still a round sphere $S^1$, $S^3$ or $S^7$ with
this $\kappa$, and the operations are then classified by Hurwitz's theorem.  Before any manifold, the chord--sagitta identity and an
elliptic/parabolic/hyperbolic trichotomy already hold for fair elements of a star algebra, and sparse comparisons on a chordal graph
with triangle defects at most $\varepsilon$ come from a potential up to $(|V|-2)\varepsilon$, the obstruction on other graphs being a
lattice gauge connection.  The key step concerns the duplication equation $f(x^2)=2f(x)^2-1$: on the line it has continuous non-cosine
solutions, none of them faithful, and on a connected Lie group every faithful continuous solution that is an even class
function is the cosine of an angle ($\re$ on $S^1$ and $\SU(2)$, the cosine of the rotation angle on $\SO(3)$).  Under a $C^2$ hypothesis, two further closing laws, minimal
dimension and completeness of comparison, select the same Lie algebra $\su(2)$; without any regularity hypothesis the
round-trip law can be replaced by an operational one: along a uniform transport, the comparison of any test point with the next point is an
affine function of its comparisons with the two previous points.  We describe the
geometry of the heap $S^3$: two flat connections with opposite torsion, their mean the Levi-Civita
connection, reflections as geodesic symmetries, and the torsion $3$-form as the obstruction to heaps on
other spheres.

On this basis we develop a calculus of minimal arcs in which differences are replaced by transports.
The recurrence $x_{n+1}=\hp{x_n}{x_{n-1}}{x_n}$ generates sampled geodesics exactly, in Chebyshev form
and over $\Q$; for fixed ends at distance $\Theta<\pi$ it is the unique minimiser of the discrete action $\sum(1-\cos d)$, and
its residual on a smooth curve is $h^2$ times the covariant acceleration.  Every step splits exactly into
a tangential rotation and a radial \emph{sagitta} $1-\cos d=2(1-\kappa)$; sums of sagittas vanish for
rectifiable paths and converge to half the quadratic variation for paths with a pathwise quadratic
variation, giving $3T/2$ and the It\^o drift for Brownian motion on $S^3$.  Levi-Civita transport along
a minimal arc is conjugation by the square root of the displacement; the product of these square roots
around a triangle, which has the algebraic form of the Bargmann--Pancharatnam invariant evaluated on the three vertices as unit vectors
of their span, has real part $(\sum\kappa-1)/(2\sqrt{\prod\kappa})=\cos(E/2)$, $E$ the spherical excess,
and the exact finite form of the lasso decomposition bounds the rotation angle of the Levi-Civita holonomy of a closed
polygon of minimal arcs by the total area of the geodesic triangles of any simplicial triangulated filling with
non-antipodal neighbours.  For rational vertices all transports and holonomies are
rational rotations.  A final section indicates which of these results are classical.
\end{abstract}

\maketitle

\section{Introduction}

\subsection{Objects}
The three-sphere $M=S^3$ of unit quaternions carries the ternary operation $\hp xyz=x\bar yz$.  Unlike
the group product, this operation does not refer to a unit element: it applies to $z$ the displacement
that carries $y$ to $x$.  A set with such an operation satisfying $\hp xxz=z=\hp zxx$ and
para-associativity is a \emph{heap} \cite{Prufer,Baer,Certaine}.  We call the points of the heap $S^3$ \emph{ternions}, the shortest
great-circle arc between two consecutive points of a sampled curve a \emph{minimal arc}, and
$v(x,y)=1-\cos d(x,y)$ the \emph{sagitta} of the step from $x$ to $y$ (Definition~\ref{def:arc}; it is
the height of the arc from $y$ through $x$ to the mirror image $\hp xyx$ above its chord).

\emph{Roadmap.}  Part~\ref{part:alg} (Section~\ref{sec:m1}) works before any manifold: in a star algebra a fair comparison already
has the chord--sagitta identity and an elliptic/parabolic/hyperbolic trichotomy by growth of powers, and on a graph sparse
comparisons are reconstructed from a potential on chordal graphs, with explicit error constants, while on other graphs the
obstruction is a lattice gauge connection (cycle holonomies modulo simultaneous conjugacy).  Part~\ref{part:axiom}
(Sections~\ref{sec:def}--\ref{sec:geom}) states the axiom and classifies its models: it
characterises the heap $S^3$ together with its comparison $\kappa(x,y)=\lVert (x+y)/2\rVert^2=\cos^2(d/2)$ by axioms that do not mention a
unit, describes the geometry of the heap, develops the theory without associativity (Mal'cev identities, Hurwitz's theorem, the
octonionic family), proves that two-sided invariance alone forces a round sphere (Section~\ref{sec:round}), and shows what one-sided
invariance gives and does not give (Section~\ref{sec:onesided}).  Part~\ref{part:arc} (Sections~\ref{sec:recur}--\ref{sec:rational}) develops a
calculus in which the difference $y-x$ of linear calculus is replaced by the transport $\hp y x\cdot$.  Part~\ref{part:final}
(Sections~\ref{sec:known}--\ref{sec:open}) places the results in the literature, records what is formalised, and lists the open problems.
% ROADMAP, temporary (delete each line when its part lands): Part "Beyond the axiom" (reflection spaces m2, indefinite and complex structures m3, self-reference d4) goes after Part I;
% phase synthesis and the rational atlas (m4) and the arc body and tropical limit (m5) join the arc calculus; Part "Arithmetic
% and holonomy" (m6) goes after it; the final part keeps known work, formalisation and open problems.
In the arc calculus the guiding observation is that linear calculus discards two second-order quantities
that are exact and computable in the heap: the radial sagitta of each step, and the failure of the
square roots of displacements to telescope.  The first is the quadratic-variation (It\^o) term, the
second is curvature.

\subsection{Results of Part~\ref{part:axiom}}
\noindent\textbf{Theorem A} (Theorem~\ref{thm:noreg}; Theorem~\ref{thm:main}).  Let $M$ be a connected manifold with a
continuous heap operation and $\kappa\colon M\times M\to[0,1]$ continuous, symmetric, faithful and invariant under
transports, with exactly one zero partner for each point and satisfying the round-trip law (axioms $(\Omega1)$--$(\Omega3)$ of
Definition~\ref{def:cs}).  Then $M$ is $S^1\subset\C$ or
$S^3\subset\HH$ and $\kappa(x,y)=\lVert(x+y)/2\rVert^2$.  No differentiability is assumed.  On the line the duplication formula
$f(2t)=2f(t)^2-1$ has continuous non-cosine solutions, but none of them is faithful; so every one-parameter subgroup closes,
on each closed one the solution is rigid, abelian subalgebras are lines and every $\ad X$ is elliptic, which leaves $\R$
and $\su(2)$.  A second proof under the additional hypothesis that $\kappa$ is $C^2$ near the diagonal
(Theorem~\ref{thm:main}) shows that all one-parameter subgroups close after the same length.

\noindent\textbf{Theorem C} (roundness without associativity; Section~\ref{sec:round}, Theorem~\ref{thm:round}, Corollary~\ref{cor:round}).  Let $M$ be a connected
locally Euclidean space of positive dimension without boundary (Hausdorff and second countable follow), $\kappa$ continuous, symmetric and equal to $1$ exactly on the diagonal, and $B$ a binary operation with a
two-sided unit $w$ whose left and right translations are continuous and preserve $\kappa$.  If $\kappa(w,B(x,x))=(2\kappa(w,x)-1)^2$ and $w$
has at most one zero partner, then $M$ is compact and $(M,\kappa)$ is the round $S^1$, $S^3$ or $S^7$ with $\kappa=\lVert(x+y)/2\rVert^2$.  The
proof is metric: the $\kappa$-preserving homeomorphisms form a Lie group (Hofmann--Kramer), the translations give Killing fields of
constant length whose integral curves are the geodesics, the duplication lemmas make the comparison a cosine of distance along each
geodesic, and an index-form estimate gives constant curvature.  For a comparison structure this is a second proof of the geometric
part of Theorem~A that uses no group structure on $M$.  Invariance under the right translations cannot be dropped as long as the
comparison is only continuous: a smooth twisted operation $xh_xyh_x^{-1}$ on $S^3$ with a continuous left-invariant comparison is not round
(Proposition~\ref{prop:onesided}; its comparison is not $C^2$ on the diagonal); under one-sided invariance roundness is equivalent to a
four-point Gram test (Proposition~\ref{prop:onesidedpos}).

\noindent\textbf{Weakening associativity} (Propositions~\ref{prop:weaknot}, \ref{prop:dim}, \ref{prop:S7}, \ref{prop:C}; Theorem~\ref{thm:rigid}).  If the
third heap identity is required only on subsystems generated by two elements, di-associativity is lost (a Fano-plane and
an $S^3$ example).  On $S^3$ with the standard comparison, continuity and invariance of $\kappa$ under left and right transports
restore the heap or its opposite, as an elementary consequence of Hurwitz's theorem; on $S^7$ the same hypotheses leave an
infinite-dimensional family, within which a single identity selects $(a\bar b)x$ up to isometry; and on a round sphere of any other
dimension than $0,1,3,7$ these hypotheses have no solution at all, even without continuity (Proposition~\ref{prop:dim}; summary table in Remark~\ref{rem:dim}).
By Theorem~C the round sphere is not an assumption: these results classify all operations with the Mal'cev identities,
continuous in the two outer arguments, and two-sided invariant $\kappa$ satisfying $(\Omega2\mathrm a$--$\mathrm d)$ and $(\Omega3)$ (Corollary~\ref{cor:round}).  With analyticity and
di-associativity instead, the Moufang identity is automatic once the tangent algebra is Malcev (Proposition~\ref{prop:mfauto}, from
theorems of Mal'cev and Kuz'min), and the comparison axioms leave $S^1$, $S^3$, $S^7$ (Proposition~\ref{prop:moufangclass}, now also a
consequence of Theorem~C).

\noindent\textbf{Theorem B} (Theorems~\ref{thm:parsimony}, \ref{thm:cc}, \ref{thm:operational}).  For pre-comparison
structures, which carry the $C^2$ hypothesis $(\Omega2\mathrm e)$, minimal dimension followed by minimal rank of $\kappa$ (the
dimension of the span of the functions $\kappa(x,\cdot)$), and completeness of comparison, select $\su(2)$ as well.  The round-trip law may be
replaced by \emph{two-readout determinacy} $(\Omega3')$: $\kappa(y,\hp xwx)$ is an affine function of $\kappa(y,x)$ and
$\kappa(y,w)$, with coefficients independent of $y$.  The proof reduces $(\Omega3')$ to d'Alembert's functional equation and uses no regularity.

\noindent\textbf{Geometry} (Section~\ref{sec:geom}).  The heap is equivalent to the pair of flat
connections $\nabla^L,\nabla^R$ of Cartan and Schouten; the reflection $s_x(w)=\hp xwx=2\ip xwx-w$ is the
geodesic symmetry; the torsion $3$-form obstructs heaps on all spheres other than $S^1,S^3$
(Theorem~\ref{thm:obstruction}).  Levi-Civita transport along a minimal arc from $u$ to $w$ is conjugation
by $(w\bar u)^{1/2}$ (Proposition~\ref{prop:LCarc}).

\subsection{Results of Part~\ref{part:arc}}
\noindent\textbf{Geodesic recurrence} (Theorems~\ref{thm:recur}, \ref{thm:action}, \ref{thm:consist}).
The recurrence $x_{n+1}=\hp{x_n}{x_{n-1}}{x_n}$ holds iff
$x_n=U_{n-1}(c)\,x_1-U_{n-2}(c)\,x_0$ for $n\ge1$, with $c=\ip{x_0}{x_1}$ and $U_k$ the Chebyshev polynomials of the
second kind, iff (for $x_1\ne\pm x_0$) the points are equally spaced on a great circle.  The sagitta action
$\sum_kv(x_k,x_{k+1})$ with fixed ends at distance $\Theta<\pi$ is at least $N(1-\cos(\Theta/N))$, with
equality only for this recurrence.  On a $C^4$ curve the residual of the recurrence is
$h^2\nabla_t\dot x-\tfrac{h^3}{2}(\lvert\dot x\rvert^2)'x+O(h^4)$.

\noindent\textbf{Sagittas and quadratic variation} (Proposition~\ref{prop:split},
Theorem~\ref{thm:sagitta}).  $y-x=\im(y\bar x)\,x-v(x,y)\,x$ exactly.  Sagitta sums tend to $0$ on
rectifiable and on H\"older-$\alpha$ paths with $\alpha>\tfrac12$, and to half the quadratic variation on
paths having one in F\"ollmer's sense (a reformulation of F\"ollmer's theory); for Brownian motion on $S^3$,
$\sum_kv_kx_{t_k}\to\tfrac32\int x\,dt$ in probability and almost surely along dyadic partitions, the It\^o drift, and
$2^{n/2}(V_{\pi_n}-\tfrac32T)\to N(0,\tfrac32T^2)$ by Jacod's theorem (Proposition~\ref{prop:clt}).

\noindent\textbf{Finite Stokes and holonomy} (Theorems~\ref{thm:stokes}, \ref{thm:triangle},
Corollary~\ref{cor:excess}).  On a shellable disk complex the boundary holonomy of any $\SU(2)$-valued edge cochain
equals the product, in shelling order, of conjugated face holonomies, at every mesh size; this is the exact finite form of
the classical lasso decomposition \cite{Bralic80,Levy03,Levy10}.  For the half-displacement cochain $h_{wu}=(w\bar u)^{1/2}$ of a
vertex map, each triangle contributes an element with the algebraic form of the Bargmann--Pancharatnam three-point invariant
\cite{SB88,MS93}, evaluated on $x,y,z$ as unit vectors of $\operatorname{span}(x,y,z)$
\[
a_{xyz}=\frac{(1+c_{xy}+c_{yz}+c_{zx})+\im(y\bar x)\times\im(z\bar x)}{4\sqrt{\kappa_{xy}\kappa_{yz}\kappa_{zx}}},
\qquad \re a_{xyz}=\cos\tfrac{E}{2},
\]
and the spin angle $\Theta=2\arccos\re H$ of the spin holonomy $H$ of a closed polygon of minimal arcs is at most the total area of the
geodesic triangles of any simplicial triangulated filling with non-antipodal neighbours (every simplicial
triangulation of a disk is shellable \cite{DK78}), hence so is the rotation angle $\min(\Theta,2\pi-\Theta)$ of the Levi-Civita holonomy.
The spin-angle bound is attained for every disk of area at most $2\pi$ in a great $2$-sphere (Theorem~\ref{thm:disk},
Corollary~\ref{cor:convex}); the rotation-angle bound is attained for such disks of area at most $\pi$.  Conversely, a $C^3$ filling of rank $2$ on a dense set with spin angle equal to its area lies in a great $2$-sphere
(Proposition~\ref{prop:converse}, via a comb form of the continuum Stokes formula, Lemma~\ref{lem:comb}).

\noindent\textbf{Rational calculus} (Proposition~\ref{prop:cayley}, Theorem~\ref{thm:rational}).  For
vertices in $S^3(\Q)$, the recurrence, all sagittas, all Levi-Civita transports along minimal arcs and all
polygon Levi-Civita holonomies (as rotations) are rational, while the spin holonomies are in general irrational; a two-chart Cayley atlas is bi-Lipschitz with constants $1$ and $2$.

Section~\ref{sec:known} lists which statements are classical.  Section~\ref{sec:lean} records the
formalisation status.

\subsection{Companion paper}
A companion paper (in preparation) studies physical readings of the
structure classified here \cite{Companion}.  It uses results from Sections~\ref{sec:def}--\ref{sec:stokes} of this paper.  Nothing in the
present paper depends on it, and we make no physical claims.

\subsection{Conventions}
$\HH$ has the Euclidean inner product $\ip xy=\re(\bar xy)$; $\im\HH\cong\R^3$ with cross product
$\times$, so that $\xi\eta=-\ip\xi\eta+\xi\times\eta$ and $[\xi,\eta]=2\,\xi\times\eta$ for
$\xi,\eta\in\im\HH$.  On $M=S^3$: $c(x,y)=\ip xy$, $d(x,y)=\arccos c(x,y)\in[0,\pi]$ (the round
distance), $\kappa=\tfrac12(1+c)$, and $\exp\xi=\cos\lvert\xi\rvert+\sin\lvert\xi\rvert\,\xi/\lvert\xi\rvert$,
so $d(1,\exp\xi)=\lvert\xi\rvert$ for $\lvert\xi\rvert\le\pi$.  Tangent vectors are written in the
\emph{right trivialisation} $T_uM=\{\xi u:\xi\in\im\HH\}$.  The rational points are $M(\Q)=S^3(\Q):=M\cap\Q^4$.

\setcounter{ppart}{-1}
\paperpart{part:alg}{Algebra and graphs before manifolds}

% Part 0 has the single Section 0, so that the numbers of Sections 2 onwards (cited by Paper II) are unchanged.
\setcounter{section}{-1}
\section{Fair comparisons in star algebras and on graphs}
\label{sec:m1}\label{m1:section}

Which comparison identities require a manifold, and when do sparse comparisons
come from a single choice of representatives? Two-sided fairness alone gives
the exact chord--sagitta identity in an associative star algebra. When the
symmetric part is a real scalar, the powers give a bounded, linear or
exponential growth criterion for the three comparison types. On a chordal
graph, triangle defects bounded by \(\varepsilon\) give a vertex potential with
uniform edge error at most \(\max\{|V|-2,0\}\varepsilon\). For arbitrary graphs
the gauge class of the edge data is determined by based cycle holonomies up
to simultaneous conjugacy. These results precede the manifold hypotheses of
Section~\ref{sec:def}.

\subsection{The chord and its sagitta}

\begin{proposition}\label{m1:algebra}
Let \(\mathcal A\) be a unital associative \(*\)-algebra with \(2\)
invertible, and call \(u\) \emph{fair} if \(u^*u=uu^*=1\).
For a fair element set \(S=(u+u^*)/2\), \(A=(u-u^*)/2\),
\(v=1-S\) and \(\Delta=u-1\). Then
\begin{equation}\label{m1:chord-eq}
 SA=AS,\qquad S^2-A^2=1,\qquad
 A^2=-v(2-v),\qquad \Delta^*\Delta=\Delta\Delta^*=2v .
\end{equation}
In particular \(S=0\) gives \(A^2=-1\).
\end{proposition}
\begin{proof}
Both \(u\) and \(u^*\) commute because each product is \(1\), so their
linear combinations \(S\) and \(A\) commute. Now
\((S+A)(S-A)=1\) gives \(A^2=S^2-1=(1-v)^2-1\).
Expansion of \((u^*-1)(u-1)\) and of the reversed product gives
\(2-u-u^*=2v\).
\end{proof}

For unit quaternions \(x,y\) take \(u=y\bar x\) and \(u^*=\bar u\).
Then \(S=c(x,y)\), \(A=\operatorname{Im}(y\bar x)\) and
\(y-x=\Delta x=Ax-vx\), with \(|y-x|^2=2v\). Thus
Proposition~\ref{prop:split} is the quaternionic instance of the two
algebraic components of the chord. Its tangent-projection interpretation
and the limiting laws of Theorem~\ref{thm:sagitta} also use the round
geometry.

\begin{example}\label{m1:algebra-boundaries}
Over \(\mathbb F_2\) with the identity involution, \(u=1\) is fair
and \(2S=u+u^*\) and \(2A=u-u^*\) both reduce to \(0=0\).
The undivided identity
\(\Delta^*\Delta=\Delta\Delta^*=2-u-u^*\) still holds. The scalar
condition on \(S\) in the growth criterion below is also substantive.
In \(M_2(\mathbb C)\) with the Hermitian adjoint,
\(u=\operatorname{diag}(1,i)\) is fair and has
\(A^2=\operatorname{diag}(0,-1)\). In \(\mathrm{Cl}_{4,0}\) with
reversion, put \(B=e_1e_2+e_3e_4\) for orthonormal generators \(e_i\).
The two summands commute and square to \(-1\), so
\(B^*=-B\) and \(B^2=-2+2e_1e_2e_3e_4\). This \(B\) cannot itself
be the odd part of a fair element in this algebra. Indeed left
multiplication on the Clifford basis is a faithful real matrix
representation with transpose as the adjoint, since each unit
generator acts by a symmetric signed permutation. The odd part of
an orthogonal matrix has operator norm at most \(1\), whereas \(B^*B\)
has eigenvalue \(4\) on the nonzero vector \(1-e_1e_2e_3e_4\). The product
\((1+e_1e_2)(1+e_3e_4)/2\) is fair and has odd part \(B/2\)
and even part \((1+e_1e_2e_3e_4)/2\).
\end{example}

\subsection{Repeated comparisons and positivity}

\begin{proposition}\label{m1:chebyshev}
For a fair \(u\) in a real or complex unital associative \(*\)-algebra,
let \(T_0=1,T_1=X,U_{-1}=0,U_0=1\) and use the recurrence
\(P_{k+1}=2XP_k-P_{k-1}\) for each family. For \(n\ge0\),
\begin{equation}\label{m1:cheb-eq}
 u^{\pm n}=T_n(S)\pm U_{n-1}(S)A,\qquad
 v(u^n)=1-T_n(1-v(u)).
\end{equation}
If the algebra is normed with continuous involution, the integer powers
are bounded if and only if both \(T_n(S)\) and \(U_{n-1}(S)A\) are
uniformly bounded for \(n\ge0\). Boundedness of nonnegative powers
suffices in this case. In a nonzero unital \(C^*\)-algebra every fair
element has \(\|u^n\|=1\) for all integers \(n\).
\end{proposition}
\begin{proof}
For the scalar Chebyshev formulas see \cite[\S18.5(i)]{m1:DLMF}.
The identity \(u+u^{-1}=2S\) gives
\(u^{\pm(n+1)}=2Su^{\pm n}-u^{\pm(n-1)}\) for \(n\ge1\).
The displayed expressions have the same recurrence and the same
initial values at \(n=0,1\). Averaging the two expressions gives the
sagitta formula and taking their half difference gives the odd term.
The triangle inequality in both directions of this decomposition gives
the boundedness criterion. Continuity of the involution gives a
constant \(C\) with \(\|u^{-n}\|=\|(u^n)^*\|\le C\|u^n\|\).
For a \(C^*\)-norm, \(\|u^n\|^2=\|(u^n)^*u^n\|=\|1\|=1\).
\end{proof}

\begin{theorem}[real scalar symmetric part]\label{m1:trichotomy}
Suppose that \(\mathcal A\ne0\) and \(S=s1\) for \(s\in\mathbb R\).
Write \(v=(1-s)1\). For \(A\ne0\), the algebra \(\mathbb R[A]\)
with involution has the following three types. The growth statements
hold for any norm on \(\operatorname{span}_{\mathbb R}\{1,A\}\), in
particular the restriction of a Banach \(*\)-algebra norm.
\begin{enumerate}
\item Elliptic, \(0<1-s<2\). The algebra is \(\mathbb C\) with
conjugation and the powers are bounded.
\item Parabolic, \(1-s\in\{0,2\}\). The algebra is the dual-number
algebra. Its powers have linear growth in \(|n|\).
\item Hyperbolic, \(1-s\notin[0,2]\). The algebra is the split
complex algebra. Its powers grow exponentially in \(|n|\), with rate
\(\rho=|s|+\sqrt{s^2-1}\) in the sense
\(\lim_{|n|\to\infty}\|u^n\|^{1/|n|}=\rho\).
\end{enumerate}
The scalar endpoints \(u=\pm1\) have bounded powers. In particular
bounded powers occur exactly for \(|s|<1\) or \(u=\pm1\).
\end{theorem}
\begin{proof}
Skew-adjointness makes \(1,A\) real-linearly independent when
\(A\ne0\), and \(A^2=(s^2-1)1\). Rescaling \(A\) by
\(\sqrt{|s^2-1|}\) when \(s^2\ne1\) gives the indicated quadratic
algebras, each with the generator sent to its negative by the
involution. For \(|s|<1\), put \(s=\cos\theta\), \(0<\theta<\pi\), and
\(J=A/\sin\theta\). The multiplication law for \(J^2=-1\) and the
angle addition formulas give \(u^n=\cos(n\theta)1+\sin(n\theta)J\)
for every integer \(n\).
For \(s=\pm1\), \(A^2=0\) gives \(u^n=s^n(1+nsA)\) for all integers
\(n\) by multiplication and inversion. Equivalence of norms on the
two-dimensional span bounds the norm above and below by positive
multiples of \(1+|n|\).

For \(|s|>1\), write \(b=\sqrt{s^2-1}\) and \(p_\pm=(1\pm A/b)/2\).
These nonzero idempotents have \(p_+p_-=0\) and \(p_++p_-=1\), so
\(u^n=(s+b)^n p_++(s-b)^n p_-\) for every integer \(n\).
The two coefficients have maximum modulus \(\rho^{|n|}\), because
\((s+b)(s-b)=1\). Norm equivalence in this basis gives fixed
positive lower and upper multiples of \(\rho^{|n|}\), whose
\(|n|\)-th roots tend to \(1\). Finally \(A=0\) implies \(s^2=1\).
\end{proof}

\begin{definition}\label{m1:positivity}
The involution is said to satisfy the scalar positivity condition
if \(a^*a=-t1\) with \(t\ge0\) implies \(a=0\).
\end{definition}

\begin{corollary}\label{m1:positive-powers}
Under the scalar hypothesis of Theorem~\ref{m1:trichotomy},
the implication in Definition~\ref{m1:positivity} for \(a=A\),
scalar positivity on \(\mathbb R[A]\), and bounded powers of \(u\)
are equivalent. A normalized faithful positive linear functional
on \(\mathcal A\) implies the ambient scalar positivity condition.
The condition also holds in every \(C^*\)-algebra.
\end{corollary}
\begin{proof}
If \(A\ne0\), the implication for \(A\) is equivalent to
\(1-s^2>0\) by \(A^*A=(1-s^2)1\). In that case
\((x+yA)^*(x+yA)=(x^2+(1-s^2)y^2)1\) with \(x,y\in\mathbb R\),
which proves the condition on the whole span. For \(A=0\)
the span consists of real scalars and its squares are nonnegative.
If a functional \(\omega\) is normalized and faithful, then
\(\omega(a^*a)=-t\ge0\) forces \(t=0\) and \(a=0\).
In a \(C^*\)-algebra the spectrum of \(a^*a\) lies in
\([0,\infty)\), so \(a^*a=-t1\) forces \(t=0\) and the norm
identity gives \(a=0\).
\end{proof}

\begin{example}\label{m1:power-boundaries}
The dual generator \(D^*=-D,\ D^2=0,\ D\ne0\) gives the fair
pair \(1\pm D\) with \(D^*D=0\). The split generator
\(J^*=-J,\ J^2=1\) gives the fair pair \(\cosh\eta\pm J\sinh\eta\),
\(\eta\ne0\), with \(J^*J=-1\). Without scalar \(S\), even a
singleton spectrum for \(S\) at \(1\) permits faster growth.
Use the real Banach algebra \(\mathbb R[t]/(t^3)\) with coefficient
\(\ell^1\) norm and \(t^*=-t\). Multiplication discards terms of
degree at least three, so the norm is submultiplicative and
\[
 u=1+t+\tfrac12 t^2,\quad u^*=1-t+\tfrac12t^2,\quad
 u^n=1+nt+\tfrac12 n^2t^2,\quad S=1+\tfrac12t^2.
\]
The power formula follows from multiplication for positive and
negative indices, starting with \(uu^*=1\). In the complexification,
\(S-\lambda1\) is invertible for \(\lambda\ne1\) by the finite
geometric series and nilpotent at \(\lambda=1\). The power norms are
\(1+|n|+n^2/2\).
\end{example}

\begin{proposition}\label{m1:time}
The real scalar comparison satisfies \(S(u^n)>0\) as a scalar
for all \(n\ge1\) exactly when \(s\ge1\).
For \(s=\cos\theta\), \(0<\theta<\pi\), the first strictly negative
reading occurs at \(n_0=\lfloor\pi/(2\theta)\rfloor+1\).
The simultaneous conditions of bounded powers and all positive
readings select \(u=1\). Positive readings and \(A^2\ne0\)
select \(s>1\).
\end{proposition}
\begin{proof}
The scalar reading is \(T_n(s)\). For \(s=\cosh\eta\ge1\) it
equals \(\cosh(n\eta)\) by the recurrence and hyperbolic addition,
and for \(s\le-1\) the first reading is negative. In the other
case it is \(\cos(n\theta)\). For \(0\le n<n_0\) the angle
is in \([0,\pi/2]\), whereas \(n_0\theta\) lies strictly between
\(\pi/2\) and \(3\pi/2\). Apply Theorem~\ref{m1:trichotomy}
and \(A^2=(s^2-1)1\) for the two selections.
\end{proof}

\begin{remark}\label{m1:matrix}
On \(M_2(\mathbb C)\) viewed as a real algebra the adjugate
\(a^\sharp=\operatorname{tr}(a)1-a\) is an involution.
Indeed it squares to the identity as a linear map, is the identity
on \(1\), and equals \(\det(a)a^{-1}\) on invertible matrices.
That expression proves reversal of products on invertible matrices
and the polynomial identity extends to all matrices by density.
The fair group is \(\mathrm{SL}(2,\mathbb C)\) by
\(aa^\sharp=a^\sharp a=\det(a)1\). Its complex scalar sagitta is
\(1-\operatorname{tr}(u)/2\). For \(u\ne\pm1\) the real values
in \((0,2)\), at \(\{0,2\}\), and outside \([0,2]\) give the three
types in Theorem~\ref{m1:trichotomy}. A nonreal value gives the
loxodromic case with nonreal trace. To see the projective
classification, the eigenvalues are \(\lambda,\lambda^{-1}\).
At a repeated eigenvalue the nonscalar Jordan block induces a
translation in suitable projective coordinates. Otherwise the
projective map is conjugate to \(z\mapsto\lambda^2 z\).
A real trace in \((-2,2)\) gives conjugate unit-modulus eigenvalues
and a real trace outside \([-2,2]\) gives real reciprocal
eigenvalues. Nonreal trace forces \(|\lambda|\ne1\) and nonreal
\(\lambda\). The convention calls positive real dilation hyperbolic and
includes a dilation with a nontrivial rotation among loxodromic maps.
There is no normalized positive functional for \(\sharp\) because
\(\operatorname{diag}(1,-1)^\sharp\operatorname{diag}(1,-1)=-1\).
Imposing Hermitian fairness as well gives \(\mathrm{SU}(2)\), whose
noncentral elements are elliptic by the unitary eigenvalues.
\end{remark}

\begin{remark}[rotations and boosts]\label{m1:spin}
On Hermitian \(2\times2\) matrices the action \(X\mapsto uXu^\dagger\)
of \(\mathrm{SL}(2,\mathbb C)\) preserves
\(\det(t1+\mathbf x\cdot\boldsymbol\sigma)=t^2-|\mathbf x|^2\),
where \(\boldsymbol\sigma\) are the Pauli matrices.
The action preserves the future cone of positive definite matrices,
and its real determinant is \(1\) by connectedness. Indeed the polar
decomposition \(u=kb\), with \(b=(u^\dagger u)^{1/2}\) positive of
determinant \(1\) and \(k\) special unitary, gives a path to \(k\),
which in turn is diagonalizable with entries \(e^{\pm i\theta}\)
and is joined to \(1\). For surjectivity onto the proper
orthochronous Lorentz group, such a transformation sends \(1\) to
some positive \(X\) of determinant \(1\). The matrix \(X^{1/2}\)
realizes that step. A remaining transformation fixing \(1\) is a
spatial rotation. For a unit vector \(\mathbf n\),
\(\exp(i\theta\,\mathbf n\cdot\boldsymbol\sigma)\) acts on its
orthogonal spatial plane by a rotation through angle of magnitude
\(2\theta\), by the Pauli multiplication rule. Since every spatial
rotation has an axis, these matrices realize the remaining step.
The kernel consists of the unitary matrices that commute with every
Hermitian matrix, hence the scalar matrices \(\pm1\).
At the identity the derivative sends a traceless complex matrix
\(Z\) to the map \(X\mapsto ZX+XZ^\dagger\). In its kernel,
\(X=1\) gives skew-adjointness, then every Hermitian \(X\)
commutes with \(Z\), so tracelessness forces \(Z=0\).
The real Lie groups both have dimension six, so the derivative
is an isomorphism. The inverse function theorem and translation by
the two kernel elements give the classical spin double cover. The group contains rotations, boosts
and null rotations. Diagonal matrices with entries \(e^{\pm i\theta}\)
rotate the spatial off-diagonal plane, those with entries
\(e^{\pm\eta/2}\) rescale the two null coordinates \(t\pm x_3\)
by \(e^{\pm\eta}\), and the nontrivial unipotent upper triangular
matrices induce unipotent maps fixing the null ray of
\(\operatorname{diag}(1,0)\). For the upper triangular matrix with
off-diagonal entry \(z\), write the diagonal entries of \(X\) as
\(a,b\) and its upper off-diagonal entry as \(w\). The new
entries are \(a+2\operatorname{Re}(z\bar w)+|z|^2 b\), \(b\)
and \(w+zb\), giving a unitriangular map in the order
\((b,\operatorname{Re}w,\operatorname{Im}w,a)\).
The quadratic-algebra generators are represented by
\(\operatorname{diag}(i,-i)\),
\(\bigl(\begin{smallmatrix}0&1\\0&0\end{smallmatrix}\bigr)\), and
\(\operatorname{diag}(1,-1)\), with squares \(-1,0,1\) respectively.
Indefinite fairness permits all these cases. Using the
positivity or time-order condition requires a choice of which
comparisons model states or coordinate changes.
\end{remark}

\subsection{Reconstructing sparse comparisons}

\begin{definition}\label{m1:gains}
Let \(G=(V,E)\) be a simple graph and \(H\) a group. A
\emph{gain assignment} chooses \(g_{ij}\in H\) on each oriented edge
with \(g_{ji}=g_{ij}^{-1}\). For a walk \(w=(x_0,\ldots,x_r)\)
write \(P_g(w)=g_{x_0x_1}\cdots g_{x_{r-1}x_r}\), with empty product
\(1\). A vertex potential is a map \(\phi:V\to H\) with exact edge
values \(\phi_i^{-1}\phi_j\). A graph is chordal if every simple
cycle of length at least four has an edge between nonconsecutive
vertices of the cycle.
\end{definition}

\begin{theorem}\label{m1:chordal}
Let \(G\) be finite chordal and \(d\) a bi-invariant pseudometric
on \(H\). If every triangle satisfies
\(d(g_{ij}g_{jk}g_{ki},1)\le\varepsilon\), with \(\varepsilon\ge0\),
there is a map \(\phi:V\to H\) such that
\begin{equation}\label{m1:stab-eq}
 d(g_{ij},\phi_i^{-1}\phi_j)
 \le \max\{|V|-2,0\}\varepsilon\qquad(\{i,j\}\in E).
\end{equation}
The same construction gives an exact potential in any group
when all triangle holonomies are \(1\).
\end{theorem}
\begin{proof}
Bi-invariance gives \(d(ab,c)=d(abc^{-1},1)\), conjugation
invariance, and \(d(a^{-1},b^{-1})=d(a,b)\) by left multiplication
by \(a\) and right multiplication by \(b\). Thus the triangle
bound is independent of its start and orientation and is equivalent
to \(d(g_{ij}g_{jk},g_{ik})\le\varepsilon\).

We give the graph argument for the simplicial-vertex step of
Dirac's theorem \cite{m1:Dirac}. A vertex is simplicial when its
neighbors form a clique. By induction a nonempty finite chordal graph
has a simplicial vertex, and a noncomplete one has two nonadjacent
such vertices. For a complete graph every vertex works. In a
disconnected graph take one from each of two components by induction.
In a connected noncomplete graph choose nonadjacent \(a,b\) and an
inclusion-minimal set \(D\subset V\setminus\{a,b\}\) separating
them. Every vertex of \(D\) has a neighbor in each of the components
\(X,Y\) of \(a,b\) in \(G-D\), since restoring that vertex alone
gives an \(a\)-to-\(b\) path. If \(p,q\in D\) were nonadjacent,
take shortest \(p\)-to-\(q\) paths with interiors in \(X\) and \(Y\).
Each has no chord by minimal length, the interiors have no edges
between them, and their union is a chordless cycle of length at
least four. Hence \(D\) is a clique. The proper induced graphs on
\(D\cup X\) and \(D\cup Y\) are chordal. If either is complete,
any vertex in its component outside \(D\) is simplicial there.
Otherwise induction supplies two nonadjacent simplicial vertices
and at least one lies outside the clique \(D\). The two selected
vertices have no additional neighbors in \(G\), so they are
nonadjacent simplicial vertices of \(G\).

Choose a simplicial vertex \(x\), solve on \(G-x\) by induction,
and keep the old values \(\psi_i\). If \(x\) has no neighbor, put
\(\phi_x=1\). Otherwise choose a neighbor \(w\) and put
\(\phi_x=\psi_w g_{wx}\). The edge \(xw\) is fitted exactly.
Every other neighbor \(j\) is adjacent to \(w\), and
\[
 d(g_{xj},\phi_x^{-1}\phi_j)
 \le \varepsilon+d(g_{wj},\psi_w^{-1}\psi_j).
\]
Here the first step replaces \(g_{xj}\) by \(g_{xw}g_{wj}\), and
left invariance gives the second term. Reverse edges have the same
error. The empty graph and graphs of at most two vertices admit
exact fits, and each later step adds at most \(\varepsilon\).
The exact assertion follows by the same induction with each
comparison an equality.
\end{proof}

For \(H\) the fair group of a unital \(C^*\)-algebra,
\(d(a,b)=\|a-b\|\) is bi-invariant. Left multiplication by a unitary
is an isometry by the \(C^*\)-identity and right multiplication is
an isometry by taking adjoints. A pseudometric suffices for
\eqref{m1:stab-eq}. At zero error exact reconstruction follows from
actual triangle equalities, or from zero distances when \(d\) is
a metric.

\begin{proposition}[sharpness of the reconstruction bound]\label{m1:sharp}
For \(n\ge3\) and \(\varepsilon>0\), the error
\((n-2)\varepsilon\) can be attained by the specified simplicial
reconstruction with a permissible choice of neighbors. For an
unconstrained choice of potential, the coefficient \(1\) at \(n=3\)
is optimal over bi-invariant metric groups. For \(n\ge4\) the
universal upper bound can already be reduced to
\((n-3)\varepsilon\). With \(H=(\mathbb R,+)\) and its usual metric,
the optimal universal coefficient at \(n\ge3\) is \((n-2)/n\).
\end{proposition}
\begin{proof}
Take the fan on \(0,\ldots,n-1\) with all spokes \(0j\) and the
path edges \(j,j+1\) for \(1\le j<n-1\). Give the outward spokes
gain \(0\) and the increasing path edges gain \(\varepsilon\).
Its only triangles are the fan faces, of defect \(\varepsilon\).
Eliminate \(n-1,n-2,\ldots,2\), each time using its predecessor
as the neighbor when reinserting it. Starting with the zero
potential on \(0,1\) gives \(\phi_j=(j-1)\varepsilon\)
for \(j\ge1\). The last spoke has error \((n-2)\varepsilon\).
For a triangle in the additive group \(\mathbb Z\) with metric
\(\varepsilon|a-b|\), put gain \(1\) on one oriented edge and \(0\)
on the other two. Error strictly less than \(\varepsilon\) would
force equality on all edges and is impossible because the oriented
sum is \(1\).

At four vertices, a connected chordal graph with a triangle has
a vertex adjacent to every other vertex. In fact the fourth vertex
is adjacent to some triangle vertex. Its star fits exactly and
every other edge then has error at most \(\varepsilon\). A triangle-free
chordal graph is a forest, where fitting tree edges successively
is exact. Disconnected four-vertex graphs have component size at
most three. Starting the induction of Theorem~\ref{m1:chordal}
at four vertices gives the stated improvement.

For real additive gains let \(r\) be the maximum, over simple
oriented cycles \(C\), of \(|\sum_C g|/|C|\), with maximum \(0\)
for a forest. Telescoping potentials gives the lower bound \(r\)
on uniform error. Give the arc \(ij\) weight \(g_{ij}+r\).
Every directed closed walk has nonnegative weight, because it
splits at repeated vertices into simple cycles and immediate
reversals, and the latter have weight \(2r\). Add a source
with zero-weight arcs to all vertices. The minimum path weight
\(\phi_j\) to each vertex exists by removing nonnegative closed
subwalks and minimizing over the finitely many simple paths.
It satisfies \(\phi_j\le\phi_i+g_{ij}+r\) for every arc, so both
orientations give \(|g_{ij}-(\phi_j-\phi_i)|\le r\).
A simple \(k\)-cycle in a chordal graph splits along chords into
\(k-2\) triangles, counted by induction using
\(k_1+k_2=k+2\) when a chord splits it. Thus
\(|\sum_C g|\le(k-2)\varepsilon\) and \(r\le(n-2)\varepsilon/n\).
The full outer cycle of the fan has sum \((n-2)\varepsilon\)
and length \(n\), giving equality for the optimal real error.
\end{proof}

\begin{proposition}[complete comparisons]\label{m1:complete}
Suppose \(X\ne\varnothing\) and \(c_{xy}\) takes values in the
fair group of \(\mathcal A\) for every pair, including repeated
points, with \(c_{yx}=c_{xy}^*\).
The third-point law \(c_{xy}c_{yz}=c_{xz}\) for every triple
is equivalent to \(c_{xy}=f_x^*f_y\) for fair \(f_x\).
It gives \(c_{xx}=1\), and anchoring at \(o\in X\) gives
\(f_x=c_{ox}\). Two representations differ by one common
left factor. The equivalence relation of identical comparison
profiles is equality of \(f_x\). Products along every finite
sampled path give \(c_{x_0x_N}\). These assertions hold also
for infinite \(X\).
\end{proposition}
\begin{proof}
An invertible idempotent \(c_{xx}\) must be \(1\), and using
the intermediate point \(o\) gives the representation. The converse
and the path assertion follow by telescoping. If \(f'\) represents
the same data, \(f_x'f_x^{-1}=f_o'f_o^{-1}\) follows from comparing
at \(o\). Equality of representatives gives equality of both profiles,
and equality of the profile at the single anchor recovers the
representative. For unit Clifford vectors the gain \(c_{xy}=xy\)
satisfies the law since \(y^2=1\).
\end{proof}

The third-point law gives the endpoint property of the heap in
Paper~I, equation~\eqref{eq:telescope}, on each finite partition with factors
multiplied in path order. Fairness on its own permits
\(c_{xx}=-1\) in the one-point real example, with \(c_{xx}^2=1\).
Thus the third-point law can supply self-pairing for the scalar
Gram input discussed below.

\subsection{Cycle holonomy and the gauge quotient}

\begin{definition}\label{m1:gauge-def}
A vertex gauge \(k:V\to H\) acts on gains by
\(g^k_{ij}=k_i^{-1}g_{ij}k_j\).
For a graph, write \(H^1(G;H)\) for the pointed set of all
gain assignments modulo these gauges, pointed by the identity
assignment. When triangles are filled in, \(H^1\) of the
resulting complex denotes the same quotient restricted to
assignments with trivial face holonomy.
\end{definition}

\begin{theorem}\label{m1:gauge}
For a finite graph with components \(G_a=(V_a,E_a)\), the gauge
quotient is
\[
 H^1(G;H)\ \cong\
 \prod_a H^{\,|E_a|-|V_a|+1}/H ,
\]
where each \(H\) acts by simultaneous conjugation. The coordinates
are fundamental cycle products transported to one root per
component. A potential exists exactly when every closed walk has
identity holonomy, and it suffices to test the simple chordless
cycles, including the triangles. On the triangle-flat locus the
obstruction is therefore carried by chordless cycles of length at
least four.
\end{theorem}
\begin{proof}
On a walk from \(i\) to \(j\), the gauge changes its product to
\(k_i^{-1}P_g(w)k_j\). For a rooted spanning tree let \(t_i\) be
the product on the tree path from the root to \(i\), with \(t_o=1\).
Gauge by \(k_i=t_i^{-1}\). The resulting tree values are \(1\)
and each directed non-tree edge \(ij\) has value
\(t_i g_{ij}t_j^{-1}\), the based fundamental cycle product.
A gauge preserving identity on all tree edges is constant on
the component. Conversely all choices of the non-tree values
occur on tree-normalized assignments. Counting the non-tree
edges proves the formula, including isolated vertices. The
identity class consists of potentials, since gauging the identity gives
\(k_i^{-1}k_j\), and gauging such an assignment by \(k_i^{-1}\)
on the vertices returns the identity.

Closed products telescope for a potential. To go from chordless
cycles to all closed walks, delete backtracks and split at repeated
vertices. A closed subwalk rooted away from the original root
contributes with conjugation by its initial access path. For a
simple cycle with a chord, start at one endpoint of the chord
and write the cycle as paths \(p:a\to b\) and \(q:b\to a\).
Its product factors as
\[
 P_g(p)P_g(q)=(P_g(p)g_{ba})(g_{ab}P_g(q)).
\]
Both factors are shorter simple cycle products. A different base
vertex changes the product by conjugation. Induction reduces to
chordless cycles. This also proves the potential criterion via the
tree normalization.
\end{proof}

\begin{corollary}[anchored abelian comparisons]\label{m1:anchored}
Let \(G\) be finite and connected and \(H\) abelian. Let \(R\) be a
nonempty set of vertices at which the gauges must be \(1\).
In a rooted spanning tree at \(o\in R\), the gauge
coordinates are the \(|E|-|V|+1\) based fundamental cycle
gains and the \(|R|-1\) rooted tree path gains to the other
anchors. The quotient is therefore isomorphic to \(H^{|E|+|R|-|V|}\).
\end{corollary}
\begin{proof}
The tree path gains between anchors are gauge invariant by the
endpoint law. So are the rooted loop gains since the root is fixed.
For two assignments with equal coordinates, write \(t_i,t_i'\)
for their tree path gains. The gauge \(k_i=t_i^{-1}t_i'\)
is \(1\) at the anchors, gives the new tree path gains \(t_i'\),
and preserves the based cycle gains. The tree normalization
in Theorem~\ref{m1:gauge} then forces all new edge gains to
be those of the second assignment. To realize prescribed
coordinates, set the \(t_i\) to the prescribed values at the
anchors and to \(1\) elsewhere. Use \(t_i^{-1}t_j\) on tree
edges, and \(t_i^{-1}h_{ij}t_j\) off the tree for the given
based cycle gains \(h_{ij}\). Reversal determines the other
orientation. The coordinate map respects multiplication
because \(H\) is abelian.
\end{proof}

\begin{proposition}\label{m1:clique}
Fill the triangles of \(G\), and let \(\Pi_a\) be the group
of rooted closed walks in each component with concatenation,
modulo backtrack moves and triangle moves \(i,j,k\leftrightarrow
i,k\). Then the triangle-flat gauge quotient is
\(\prod_a\operatorname{Hom}(\Pi_a,H)/H\).
Triangle flatness forces a potential for every coefficient group
if and only if every \(\Pi_a\) is trivial. If \(d\) is bi-invariant
and the triangle defects are bounded by \(\varepsilon\), a
closed walk reduced to a point by \(N\) triangle moves has
\(d(P_g(w),1)\le N\varepsilon\), with moves counted in either
direction.
\end{proposition}
\begin{proof}
Concatenation respects the moves and reversal gives inverses.
A triangle-flat assignment gives a homomorphism by taking the
product, and a gauge conjugates it at the root. Conversely, given
a homomorphism \(\rho:\Pi_a\to H\), take rooted tree paths
\(\gamma_i\) and set
\(g_{ij}=\rho([\gamma_i(ij)\gamma_j^{-1}])\).
The tree paths cancel on composition, so the assignment reverses
correctly and satisfies the triangle law. It is identity on tree
edges and realizes \(\rho\) on every rooted loop. The residual
gauge is common conjugation as above. Taking \(H=\Pi_a\) and
\(\rho=\mathrm{id}\) proves necessity for the assertion over
all groups. For the estimate, a backtrack leaves the product
unchanged. A triangle move replaces a subproduct by one at
distance at most \(\varepsilon\), also after multiplication
by the prefix and suffix. Sum these distances along the reduction.
\end{proof}

\begin{remark}[topological interpretation]\label{m1:topology}
The classical edge-path description identifies \(\Pi_a\) with
\(\pi_1\) of the clique complex \cite[Ch.~3]{m1:Spanier}.
Only its \(2\)-skeleton enters the relations. A chordal
component has contractible clique complex. At each simplicial
vertex with a neighbor the simplex on the vertex and its
neighbors retracts linearly onto the neighbor face by sending
the vertex to one chosen neighbor, keeping that face fixed.
Deletion preserves connectivity since paths crossing the vertex
can detour through its clique. Iteration ends at a point.
Abstract gain balance and switching are classical
\cite{m1:Zaslavsky}. The spanning-tree basis for reduced loops is
also the one in L\'evy's Lemma~1.3.33 \cite{Levy10}. Angular synchronization uses the same
relative-label model over the circle \cite{m1:Singer}.
Lerman and Shi characterize exact compact-group synchronization
by cycle consistency and use inconsistent cycles to estimate
corruptions \cite[\S3, Proposition~3]{m1:LermanShi}. The deterministic estimate
\eqref{m1:stab-eq} uses uniform triangle defects and chordality. Duncan and Kileel
extend the exact cycle-consistency criterion to hyperedge data
\cite[Theorems~3.1 and 3.5]{m1:DuncanKileel}.
\end{remark}

\begin{example}\label{m1:cycles}
On a four-cycle there are no triangles, so making three edge
values \(1\) and the fourth \(h\) realizes any holonomy \(h\).
For \(h\ne1\) there is no potential. Adding a hub gives an
induced outer four-cycle but makes triangle flatness sufficient
by taking the hub values as the potential, with hub value \(1\).
For a separate example, join two triangles at one vertex, with
no filled faces, and use \(H=S_3\).
Tree normalization can give based tuples \(((12),(12))\)
and \(((12),(23))\). They have the same separate conjugacy
classes on the simple cycles. Only the first tuple has equal
coordinates, so the two assignments lie in different gauge classes.
\end{example}

\begin{remark}[finite Stokes and gauge fields]\label{m1:stokes-link}
For a simple \(k\)-cycle in a chordal graph, the chord
factorization in the proof of Theorem~\ref{m1:gauge} gives a
boundary product of \(k-2\) triangle factors transported to the
base vertex. Applying bi-invariance gives the holonomy bound
\((k-2)\varepsilon\). It equals the fan's boundary defect in
Proposition~\ref{m1:sharp}. The finite Stokes theorem
(Theorem~\ref{thm:stokes}) gives the ordered face factorization
for shellable disks with a specified \(\mathrm{SU}(2)\)-cochain.
Its transport convention multiplies later edges on the left.
Putting its transport from \(i\) to \(j\) equal to \(g_{ij}^{-1}\)
gives its path holonomy \(P_g(w)^{-1}\).
In our convention transporting a face based at \(a\) to the base \(o\) along
\(\gamma:o\to a\) gives
\(P_g(\gamma)P_g(\partial f)P_g(\gamma)^{-1}\).
Lattice link variables and face products provide the
corresponding finite kinematics for the bundles of Paper~II
\cite[Sec.~\emph{Fields of comparison structures}, Theorem on Hypothesis F]{Companion}.
\end{remark}

\begin{remark}[continuum hypothesis]\label{m1:continuum-input}
To obtain the gauge fields of Paper~II, Hypothesis F supplies
a base manifold, a bundle of comparison structures, a clock
and a clock-preserving connection. The resulting structure groups
before and after the clock reduction are \(\mathrm{SO}(4)\)
and \(\mathrm U(2)\), respectively, in that model.
\end{remark}

\begin{remark}[continuum programme]\label{m1:continuum-programme}
A continuum reconstruction from sparse gain data would need a
refinement of the underlying graph, compatibilities or estimates
on links under subdivision, compactness modulo vertex gauge,
and a proof identifying normalized face holonomies with the
curvature of a limiting connection. A field equation would also
require a dynamical law and a convergence argument for that law.
The gauge quotient above specifies the finite data for such a
limit. A separate finite problem is to determine the optimal
stability constants for arbitrary bi-invariant metric groups at
each graph size beyond the cases where sharpness was proved above.
\end{remark}

\subsection{From potentials to Gram geometry}

\begin{proposition}\label{m1:gram}
Suppose \(c_{xy}=f_x^*f_y\) as in Proposition~\ref{m1:complete}.
Let \(\omega\) be a normalized positive linear functional with
\(\omega(a^*)=\overline{\omega(a)}\), or the corresponding real
condition for a real algebra. Then
\(G_{ij}=\omega(S(c_{x_ix_j}))\) is a real positive semidefinite
matrix with diagonal \(1\), and
\[
 d(x,y)^2=2\omega(v(c_{xy}))
          =\omega((f_y-f_x)^*(f_y-f_x))
\]
is the square of a Hilbert-space pseudodistance. If the reads
\(S(c_{x_ix_j})=s_{ij}1\) are real scalars, scalar positivity
alone gives positive semidefiniteness of \((s_{ij})\).
Conversely, every real positive semidefinite matrix with
diagonal \(1\) is an even readout of flat fair comparisons.
\end{proposition}
\begin{proof}
For \(\alpha_i\in\mathbb R\), the element \(a=\sum_i\alpha_i f_{x_i}\)
satisfies \(a^*a=\sum_{ij}\alpha_i\alpha_j S(c_{x_ix_j})\) since the
odd parts cancel in transposed pairs. Applying \(\omega\) gives
the Gram inequality. For scalar reads \(a^*a=(\alpha^{\mathsf T}
(s_{ij})\alpha)1\), and a strictly negative scalar would force
\(a=0\) by Definition~\ref{m1:positivity}, a contradiction.
On the real underlying vector space the symmetric bilinear form
\(\operatorname{Re}\omega(a^* b)\) is positive semidefinite.
Cauchy--Schwarz follows from nonnegativity of its quadratic on
\(a+tb\) for all real \(t\), so its nullspace is orthogonal to
everything. Quotient by that nullspace and complete to obtain a
real Hilbert space. Since \(f_y-f_x=f_x(c_{xy}-1)\),
Proposition~\ref{m1:algebra} gives the distance formula.
For the converse, factor the matrix as the Gram matrix of
real unit vectors \(f_i\) by diagonalization. In the Euclidean
Clifford algebra with reversion these vectors are fair and
\(S(f_i f_j)=(f_i f_j+f_j f_i)/2=\langle f_i,f_j\rangle1\).
The scalar-part functional is normalized and positive because
the scalar part of \(a^*a\) is the sum of squared coefficients
in the orthonormal Clifford monomial basis.
\end{proof}

Third-point redundancy, scalar even reads and scalar positivity
supply the normalized self-read, elliptic scalar square and Gram
positivity. A positive
functional supplies the Gram inequality without scalar even
reads. It gives a Hilbert norm after quotient and hence supplies
the parallelogram law on the linear realization by expanding
\(\|x+y\|^2+\|x-y\|^2=2\|x\|^2+2\|y\|^2\).
Conversely the Jordan--von Neumann theorem recovers the inner
product of a normed linear space with that law
\cite{m1:JordanVonNeumann}. In \(\ell^1_2\), the two sides for the
two coordinate vectors are \(8\) and \(4\).
A unit-sphere constraint concerns the induced geometry on this
realization, as in Theorem~\ref{thm:round} of the manifold theory.

\begin{proposition}\label{m1:distance}
Let \(D\) be a symmetric real matrix with zero diagonal, and put
\(J=I-\mathbf 1\mathbf 1^{\mathsf T}/m\) for \(m\ge1\).
The entries of \(D\) are squared Euclidean distances exactly when
\(-JDJ/2\) is positive semidefinite. For a complex Gram matrix the
full labeled entries determine any finite vector realization
up to a unitary isomorphism of the spans, which extends to an
ambient unitary when the ambient spaces have the same finite
dimension.
\end{proposition}
\begin{proof}
For a Gram matrix \(K\), the associated squared distance is
\(D_{ij}=K_{ii}+K_{jj}-2K_{ij}\). Then
\(-JDJ/2=JKJ\), the Gram matrix of centered vectors.
Conversely factor \(-JDJ/2\) as a real Gram matrix. Its induced
distance matrix \(D'\) has \(JD'J=JDJ\), and
\(J(D-D')J=0\) implies
\((D-D')_{ij}=r_i+r_j-r\), where \(r_i\) are row means and \(r\)
the overall mean. The zero diagonal forces \(2r_i=r\) for all
\(i\) and the difference vanishes. This is the finite
distance criterion of Schoenberg \cite{m1:Schoenberg}.
For the complex assertion, map
\(\sum_i a_i x_i\mapsto\sum_i a_i y_i\).
Squared norms of both linear combinations are \(a^*Ka\), so
the map is well-defined, onto the second span and isometric.
Polarization gives preservation of the complex inner product,
and a bijection of orthonormal bases of the complements gives
the extension. In particular, for the squared distances of a
real Gram matrix, \(a^{\mathsf T}Da=-2a^{\mathsf T}Ka\le0\)
whenever \(\sum_i a_i=0\).
\end{proof}

The chordal partial Gram completion theorem of Grone, Johnson,
S\'a and Wolkowicz \cite{m1:Grone} concerns consistent clique
submatrices with positive semidefinite Gram realizations.
Theorem~\ref{m1:chordal} concerns group-valued edge products
and their triangle defects. The shared graph condition permits
the common simplicial elimination step.

\begin{proposition}[two readout identities]\label{m1:readout-identities}
For real vectors \(x,y\) in a finite orthonormal coordinate
system,
\[
 \|x\|^2\|y\|^2-\langle x,y\rangle^2
 =\sum_{i<j}(x_i y_j-x_j y_i)^2 .
\]
A complex polynomial in \(z,\bar z\) invariant under common
phase is a polynomial in the quadratic entries \(\bar z_i z_j\).
If it is homogeneous of degree two and real-valued, it has
the form \(z^\dagger Mz\) for a unique Hermitian matrix \(M\),
and is nonnegative everywhere exactly when \(M\) is positive
semidefinite.
\end{proposition}
\begin{proof}
Expanding the squared wedges and combining the two orders
of the indices gives the left-hand side. For the phase
assertion, the change from real and imaginary variables to the
formal variables \(z,\bar z\) is complex-linearly invertible.
Polynomial identities on the real coordinate space therefore
identify coefficients in \(z,\bar z\). The monomial
\(z^\alpha\bar z^\beta\) acquires factor
\(e^{i\theta(|\alpha|-|\beta|)}\) under phase rotation.
Invariance at every \(\theta\) forces equal counts for every
monomial of nonzero coefficient. Pairing the factors gives a
product of the indicated quadratic entries. In degree two
these entries are independent monomials. Conjugating the value
swaps their indices and conjugates the coefficients, so reality
is precisely Hermiticity. The sign criterion is the definition
of a positive semidefinite Hermitian form.
\end{proof}

\begin{example}\label{m1:negative-gram}
Local ellipticity alone does not give a Gram matrix, even
with complete flat gains. In \(M_2(\mathbb R)\) with the
adjugate involution put
\[
 J_1=\begin{pmatrix}0&-1\\1&0\end{pmatrix},\qquad
 J_2=\begin{pmatrix}0&-2\\1/2&0\end{pmatrix},\qquad
 f_0=1,\quad f_1=(1+J_1)/\sqrt2,\quad f_2=(1-J_2)/\sqrt2 .
\]
Both \(J_k\) square to \(-1\) and are odd, so the \(f_i\) are
fair. Half-traces of \(f_i^\sharp f_j\) on edges \(01,02,12\)
are \(1/\sqrt2,1/\sqrt2,-1/8\), the last because
\(J_1J_2=\operatorname{diag}(-1/2,-2)\). The determinant
of the unit-diagonal read table is \(-9/64\) by expansion.
The group generated by the comparisons is unbounded since
it contains \(f_1^2 f_2^2=-J_1J_2=\operatorname{diag}(1/2,2)\).

Conversely, positivity in the ambient algebra leaves room
for a nonflat comparison. Take complex circle gains
\(c_{01}=c_{02}=e^{i\pi/6}\), \(c_{12}=-1\), inverse gains on
reversed edges and \(c_{ii}=1\). The triangle product is
\(-1\), and replacing \(c_{12}\) by the product through \(0\)
has chord error \(2\). The real reads on \(01,02,12\) are
\(\sqrt3/2,\sqrt3/2,-1\). The table has eigenvalue \(2\)
on \((0,1,-1)\). The complementary block in the orthonormal
basis \((1,0,0),(0,1,1)/\sqrt2\) has eigenvalues
\((1\pm\sqrt7)/2\). The determinant is \(-3\). The two
chords through \(0\) have length sum
\(2\sqrt{2-\sqrt3}=\sqrt6-\sqrt2<2=d(1,2)\).
\end{example}

\begin{proposition}\label{m1:near-gram}
For \(m\ge2\) labeled points with normalized diagonal and complete
fair gains in a unital \(C^*\)-algebra with a state \(\omega\),
fix one of the labels \(o\) and let
\(\delta_o=\max_{ij}\|c_{ij}-c_{io}c_{oj}\|\).
Then the real read matrix satisfies
\(\lambda_{\min}(G)\ge -(m-2)\delta_o\).
For an external anchor the bound is \(-(m-1)\delta_o\).
The single-label matrix is \((1)\).
\end{proposition}
\begin{proof}
Subtract from \(G\) the positive semidefinite Gram matrix of
the anchored values \(f_i=c_{oi}\) in Proposition~\ref{m1:gram}.
The symmetric error matrix \(F\) has entries bounded by
\(\delta_o\), by \(|\omega(a)|\le\|a\|\), and zero diagonal.
With an internal anchor its row and column vanish. For a
real vector \(z\), apply \(2|z_i z_j|\le z_i^2+z_j^2\) to
\(\sum_{ij}|F_{ij}||z_i z_j|\).
This bounds the quadratic form by the largest absolute row sum
times \(\|z\|^2\), proving both bounds.
\end{proof}

\begin{remark}\label{m1:unitarize}
There is a stronger sufficient input for a family of comparisons
represented on a finite-dimensional real or complex vector space.
Assume the subgroup \(\Gamma\) generated by the fair gains is bounded. Its matrix closure is a compact group. In fact the
inverses lie in the same bounded set, so for a convergent sequence
a subsequence of inverses converges to an inverse of the limit.
Average a positive definite inner product over that compact
group with normalized Haar measure to get a positive definite
invariant form. Each \(g^*\) for \(g\in\Gamma\) is then the
adjoint of \(g\) for this form. By real linearity the involution
agrees with that adjoint on the real span of \(\Gamma\).
Thus scalar positivity holds on the span, and for flat
fair gains with real scalar even reads Proposition~\ref{m1:gram}
applies. A faithful operator realization is understood here.
\end{remark}

\begin{example}[spherical excess]\label{m1:spherical}
Let \(a=e_1\), \(b=\cos\epsilon\,e_1+\sin\epsilon\,e_2\) and
\(p=\cos\epsilon\,e_1+\sin\epsilon\,e_3\) on the unit \(2\)-sphere,
\(0<\epsilon\le\pi/2\). The Gram determinants of
\((a,b,e_3)\) and \((a,b,p)\) are \(\sin^2\epsilon\) and
\(\sin^4\epsilon\), by squaring the corresponding column
determinants. In oriented dimension three the Hodge star
identifies unit plane bivectors with normal directions. In dimension two the star of a \(2\)-form is a
scalar and two independent directions have no common nonzero
normal in the plane. In dimensions above three the star of a
bivector has degree greater than one.
The angle at \(a\) is \(\pi/2\). Writing \(c=\cos\epsilon\),
the cosine of the angle \(\beta\) at each of \(b,p\) is
\(c/\sqrt{1+c^2}\) by projecting the other two vertices
orthogonally to the tangent plane. Hence
\(\tan(\beta-\pi/4)=(1-c)/(1+c)=\tan^2(\epsilon/2)\),
with the endpoint \(c=0\) obtained directly at \(\beta=\pi/2\).
Girard's theorem gives area equal to the spherical excess,
here
\[
 \Omega=\pi/2+2\beta-\pi=2\arctan(\tan^2(\epsilon/2))
        =\epsilon^2/2+O(\epsilon^4)>0 .
\]
Taylor expansion at \(0\) gives the leading term. The exact
formula is compatible with Theorem~\ref{thm:disk}.
The third side has cosine \(c^2\), so equality of the three
side lengths in this range occurs only at \(\epsilon=\pi/2\).
An affine Euclidean triangle has zero excess.
\end{example}

\begin{remark}[Gram phases and order]\label{m1:phase}
For unit vectors \(\psi_i\) in a complex Hilbert space use
\(K_{ij}=\langle\psi_i,\psi_j\rangle\), antilinear in the first
slot. On the graph of nonzero overlaps the gains
\(g_{ij}=K_{ij}/|K_{ij}|\) reverse by conjugation. Rephasing
the vectors gives the vertex gauge law. For
\(P_i=\psi_i\psi_i^\dagger\) the trace of \(P_1P_2P_3\) is
\(K_{12}K_{23}K_{31}\) by rank-one multiplication. Its phase
is the triangle holonomy in our product order. For the octant
\(e_3,e_1,e_2\) use
\((1,0)^{\mathsf T},(1,1)^{\mathsf T}/\sqrt2,
(1,i)^{\mathsf T}/\sqrt2\), whose overlap product is \((1+i)/4\).
Replacing \(e_2\) by \(-e_2\) conjugates this value.
The phase of the returned state for horizontal transport is the
inverse phase, as in Paper~II's proposition on the Hopf fibre and
its holonomy \cite{Companion,MS93}. The octant has solid angle
\(\pi/2\). The three Bloch overlap magnitudes are the same for
both orientations, so squared moduli alone lose this datum.
Full Gram reconstruction concerns reindexing or reading the
table. For noncommuting transports, consider the quarter-turns
\(R_x,R_y\) about the positive \(x,y\) axes give \(R_xR_y e_3=e_1\) and
\(R_yR_x e_3=-e_2\) by the right-hand rule.
Also \(P_0 P_+\ne P_+P_0\) for the rank-one projections onto
\((1,0)^{\mathsf T}\) and \((1,1)^{\mathsf T}/\sqrt2\).
\end{remark}

\begin{example}[overlap data and contextuality]\label{m1:overlap-boundaries}
For two unit vectors, an inner product of \(1\) is equivalent
to equality, since their squared difference is
\(2-2\operatorname{Re}K_{ij}\). For
\(\psi(z)=(\sqrt{1-t},\sqrt t\,z)\), \(0\le t\le1\),
\(|z|=1\), the overlap with \(\psi(w)\) is
\(1-t+t\bar z w\). The endpoints remain valid by direct
substitution.
Global Gram realization is an issue about one vector family.
An obstruction to assigning common signs to overlapping
commuting operator contexts can occur even with legitimate
matrices, as in the Peres--Mermin square \cite{m1:Mermin}.
With \(X,Y,Z\) the Pauli matrices and \(I\) the
identity, take
\[
 \begin{pmatrix}
 X\otimes I&I\otimes X&X\otimes X\\
 I\otimes Y&Y\otimes I&Y\otimes Y\\
 X\otimes Y&Y\otimes X&Z\otimes Z
 \end{pmatrix}.
\]
The rows and columns are commuting triples. The product in
each row and in the first two columns is \(I\otimes I\), and
in the last column it is \(-I\otimes I\). Indeed the first
two entries multiply to the third except in the last column,
where \(XY=iZ\) gives its negative. Since the first two entries
commute and square to identity, the third commutes with them.
An assignment of signs with these products would multiply the
six context products to \(-1\), whereas each assigned sign
occurs twice and the product is \(1\).
\end{example}

\begin{example}[sequential instruments]\label{m1:sequential}
Write \(P_1=I-P_0\) and \(X=\bigl(\begin{smallmatrix}0&1\\1&0
\end{smallmatrix}\bigr)\). The two Kraus families \(P_b\) and
\(XP_b\), \(b=0,1\), have the same effects \(P_b\) because \(X\)
is unitary. With initial state \(P_0\), branch \(0\) is \(P_0\)
in the first instrument and \(P_1\) in the second. The second
effect \(P_1\) therefore gives joint weights
\(\operatorname{Tr}(P_0P_1)=0\) and
\(\operatorname{Tr}(P_1^2)=1\).
\end{example}

\begin{remark}[how Part~0 connects to the rest]\label{m1:links}
(i) Proposition~\ref{m1:algebra} is the algebraic content of the sagitta: for $u=y\bar x$ on $S^3$ it gives the exact split of
Proposition~\ref{prop:split} (where $v(x,y)=1-\re(y\bar x)$), and the sagitta sums of Theorem~\ref{thm:sagitta} add up the quantity $v=1-S$ of
\eqref{m1:chord-eq}; the Chebyshev form \eqref{m1:cheb-eq} is the algebraic form of the geodesic recurrence of
Theorem~\ref{thm:recur}.  (ii) The cycle holonomies of Theorem~\ref{m1:gauge} and the chord factorisation of
Remark~\ref{m1:stokes-link} are the graph form of the finite Stokes theorem (Theorem~\ref{thm:stokes}) and of its continuum comb
order (Lemma~\ref{lem:comb}); an $\SU(2)$-cochain of Section~\ref{sec:stokes} is a gain assignment (Definition~\ref{m1:gains}) on the
$1$-skeleton, with the inversion convention of Remark~\ref{m1:stokes-link}.  (iii) The Gram statements of Propositions~\ref{m1:gram}, \ref{m1:distance} and \ref{m1:near-gram} are
the finite counterparts of the positive semidefiniteness and rank in Theorems~\ref{thm:main} and \ref{thm:round}, and of the four-point
Gram test of Proposition~\ref{prop:onesidedpos}(ii); Example~\ref{m1:negative-gram} shows why such a test is needed as an input.
\end{remark}

\setcounter{section}{1}

\paperpart{part:axiom}{The axiom and its classification}

\section{Comparison structures}\label{sec:def}

We first state the transport axiom alone, because it already determines a smooth structure.

\smallskip\noindent$(\Omega1)$ \emph{Origin-free transport.}  $M$ is a connected topological manifold of
finite dimension $n\ge1$ and $\hp\cdot\cdot\cdot\colon M^3\to M$ is continuous with
\[
\hp xxz=z,\qquad \hp xzz=x,\qquad \hp{\hp xyz}uv=\hp xy{\hp zuv}.
\]

\begin{lemma}[Certaine \cite{Certaine}]\label{lem:certaine}
Assume $(\Omega1)$ and fix $o\in M$.  The product $x\cdot y:=\hp xoy$ makes $M$ a topological group $G_o$
with unit $o$ and inverse $x^{-1}=\hp oxo$, and $\hp xyz=xy^{-1}z$.  For $o'\in M$ the map
$x\mapsto\hp xo{o'}$ is an isomorphism $G_o\to G_{o'}$.
\end{lemma}
\begin{proof}
Associativity is the third identity with $y=u=o$.  Unit: $\hp ooy=y$, $\hp xoo=x$.  Inverse:
$\hp xo{\hp oxo}=\hp{\hp xoo}xo=\hp xxo=o$ and $\hp{\hp oxo}ox=\hp ox{\hp oox}=\hp oxx=o$.  Then
$xy^{-1}z=\hp xo{\hp{\hp oyo}oz}=\hp xo{\hp oyz}=\hp{\hp xoo}yz=\hp xyz$.  Finally, for
$\phi(x)=\hp xo{o'}$: $\phi(\hp xoy)=\hp{\hp xoy}o{o'}=\hp xo{\phi(y)}$ and
$\hp{\phi(x)}{o'}{\phi(y)}=\hp xo{\hp{o'}{o'}{\phi(y)}}=\hp xo{\phi(y)}$, so $\phi$ is a homomorphism
$G_o\to G_{o'}$; it is continuous with continuous inverse $x\mapsto\hp x{o'}o$.
\end{proof}

\begin{lemma}\label{lem:lie}
Assume $(\Omega1)$.  There is a unique real-analytic manifold structure on $M$, compatible with its
topology, for which every $G_o$ is a Lie group; the ternary operation is analytic for it, and so are all
left transports $\ell_{a,b}:=\hp ab\cdot$ and right transports $r_{a,b}:=\hp\cdot ab$.
\end{lemma}
\begin{proof}
A locally Euclidean topological group admits a unique compatible analytic structure making it a Lie
group \cite{Gleason52,MZ}; let $\mathcal A_o$ be that of $G_o$.  The map $\phi(x)=\hp xo{o'}$ is right
multiplication by $o'$ in $G_o$, hence an analytic diffeomorphism of $(M,\mathcal A_o)$, and it is a
continuous isomorphism $G_o\to G_{o'}$, hence analytic $(M,\mathcal A_o)\to(M,\mathcal A_{o'})$.  So
$\mathrm{id}_M=\phi\circ\phi^{-1}$ is analytic from $\mathcal A_o$ to $\mathcal A_{o'}$, and by symmetry
$\mathcal A_o=\mathcal A_{o'}$.  Analyticity of $\hp xyz=xy^{-1}z$ and of the transports follows.
\end{proof}

\begin{definition}\label{def:cs}
A \emph{comparison structure} is a triple $(M,\hp\cdot\cdot\cdot,\kappa)$ satisfying $(\Omega1)$ and
\begin{enumerate}[label=$(\Omega\arabic*)$,leftmargin=3.2em,start=2]
\item \emph{Faithful invariant comparison.}  $\kappa\colon M\times M\to[0,1]$ is continuous and
\begin{enumerate}[label=(\alph*)]
\item $\kappa(x,y)=\kappa(y,x)$;
\item $\kappa(x,y)=1$ if and only if $x=y$;
\item $\kappa(\hp xuv,\hp yuv)=\kappa(\hp uvx,\hp uvy)=\kappa(x,y)$;
\item for every $x$ there is exactly one $x^\perp$ with $\kappa(x,x^\perp)=0$.
\end{enumerate}
\item \emph{Round-trip law.}  $\kappa\bigl(w,\hp xwx\bigr)=\bigl(2\kappa(w,x)-1\bigr)^2$ for all $w,x\in M$.
\end{enumerate}
An \emph{isomorphism} is a homeomorphism preserving the ternary operation and $\kappa$; the automorphisms of $M$ form the
group $\Aut(M)$.
\end{definition}

We also use the following regularity condition, which is \emph{not} part of the definition:

\smallskip\noindent$(\Omega2\mathrm e)$ \emph{for some $x_*\in M$ the function $\kappa(x_*,\cdot)$ is of class $C^2$ on a neighbourhood of
$x_*$, for the analytic structure of Lemma~\ref{lem:lie}.}

\smallskip\noindent It enters the first proof of Theorem~\ref{thm:main} (Section~\ref{sec:class}, the $C^2$ route) and the definition of
pre-comparison structures (Section~\ref{sec:closing}); Theorem~\ref{thm:noreg} shows that every comparison structure satisfies it.  (An earlier version
listed it as clause (e) of $(\Omega2)$.)  By Lemma~\ref{lem:lie}, $(\Omega2\mathrm e)$ refers to a structure defined by $(\Omega1)$ alone.  Since transports are analytic diffeomorphisms preserving $\kappa$ and act
transitively (for $\ell=\ell_{x,x_*}$, $\kappa(x,\cdot)=\kappa(x_*,\cdot)\circ\ell^{-1}$), $(\Omega2\mathrm e)$ then
holds at every point.  The proofs below use $(\Omega2\mathrm e)$ only through the second-order (Peano) expansion of
$\kappa(x_*,\cdot)$ at $x_*$; twice Fr\'echet differentiability at the single point $x_*$ would suffice.  Positive semidefiniteness of $\kappa$ is
not assumed; it is derived (Theorem~\ref{thm:main}).  Write $s_x(w):=\hp xwx$ for the \emph{reflection} of $w$ in $x$ and
$c:=2\kappa-1$ for the \emph{correlation}.  By $(\Omega2\mathrm c)$,
$\kappa(x,s_xw)=\kappa(\hp xww,\hp xwx)=\kappa(w,x)$, so $(\Omega3)$ says
$\kappa(w,s_xw)=c(w,x)\,c(x,s_xw)$: the comparison of a point with its mirror image factorises through
the mirror.

\begin{example}\label{ex:models}
Let $A=\C$ or $\HH$, $M=\{x\in A:\lvert x\rvert=1\}$, $\hp xyz=x\bar yz$ and
$\kappa(x,y)=\tfrac12(1+\re\bar xy)=\lVert\tfrac12(x+y)\rVert^2=1-\tfrac14\lVert x-y\rVert^2$.  Laws
$(\Omega1)$, $(\Omega2)$ are immediate, with $x^\perp=-x$.  Since $x\bar w+w\bar x=2\ip xw$, one has
$x\bar wx=2\ip xw x-w$, hence $\ip w{s_xw}=2\ip xw^2-1$ and $\kappa(w,s_xw)=\ip xw^2=c(w,x)^2$.
\end{example}

\section{Classification}\label{sec:class}

\begin{theorem}\label{thm:main}
Every comparison structure is isomorphic to exactly one of the two structures of
Example~\ref{ex:models}.  In particular $\dim M\in\{1,3\}$, $M$ is compact, $\kappa$ is positive
semidefinite of rank $\dim M+2$, and the transport is noncommutative ($\hp xyz\ne\hp zyx$ for some
$x,y,z$) if and only if $M\cong S^3$.
\end{theorem}

Up to the proof of Theorem~\ref{thm:main} below we assume $(\Omega2\mathrm e)$ in addition; Section~\ref{sec:regularity} removes it.
Fix a comparison structure and $o\in M$; write $G=G_o$, $e=o$, $\g$ for the Lie algebra and
$\varphi(g):=\kappa(e,g)$.

\begin{lemma}\label{lem:class}
$\kappa(x,y)=\varphi(x^{-1}y)$, $\varphi(hgh^{-1})=\varphi(g)=\varphi(g^{-1})$, and $\varphi(g)=1$ iff
$g=e$.  The form $h(X,Y):=-\mathrm{Hess}_e\varphi\,(X,Y)$ on $\g$ is symmetric, positive semidefinite and
$\Ad$-invariant, and $\varphi(\exp tX)=1-\tfrac12h(X,X)t^2+o(t^2)$ as $t\to0$.
\end{lemma}
\begin{proof}
The left transport $z\mapsto\hp exz=x^{-1}z$ gives $\kappa(x,y)=\kappa(e,x^{-1}y)$; the right transport
$z\mapsto\hp zeg=zg$ gives $\varphi(g^{-1}x^{-1}yg)=\varphi(x^{-1}y)$; symmetry gives
$\varphi(g^{-1})=\varphi(g)$.  By $(\Omega2\mathrm e)$, $\varphi$ is $C^2$ near $e$, where it attains its
maximum $1$; so its differential vanishes there and its Hessian is negative semidefinite.  Conjugation
fixes $e$ and has differential $\Ad$, which gives $\Ad$-invariance of $h$; the expansion is Taylor's
formula.
\end{proof}

\begin{lemma}[duplication]\label{lem:dup}
Let $\psi\colon\R\to[-1,1]$ satisfy $\psi(2t)=2\psi(t)^2-1$ for all $t$ and $\psi(t)=1-at^2+o(t^2)$ as
$t\to0$, with $a\ge0$.  Then $\psi(t)=\cos(\sqrt{2a}\,t)$ for all $t$.
\end{lemma}
\begin{proof}
Fix $t$, put $\theta_n:=\arccos\psi(t/2^n)\in[0,\pi]$.  Since $\cos2\theta=2\cos^2\theta-1$, induction
gives $\psi(t)=\cos(2^n\theta_n)$.  Now $1-\cos\theta_n=at^24^{-n}+o(4^{-n})\to0$, so $\theta_n\to0$,
$\theta_n^2=2(1-\cos\theta_n)(1+o(1))$, and $2^n\theta_n\to\sqrt{2a}\,\lvert t\rvert$.
\end{proof}

\begin{lemma}\label{lem:closing}
For $X\in\g$ and $\lVert X\rVert^2:=h(X,X)$ one has $\varphi(\exp tX)=\cos^2\bigl(\lVert X\rVert t/\sqrt2\bigr)$
for all $t\in\R$.  Hence $h$ is positive definite, and for $X\ne0$, $\exp tX=e$ iff
$t\lVert X\rVert\in\sqrt2\,\pi\Z$.
\end{lemma}
\begin{proof}
By Lemma~\ref{lem:certaine}, $\kappa(w,\hp xwx)=\varphi(w^{-1}xw^{-1}x)=\varphi(g^2)$ with $g=w^{-1}x$, so
$(\Omega3)$ reads $\varphi(g^2)=(2\varphi(g)-1)^2$.  Then $\psi_X(t):=2\varphi(\exp tX)-1$ satisfies
$\psi_X(2t)=2\psi_X(t)^2-1$ and $\psi_X(t)=1-\lVert X\rVert^2t^2+o(t^2)$; Lemma~\ref{lem:dup} gives
$\psi_X(t)=\cos(\sqrt2\lVert X\rVert t)$.  If $\lVert X\rVert=0$ then $\varphi(\exp tX)=1$, so $\exp tX=e$
for all $t$ by faithfulness, and $X=0$.  The last claim is faithfulness.
\end{proof}

\begin{lemma}[rank one]\label{lem:rankone}
Every abelian subalgebra of $\g$ has dimension at most one.  Hence $\g\cong\R$ or $\g\cong\su(2)$.
\end{lemma}
\begin{proof}
Let $\mathfrak a\subset\g$ be abelian.  Then $\exp|_{\mathfrak a}$ is a homomorphism $(\mathfrak a,+)\to G$
whose kernel $\Gamma$ contains, by Lemma~\ref{lem:closing}, the sphere $S$ of $h$-radius $\sqrt2\pi$.  If
$\dim\mathfrak a\ge2$, $S$ is connected, so $\{\lVert v-w\rVert:v,w\in S\}=[0,2\sqrt2\pi]$ and the subgroup
$\Gamma$ contains nonzero vectors of arbitrarily small norm, contradicting injectivity of $\exp$ near $0$.

As $h$ is an $\Ad$-invariant inner product, $\g=\z\oplus\kk$ with $\z$ the centre and $\kk$ semisimple of
compact type \cite[\S7, Lemmas~7.4--7.5]{Milnor}.  By the first part $\dim\z\le1$, and $\z\ne0\ne\kk$ is impossible
because $\z\oplus\R Y$ ($0\ne Y\in\kk$) would be abelian of dimension two.  If $\g=\kk$, its maximal abelian
subalgebras are lines, so $\kk$ is compact semisimple of rank one, that is $\su(2)$ \cite[Ch.~IV]{Knapp}.
\end{proof}

\begin{proof}[Proof of Theorem~\ref{thm:main}]
We first assume $(\Omega2\mathrm e)$; by Theorem~\ref{thm:noreg} every comparison structure satisfies it, which removes the
assumption.  If $\g\cong\R$ then $G\cong\R$ or $S^1$, and $\R$ is excluded because one-parameter subgroups close
(Lemma~\ref{lem:closing}).  If $\g\cong\su(2)$ then $G\cong\SU(2)$ or $\SO(3)$.  In all cases $G$ is compact
and connected, so $\exp$ is onto and Lemma~\ref{lem:closing} determines $\varphi$.  For $G=S^1$ the
subgroup $e^{i\theta}$ first returns at $\theta=2\pi$, so $\theta=\sqrt2\lVert X\rVert t$ and
$\varphi(e^{i\theta})=\cos^2(\theta/2)=\lvert(1+e^{i\theta})/2\rvert^2$.  For $G=\SU(2)$, the unit
quaternions, $\exp tX=\cos\tau+\hat u\sin\tau$ first returns at $\tau=2\pi$, so
$\varphi=\cos^2(\tau/2)=\tfrac12(1+\re g)$.  For $G=\SO(3)$ a rotation by $\alpha$ first returns at
$\alpha=2\pi$, so $\varphi=\cos^2(\alpha/2)$ vanishes on all rotations by $\pi$, a projective plane,
contradicting $(\Omega2\mathrm d)$.  In the remaining cases $\kappa(x,y)=\varphi(x^{-1}y)=\lVert(x+y)/2\rVert^2$.
Positive semidefiniteness: $\kappa(x,y)=\ip{\Phi(x)}{\Phi(y)}$ with $\Phi(x)=(1,x)/\sqrt2\in\R\oplus A$,
and $\Phi(M)$ spans $\R\oplus A$.
\end{proof}

\begin{remark}[d'Alembert's equation]\label{rem:dalembert}
With $\psi:=2\varphi-1$ the interference law of Proposition~\ref{prop:linear}(iv) reads
$\psi(ab)+\psi(ab^{-1})=2\psi(a)\psi(b)$, d'Alembert's functional equation on $G$.  Davison
\cite{Davison08} solved it on compact groups: the solutions come from continuous homomorphisms into $\SU(2)$, and
only compact groups isomorphic to compact subgroups of $\SU(2)$ carry a \emph{basic} solution ($f(xy)=f(y)$ for all $y$
only if $x=e$; faithfulness in our sense implies this); pre-d'Alembert functions on topological monoids are of the form
$\tfrac12\operatorname{tr}\circ h$ with $h$ a continuous homomorphism into $\mathrm{Mat}_2(\C)$ \cite[Thm.~4.12]{Davison09};
see also \cite{Stetkaer}.  Corollary~\ref{cor:interference} below derives the
classification from the two-variable law without any regularity hypothesis and without these theorems.  The law $(\Omega3)$ is the diagonal case $a=b$, the duplication formula
$\psi(a^2)=2\psi(a)^2-1$.  Theorem~\ref{thm:main} together with Proposition~\ref{prop:linear}(iv) shows
that for a comparison structure, that is, under all of $(\Omega1)$--$(\Omega3)$, the duplication formula
already implies the full equation.  The binary-alternative law $(\Omega2\mathrm d)$ cannot be dropped here:
on $\SO(3)$ the function $\psi=\cos\alpha$ ($\alpha$ the rotation angle) is faithful, smooth and satisfies the
duplication formula, but for the rotations $a,b$ by $\pi/2$ about two orthogonal axes,
$\psi(ab)+\psi(ab^{-1})=-1\ne0=2\psi(a)\psi(b)$.  We did not find the implication in the literature we
searched; the converse direction (setting $b=a$) is classical.
\end{remark}

\begin{remark}
Dropping connectedness admits $S^0=\{\pm1\}\subset\R$ with the same formulas.  Law $(\Omega3)$ cannot be
dropped: on $M=S^3\times S^3$ with the componentwise heap, $\kappa(x,y)=\tfrac14\bigl(2+c(x_1,y_1)+c(x_2,y_2)\bigr)$
satisfies $(\Omega1)$ and $(\Omega2)$, its unique zero partner being $x^\perp=(-x_1,-x_2)$, but violates $(\Omega3)$
(for $w=(w_1,x_2)$ the law would require $\tfrac14(3+c')=\tfrac14(1+c(w_1,x_1))^2$ with
$c'=2c(w_1,x_1)^2-1$, i.e.\ $(1-c(w_1,x_1))^2=0$, which fails unless $w_1=x_1$).  Without $(\Omega2\mathrm d)$,
$\SO(3)$ with $\kappa=\cos^2(\alpha/2)$ is a further solution, of rank $10$.  Flat space is excluded by
$(\Omega3)$ with faithfulness: the cosine of Lemma~\ref{lem:closing} must return to $1$.
\end{remark}

\subsection{Regularity is automatic}\label{sec:regularity}
Hypothesis $(\Omega2\mathrm e)$ enters the proof above only through the quadratic form $h$.  On the line it cannot be
dropped from Lemma~\ref{lem:dup}: for $q$ continuous, positive and $1$-periodic, $f(t)=\cos\bigl(\lvert t\rvert\,q(\log_2\lvert t\rvert)\bigr)$,
$f(0)=1$, is a continuous solution of $f(2t)=2f(t)^2-1$ that is not a cosine.  The group structure nevertheless
makes $(\Omega2\mathrm e)$ redundant.

\begin{lemma}\label{lem:noline}
There is no continuous $p\colon[0,\infty)\to[-1,1]$ with $p(2t)=2p(t)^2-1$ for all $t$, $p(0)=1$, and $p(t)<1$ for $t>0$.
\end{lemma}
\begin{proof}
For $t>0$, $p(t)=-1$ would give $p(2t)=1$, and $p(t)=0$ would give $p(2t)=-1$; so $p$ has no zero, and $p>0$ by
continuity.  Write $p(t)=\cos\alpha(t)$ with $\alpha(t)\in(0,\pi/2)$ for $t>0$.  Then $p(2t)=\cos2\alpha(t)$ with
$2\alpha(t)\in(0,\pi)$, so $\alpha(2t)=2\alpha(t)$ and $\alpha(2^nt)=2^n\alpha(t)\to\infty$, a contradiction.
\end{proof}

\begin{lemma}[circle rigidity]\label{lem:circle}
Let $\psi\colon\R/2\pi\Z\to[-1,1]$ be continuous and even, with $\psi(\theta)=1$ only for $\theta=0$ and
$\psi(2\theta)=2\psi(\theta)^2-1$ for all $\theta$.  Then $\psi=\cos$.
\end{lemma}
\begin{proof}
From $\psi(0)=\psi(2\pi)=2\psi(\pi)^2-1$ and $\psi(\pi)\ne1$ we get $\psi(\pi)=-1$.  If $\psi(\theta)=-1$ then $\psi(2\theta)=1$,
so $\theta\in\{0,\pi\}$, i.e.\ $\theta=\pi$; if $\psi(\theta)=0$ then $\psi(2\theta)=-1$, so $\theta=\pm\pi/2$.  Let
$h(\theta)=\arccos\psi(\theta)$ for $\theta\in[0,\pi]$: $h$ is continuous, $h(0)=0$, $h(\pi/2)=\pi/2$, $h(\pi)=\pi$, and
$h(\theta)=\pi/2$ only at $\theta=\pi/2$; hence $h<\pi/2$ on $[0,\pi/2)$.  For $s\in[0,\pi/2]$, $\cos h(2s)=\cos2h(s)$ with both
$h(2s),2h(s)\in[0,\pi]$, so $h(2s)=2h(s)$.  Also $h>\pi/2$ on $(\pi/2,\pi]$, by the intermediate value theorem, since
$h(\pi)=\pi$ and $h=\pi/2$ only at $\pi/2$.  For $s\in[0,\pi/2]$, evenness and periodicity give
$\psi(2(\pi-s))=\psi(-2s)=\psi(2s)=\cos2h(s)$, while the duplication law gives $\psi(2(\pi-s))=\cos2h(\pi-s)=\cos(2\pi-2h(\pi-s))$
with $2\pi-2h(\pi-s)\in[0,\pi]$ because $h(\pi-s)\ge\pi/2$; injectivity of $\cos$ on $[0,\pi]$ gives $h(\pi-s)=\pi-h(s)$.  Put $g(s)=h(s)-s$; then $g(2s)=2g(s)$ on $[0,\pi/2]$ and $g(\pi-s)=-g(s)$, so
with the tent map $F(s)=2s$ on $[0,\pi/4]$, $F(s)=\pi-2s$ on $[\pi/4,\pi/2]$ one gets $\lvert g(s)\rvert=\tfrac12\lvert g(F(s))\rvert$
on $[0,\pi/2]$.  Iterating, $\lvert g\rvert\le2^{-n}\max\lvert g\rvert$ for all $n$, so $g=0$ and $\psi=\cos$ on $[0,\pi]$, hence
everywhere.
\end{proof}

\begin{lemma}\label{lem:elliptic}
Let $\g\ne0$ be a real Lie algebra in which every $\ad X$ is diagonalisable over $\C$ with purely imaginary
eigenvalues and every abelian subalgebra has dimension at most one.  Then $\g\cong\R$ or $\g\cong\su(2)$.
\end{lemma}
\begin{proof}
Let $\mathfrak r$ be the radical.  By Lie's theorem applied to the solvable algebra $\ad_\g(\mathfrak r)$ acting on
$\g\otimes\C$, every $\ad_\g X$ with $X\in[\mathfrak r,\mathfrak r]$ is nilpotent; being diagonalisable, it is $0$, so
$[\mathfrak r,\mathfrak r]$ is central.  If the centre $\z$ contains $Z\ne0$, then $\R Z+\R Y$ is abelian for every $Y$, so
$\g=\R Z$.  Otherwise $[\mathfrak r,\mathfrak r]=0$, so $\mathfrak r$ is abelian of dimension at most one; if
$\mathfrak r=\R R$, then $[X,R]=\lambda(X)R$ with $\lambda(X)$ a real eigenvalue of $\ad X$, hence $0$, and $R$ is central,
which is excluded.  So $\g$ is semisimple.  Its Killing form $B(X,X)=\operatorname{tr}(\ad X)^2$ is a sum of squares of
imaginary numbers, hence $\le0$, and nondegenerate, hence negative definite; so $\g$ is compact semisimple, its
maximal abelian subalgebras have dimension $\rk\g=1$, and $\g\cong\su(2)$ \cite[Ch.~IV]{Knapp}.
\end{proof}

\begin{theorem}\label{thm:noreg}
Let $(M,\hp\cdot\cdot\cdot,\kappa)$ satisfy $(\Omega1)$, $(\Omega2\mathrm{a,b,c})$ (in particular $\kappa$ is continuous with
values in $[0,1]$) and $(\Omega3)$.  For $o\in M$, the group $G_o$ is isomorphic to $S^1$, $\SU(2)$ or $\SO(3)$, and
$\psi=2\kappa(o,\cdot)-1$ is, respectively, $\re$, $\re$, or $\cos\alpha$ with $\alpha$ the rotation angle.  In particular
$\kappa$ is real-analytic, so $(\Omega2\mathrm e)$ holds, and under $(\Omega2\mathrm d)$ the structure is isomorphic to one of
Example~\ref{ex:models}.  In particular every comparison structure satisfies $(\Omega2\mathrm e)$, and Theorem~\ref{thm:main} holds
as stated, for Definition~\ref{def:cs} without any regularity clause.
\end{theorem}
\begin{proof}
Let $G=G_o$ (a Lie group by Lemma~\ref{lem:lie}).  The first assertions of Lemma~\ref{lem:class}, whose proof uses only
$(\Omega2\mathrm{a,b,c})$, show that $\psi$ is a continuous class function with values in $[-1,1]$, $\psi(g^{-1})=\psi(g)$,
$\psi(g)=1$ iff $g=e$; and $(\Omega3)$ reads $\psi(g^2)=2\psi(g)^2-1$ (proof of Lemma~\ref{lem:closing}).

\emph{Step 1.}  For $X\in\g\setminus0$, Lemma~\ref{lem:noline} applied to $p(t)=\psi(\exp tX)$ shows that $\exp T X=e$ for some
$T>0$.  Let $T_X$ be the least such $T$ (it exists because $\exp$ is injective near $0$).  Then $C_X=\exp(\R X)$ is a closed
subgroup and $\theta\mapsto\exp(\theta T_XX/2\pi)$ an isomorphism $\R/2\pi\Z\to C_X$.

\emph{Step 2.}  Lemma~\ref{lem:circle}, applied to $\theta\mapsto\psi(\exp(\theta T_XX/2\pi))$, gives
$\psi(\exp(\theta T_XX/2\pi))=\cos\theta$.

\emph{Step 3: abelian subalgebras are lines.}  Let $\mathfrak a$ be abelian with $\dim\mathfrak a\ge2$.  The kernel $\Gamma$ of
the homomorphism $\exp\colon\mathfrak a\to G$ is discrete and, by Step~1, meets every line through $0$ in a nonzero point.
But a discrete subgroup of $\mathfrak a$ is countable, and each of its nonzero points lies on exactly one line through $0$,
so all but countably many of the (uncountably many) lines through $0$ miss $\Gamma\setminus\{0\}$, a contradiction.

\emph{Step 4.}  $e^{T_X\ad X}=\Ad(\exp T_XX)=1$, so $\ad X$ is diagonalisable with eigenvalues in $(2\pi i/T_X)\Z$.

\emph{Step 5.}  By Lemma~\ref{lem:elliptic}, $\g\cong\R$ or $\su(2)$, so $G\cong\R$, $S^1$, $\SU(2)$ or $\SO(3)$, and $\R$ is
excluded by Step~1.  These groups are compact and connected, so every element lies on some $C_X$.  On $S^1$, Step~2 gives
$\psi(e^{i\theta})=\cos\theta$; on $\SU(2)$, $C_X=\{\cos\theta+u\sin\theta\}$ and $\psi=\cos\theta=\re$; on $\SO(3)$, $C_X$ is the
group of rotations about an axis, parametrised by the angle, and $\psi=\cos\alpha$.  In the first two cases $\kappa$ is the
kernel of Example~\ref{ex:models}; in the third, $\kappa=\cos^2(\alpha/2)$ vanishes on all rotations by $\pi$, which
$(\Omega2\mathrm d)$ excludes (Remark~\ref{rem:quotient}).
\end{proof}

\begin{corollary}[the interference law]\label{cor:interference}
Let $(M,\hp\cdot\cdot\cdot,\kappa)$ satisfy $(\Omega1)$, $(\Omega2\mathrm{a,b,c})$ and the interference law
$c(y,s_xw)+c(y,w)=2c(w,x)c(y,x)$ for all $w,x,y$.  Then it is isomorphic to one of the structures of
Example~\ref{ex:models}; neither $(\Omega2\mathrm d)$ nor $(\Omega2\mathrm e)$ is needed.
\end{corollary}
\begin{proof}
With $y=w$ the law gives $c(w,s_xw)=2c(w,x)^2-1$, which is $(\Omega3)$.  By Theorem~\ref{thm:noreg} only the case
$(\SO(3),\cos\alpha)$ remains to be excluded.  In group form (put $x=e$, $w=u$, $z=y^{-1}$) the law reads
$\psi(zu)+\psi(zu^{-1})=2\psi(z)\psi(u)$, d'Alembert's equation, and $\cos\alpha$ violates it (Remark~\ref{rem:dalembert}).
\end{proof}

Once compactness is known, Davison's theorem \cite{Davison08} gives a second route to Corollary~\ref{cor:interference}:
a faithful solution of d'Alembert's equation on a compact group forces an embedding into $\SU(2)$.

\begin{remark}[the solution excluded by $(\Omega2\mathrm d)$]\label{rem:quotient}
The third case of Theorem~\ref{thm:noreg} is the quotient of the model by its centre.  For unit quaternions $x,w$ with
$c=\ip xw$, the relative rotation $\Ad(\bar wx)$ has angle $\alpha$ with $\tfrac12(1+\cos\alpha)=c^2$, so this solution is
$\kappa'([x],[w])=c(x,w)^2$ on $S^3/\{\pm1\}\cong\SO(3)$ with the induced heap.  It satisfies $(\Omega1)$, $(\Omega2\mathrm{a,b,c})$ and
$(\Omega3)$, the last because $\ip w{x\bar wx}=2c^2-1$; its zero set from a point $[x]$ is the real projective plane of classes
of $x^\perp\cap S^3$, and the images of an orthonormal basis of $\R^4$ are four mutually orthogonal classes.  Among the solutions
of $(\Omega1)$, $(\Omega2\mathrm{a,b,c})$, $(\Omega3)$, the binary clause $(\Omega2\mathrm d)$ therefore excludes exactly this quotient, and it
is the clause that separates $\kappa=\tfrac12(1+c)$ on $S^3$ from its quotient by the centre.  In the companion paper
\cite{Companion} the same quotient appears as the real reading of $S^3$, with unobservable subgroup $\{\pm1\}$ and weight
$c^2$, and is excluded there in two ways: at the level of states by binarity, and by the input that the unobservable subgroup contains a
clock (a circle subgroup), which $\{\pm1\}$ does not.  $(\Omega2\mathrm d)$ and binarity exclude the same object at two levels.
\end{remark}

\section{Two other closing laws}\label{sec:closing}

Call a triple satisfying $(\Omega1)$, $(\Omega2\mathrm{a,b,c,e})$ with nondegenerate $h$ a
\emph{pre-comparison structure}.  Lemmas~\ref{lem:certaine}--\ref{lem:class} apply, and as in the proof of
Lemma~\ref{lem:rankone}, $\g=\z\oplus\kk$ with $\kk$ of compact type.

\begin{theorem}[parsimony]\label{thm:parsimony}
Let $\mathcal C$ be the class of noncommutative pre-comparison structures with $\kappa$ positive
semidefinite and $\inf\kappa=0$.  Then
\begin{enumerate}[label=(\roman*)]
\item $\min_{\mathcal C}\dim M=3$, attained exactly when $G\cong\SU(2)$ or $\SO(3)$;
\item among members of dimension $3$, $\min\rk\kappa=5$, where $\rk\kappa:=\dim\operatorname{span}\{\kappa(x,\cdot):x\in M\}$, attained exactly by $S^3$ with
$\kappa=\lVert(x+y)/2\rVert^2$.
\end{enumerate}
\end{theorem}
\begin{proof}
(i) Noncommutativity gives $\kk\ne0$, so $\dim M\ge\dim\kk\ge3$, with equality iff $\g=\su(2)$.
(ii) A continuous positive semidefinite class function on a compact group is a uniformly convergent sum
$\varphi=\sum_jc_j\chi_j$ of irreducible characters with $c_j\ge0$ \cite{PW,Folland}, and the
reproducing-kernel space of $\varphi(x^{-1}y)$ has dimension $\sum_{c_j>0}d_j^2$.  On $\SO(3)$ the
irreducible dimensions are $1,3,5,\dots$, and a nonconstant $\varphi$ has rank at least $9$.  On $\SU(2)$
they are $1,2,3,\dots$, and $\chi_{1/2}(g)=2\cos d(e,g)$.  Rank at most $5$ leaves the supports $\{0\}$
(constant, not faithful), $\{\tfrac12\}$ ($\varphi=\cos d$, negative for $d>\pi/2$) and $\{0,\tfrac12\}$;
for the last, $\varphi(e)=1$ and $\min\varphi=c_0-2c_{1/2}=0$ give $\varphi=\tfrac12(1+\cos d)$.
\end{proof}

Let $\Aut^{\pm}(M)$ be the group of homeomorphisms preserving $\kappa$ and either preserving the transport
or reversing it, $f\hp xyz=\hp{fz}{fy}{fx}$.  A pre-comparison structure is \emph{complete} if whenever
$\kappa(x,y)=\kappa(x',y')$ some $f\in\Aut^\pm(M)$ has $f(x)=x'$, $f(y)=y'$: the value of $\kappa$ is the only
invariant of a pair.

\begin{theorem}\label{thm:cc}
A complete noncommutative pre-comparison structure has $\g\cong\su(2)$.  Conversely $S^3$ is complete.
\end{theorem}
\begin{proof}
The stabiliser of $e$ in $\Aut^\pm(M)$ consists of continuous automorphisms $\alpha$ of $G$ with
$\varphi\circ\alpha=\varphi$ and their composites with inversion; these intertwine $\exp$ with the linear
maps $\pm d\alpha$, which preserve $h$.  Let $K:=\{\pm\beta:\beta\in\Aut(\g)\cap\OO(\g,h)\}$, a compact
group.  Choose $r_1>0$ with $\exp$ injective on the ball $B_{r_1}$ and $F:=\varphi\circ\exp\le1-\tfrac14\lVert
X\rVert^2$ there (Lemma~\ref{lem:class}), and fix $c\in(1-r_1^2/4,1)$.  Every ray from $0$ meets
$F^{-1}(c)\cap B_{r_1}$.  For two such points $X,X'$, completeness gives a stabiliser element taking $\exp X$
to $\exp X'$, so $\exp X'=\exp kX$ with $k\in K$, $\lVert kX\rVert<r_1$, and $X'=kX$.  Hence these points form
a sphere $S_\rho$ on which $K$ is transitive.  $\Aut(\g)$ preserves $\z$ and $\kk=[\g,\g]$, so
$\g=\z$ or $\g=\kk$, and noncommutativity gives $\g=\kk$.  The identity component of $\Aut(\kk)$ is
$\Ad(G)$ and $K$ has finitely many components, while $\Ad$-orbits have dimension at most
$\dim\kk-\rk\kk$.  If $\rk\kk\ge2$, the sphere $S_\rho$ of dimension $\dim\kk-1$ would be a finite union
of submanifolds of smaller dimension.  So $\rk\kk=1$ and $\kk\cong\su(2)$.  The converse is
Proposition~\ref{prop:aut}(iii).
\end{proof}

\begin{corollary}
Let $(M,\hp\cdot\cdot\cdot,\kappa)$ be a noncommutative pre-comparison structure.  Then $\g\cong\su(2)$ in each of the
following cases: (a) $(\Omega3)$ holds (Lemma~\ref{lem:rankone}, whose proof does not use $(\Omega2\mathrm d)$);
(b) the structure belongs to the class $\mathcal C$ of Theorem~\ref{thm:parsimony} and has minimal dimension in it;
(c) it is complete (Theorem~\ref{thm:cc}).
\end{corollary}

\subsection{An operational form of the round-trip law}\label{sec:operational}
The round-trip law compares a point with its own mirror image.  The following law refers instead to a test
point $y$ observed along a uniform transport $\dots,w,x,s_x(w),\dots$ (Theorem~\ref{thm:recur}): the comparison of
$y$ with the next point is determined, affinely, by its comparisons with the two previous ones.

\smallskip\noindent$(\Omega3')$ \emph{Two-readout determinacy.}  For all $w,x\in M$ there are real numbers
$A,B,D$, depending only on $(w,x)$, such that
\[
\kappa\bigl(y,s_x(w)\bigr)=A\,\kappa(y,x)+B\,\kappa(y,w)+D\qquad\text{for every }y\in M .
\]

\begin{lemma}\label{lem:opmodel}
In the structures of Example~\ref{ex:models}, $(\Omega3')$ holds with $A=2c(w,x)$, $B=-1$ and $D=v(x,w)=1-c(w,x)$;
for $w\ne\pm x$ these coefficients are unique.
\end{lemma}
\begin{proof}
With $c=2\kappa-1$ the claim reads $c(y,s_xw)=2c(w,x)c(y,x)-c(y,w)$, which is $s_x(w)=2\ip xwx-w$ paired with $y$.
For $w\ne\pm x$ the functions $y\mapsto\ip yx$, $\ip yw$ and $1$ are linearly independent on $M$.
\end{proof}

\begin{theorem}\label{thm:operational}
Let $(M,\hp\cdot\cdot\cdot,\kappa)$ satisfy $(\Omega1)$, $(\Omega2\mathrm{a,b,c})$ and $(\Omega3')$.  Fix $o\in M$, let
$G=G_o$ and $\psi:=2\kappa(o,\cdot)-1$.  Then there are a constant $m\in[0,1)$ and an isomorphism of $G$ with $S^1$ or
$\SU(2)$ under which $\psi=m+(1-m)\re$.  If in addition $\kappa$ has a zero (in particular under $(\Omega2\mathrm d)$),
then $m=0$, $(\Omega3)$ holds, and the structure is isomorphic to one of Example~\ref{ex:models}.  Of $(\Omega2)$ only
continuity, the range $[0,1]$, (a), (b), (c) and the existence of one zero of $\kappa$ are used; $(\Omega2\mathrm e)$ is not
assumed.  Thus $(\Omega1)$,
$(\Omega2)$, $(\Omega3')$ characterise the same two structures as $(\Omega1)$--$(\Omega3)$.
\end{theorem}
\begin{proof}
\emph{Step 0: group form.}  Transports preserve $c=2\kappa-1$ and commute with reflections: for
$\ell(z)=x^{-1}z$ one has $\ell(s_xw)=w^{-1}x=s_e(\ell w)$.  So $(\Omega3')$ at $(w,x)$ is equivalent to $(\Omega3')$ at
$(u,e)$ with $u=x^{-1}w$ and the same coefficients.  Writing the law with $c$ (coefficients $a=A$, $b=B$,
$d=2D+A+B-1$) and $c(y,v)=\psi(y^{-1}v)$, $z=y^{-1}$: for every $u\in G$ there are $a,b,d$ with
\begin{equation}\label{eq:op1}
\psi(zu^{-1})=a\,\psi(z)+b\,\psi(zu)+d\qquad(z\in G).
\end{equation}
By the first assertions of Lemma~\ref{lem:class} (which use only $(\Omega2\mathrm{a,c})$), $\psi$ is a class function with $\psi(g^{-1})=\psi(g)$, so $\psi(z^{-1}u^{-1})=\psi(uz)=\psi(zu)$ and
$\psi(z^{-1}u)=\psi(u^{-1}z)=\psi(zu^{-1})$; replacing $z$ by $z^{-1}$ in \eqref{eq:op1} gives, with the same coefficients,
$\psi(zu)=a\psi(z)+b\psi(zu^{-1})+d$.

\emph{Step 1: $b=-1$ off involutions.}  Subtracting, $(1+b)\bigl(\psi(zu^{-1})-\psi(zu)\bigr)=0$.  If $b\ne-1$, then
$\psi(zu^{-1})=\psi(zu)$ for all $z$, and $z=u$ gives $\psi(u^2)=\psi(e)=1$, so $u^2=e$ by faithfulness.  Let
$\Omega:=\{u\in G:u^2\ne e\}$.  For $u\in\Omega$, therefore,
\begin{equation}\label{eq:op2}
\Psi_u(z):=\psi(zu)+\psi(zu^{-1})=a(u)\,\psi(z)+d(u)\qquad(z\in G),
\end{equation}
and $(a(u),d(u))$ is unique because $\psi$ is not constant ($\psi<1$ off $e$ and $\dim G\ge1$).  The set $\Omega$ is
open and dense: its complement is closed, and if $u\mapsto u^2$ were constant on an open set it would be constant on
the connected analytic manifold $G$ (Lemma~\ref{lem:lie}), whereas $\exp(X)^2=\exp(2X)\ne e$ for small $X\ne0$.

\emph{Step 2: extension to all $u$.}  Fix $z_2\ne e$, so $\psi(z_2)<1$, and define for all $u\in G$
\[
a(u):=\frac{\Psi_u(e)-\Psi_u(z_2)}{1-\psi(z_2)},\qquad d(u):=\Psi_u(e)-a(u).
\]
These are continuous in $u$ and agree with the coefficients of \eqref{eq:op2} on $\Omega$.  For fixed $z$ both sides
of \eqref{eq:op2} are continuous in $u$, so \eqref{eq:op2} holds for all $u\in G$.  At $z=e$: $2\psi(u)=a(u)+d(u)$.

\emph{Step 3: the constant term.}  Put $Q(u,v):=\psi(uv)+\psi(uv^{-1})+\psi(u^{-1}v)+\psi(u^{-1}v^{-1})$.  Since $\psi$ is a
class function and even, $Q(u,v)=Q(v,u)$.  By \eqref{eq:op2} at $z=u$ and $z=u^{-1}$, $Q(u,v)=2a(v)\psi(u)+2d(v)$.  Hence
$a(v)\psi(u)+d(v)=a(u)\psi(v)+d(u)$, and with $\psi=\tfrac12(a+d)$,
\begin{equation}\label{eq:op3}
d(u)\bigl(a(v)-2\bigr)=d(v)\bigl(a(u)-2\bigr)\qquad(u,v\in G).
\end{equation}
If $a\equiv2$, then $d=2\psi-2$, and $\chi:=\psi-1$ satisfies $\chi(zu)+\chi(zu^{-1})=2\chi(z)+2\chi(u)$; with $z=u^n$ this
gives $\chi(u^{n+1})-2\chi(u^n)+\chi(u^{n-1})=2\chi(u)$, so $\chi(u^n)=n^2\chi(u)$.  As $\chi$ takes values in $[-2,0]$,
$\chi\equiv0$, contradicting faithfulness.  So $a(v_0)\ne2$ for some $v_0$; with $m:=d(v_0)/(2-a(v_0))$,
\eqref{eq:op3} gives $d(u)=m\,(2-a(u))$ for all $u$.

\emph{Step 4: reduction to d'Alembert's equation.}  Then $\chi:=\psi-m$ satisfies $\chi(zu)+\chi(zu^{-1})=a(u)\chi(z)$.  At
$z=e$, $2\chi(u)=a(u)(1-m)$.  If $m=1$ this forces $\chi\equiv0$, i.e.\ $\psi\equiv1$, which is impossible.  So
$\psi_0:=(\psi-m)/(1-m)$ satisfies $\psi_0(e)=1$, $a=2\psi_0$, and
\begin{equation}\label{eq:op4}
\psi_0(zu)+\psi_0(zu^{-1})=2\psi_0(z)\psi_0(u)\qquad(z,u\in G),
\end{equation}
d'Alembert's equation; with $z=u$ it gives the duplication formula $\psi_0(u^2)=2\psi_0(u)^2-1$.  If $m>1$, then
$\psi_0=(m-\psi)/(m-1)\ge1$; if $\psi_0(u)=1+\varepsilon$ with $\varepsilon>0$, the duplication formula gives
$\psi_0(u^{2^k})\ge1+4^k\varepsilon$, contradicting boundedness of $\psi$; so $\psi_0\equiv1$ and $\psi\equiv1$, impossible.  Hence
$m<1$, $\psi_0\le1$, and the duplication formula gives $\psi_0\ge-1$ (otherwise $\psi_0(u^2)>1$).

\emph{Step 5: classification.}  The function $\kappa_0(x,y):=\tfrac12\bigl(1+\psi_0(x^{-1}y)\bigr)$ is continuous with values
in $[0,1]$, symmetric, faithful and bi-invariant, and satisfies the round-trip law, which is the duplication formula.  So
Theorem~\ref{thm:noreg} applies to $\kappa_0$: $G\cong S^1$, $\SU(2)$ or $\SO(3)$, with $\psi_0=\re$ in the first two cases
and $\psi_0=\cos\alpha$ ($\alpha$ the rotation angle) on $\SO(3)$.  The
last function violates \eqref{eq:op4} (Remark~\ref{rem:dalembert}), so $G\cong S^1$ or $\SU(2)$ and
$\psi=m+(1-m)\re$, with range $[2m-1,1]$.  Since $\kappa\ge0$, $m\ge0$.  If $\kappa$ has a zero, $\psi$ attains $-1$, so
$m=0$; then $\psi=\re$, $\kappa$ is the kernel of Example~\ref{ex:models}, and Theorem~\ref{thm:main} applies.
\end{proof}

\begin{remark}\label{rem:operational}
(1) The constant $m$ is not determined by $(\Omega1)$, $(\Omega2\mathrm{a,b,c})$, $(\Omega3')$: for $m\in(0,1)$,
$\kappa=\tfrac12\bigl(1+m+(1-m)c\bigr)$ on $S^3$ satisfies these laws (with $D$ changed accordingly) and has minimum
$m>0$; for $m=\tfrac12$ this is $\kappa=\tfrac34+\tfrac14\cos d$.  Existence of a zero of $\kappa$ is the only part of
$(\Omega2\mathrm d)$ used.  (2) $\kappa=\cos^4(d/2)$ on $S^3$ satisfies $(\Omega1)$ and $(\Omega2)$ but not $(\Omega3')$: otherwise, having a zero, it
would be the kernel of Example~\ref{ex:models}.
(3) One readout does not suffice: if $c(y,s_xw)=a\,c(y,x)+d$ for all $y$ (that is, $b=0$ in \eqref{eq:op1}), Step~1 gives
$u^2=e$; so a one-readout law cannot hold for all $w,x$ when $\dim M\ge1$.  (4) For involutive steps ($u^2=e$, in the
model $w=-x$, $s_x(w)=w$) the coefficients are not unique and $b\ne-1$ is possible; Step~2 removes this case by
continuity.  These statements are checked numerically in \texttt{check\_arc\_calculus.py}.
\end{remark}

\begin{remark}[relation to Wilson's equation]
Equation \eqref{eq:op2} is Wilson's functional equation $f(zu)+f(zu^{-1})=2f(z)g(u)$ with an additional term
$d(u)$; Wilson's equation and its variants on groups are studied in \cite{Stetkaer}.  Because $\psi$ is a class
function with $\psi(e)=1$, Steps~3 and~4 reduce it to d'Alembert's equation \eqref{eq:op4} by elementary means, and
since $G$ is compact by Theorem~\ref{thm:operational}, Davison's theorem for compact groups (Remark~\ref{rem:dalembert})
identifies $\psi_0$ with $\tfrac12\operatorname{tr}\circ\rho$ for a homomorphism $\rho\colon G\to\SU(2)$; the right translates of $\psi$ then span a space of dimension at most $5$ (constants
and matrix coefficients of $\rho$), which is the content of Levi-Civita's equation for $\psi$.  Related reductions on abelian groups are in \cite{Stetkaer98} and, under the condition $f(xyz)=f(xzy)$, in
\cite[Lemma~1]{CENSZ95}.  We did not find $(\Omega3')$, or the reduction of an affine two-readout law to d'Alembert's
equation through centrality, in the searched scope.
\end{remark}

\section{Geometry of the heap \texorpdfstring{$S^3$}{S3}}\label{sec:geom}

From here on $M=S^3\subset\HH$ with $\hp xyz=x\bar yz$ and $\kappa=\cos^2(d/2)$.  Every statement is
invariant under $\Aut(M)$ and so does not depend on a base point; quaternion formulas are used for
computation only.

\begin{proposition}\label{prop:linear}
Let $V$ be the reproducing-kernel Hilbert space of $c=2\kappa-1$.
\begin{enumerate}[label=(\roman*)]
\item $\dim V=4$, $x\mapsto c(x,\cdot)$ embeds $M$ as the unit sphere of $V$, and $c(x,y)=\ip xy_V$.
\item $\kappa(x,y)=\lVert\tfrac12(x+y)\rVert^2$; the \emph{midpoint} of the minimal arc from $x$ to
$y\ne-x$ is $m(x,y)=(x+y)/(2\sqrt{\kappa(x,y)})$.
\item $s_x(w)=2\ip xw\,x-w$.
\item \emph{(interference law)} $c(y,s_xw)+c(y,w)=2\,c(w,x)\,c(y,x)$ for all $w,x,y$.
\end{enumerate}
\end{proposition}
\begin{proof}
In the model $c(x,y)$ is the Euclidean inner product of $\R^4$; (ii)--(iv) follow from
Example~\ref{ex:models}, the midpoint being the normalised sum.
\end{proof}

For $a,b\in M$ put $\ell_{a,b}(z):=\hp abz=a\bar bz$ and $r_{a,b}(z):=\hp zab=z\bar ab$; both carry $b$ to
$a$, resp.\ $a$ to $b$.  Let $\cL,\cR$ be the sets of left and right transports.

\begin{proposition}\label{prop:aut}
\begin{enumerate}[label=(\roman*)]
\item $\cL$ and $\cR$ are commuting groups, each acting simply transitively and isomorphic to $\SU(2)$;
$\cL\cap\cR=\{\mathrm{id},-\mathrm{id}\}$.
\item $\Aut(M)=\cL\cR\cong\SO(4)$ and $\Aut^\pm(M)\cong\OO(4)$.
\item $M$ is complete: $\Aut(M)$ is transitive on pairs with equal $\kappa$.
\end{enumerate}
\end{proposition}
\begin{proof}
(i) is standard for left and right multiplication on $\HH$.  (ii) If $f\in\Aut(M)$, choose $\ell\in\cL$
with $\ell f(1)=1$; then $\ell f$ is a continuous automorphism of $\SU(2)$, hence inner, so $f\in\cL\cR$;
every element of $\SO(4)$ is $z\mapsto azb$; inversion $z\mapsto\bar z$ reverses the transport and has
determinant $-1$.  (iii) $\SO(4)$ is transitive on pairs of unit vectors with given inner product.
\end{proof}

\begin{proposition}[Cartan--Schouten \cite{CS}; Loos \cite{Loos}]\label{prop:conn}
\begin{enumerate}[label=(\roman*)]
\item The differentials of $\ell_{b,a}$ and $r_{a,b}$, $T_aM\to T_bM$, are path-independent parallel
transports of two flat metric connections $\nabla^L,\nabla^R$.  On left-invariant fields their torsions
are $T^L(X,Y)=-[X,Y]$ and $T^R(X,Y)=[X,Y]$: parallel, totally skew, opposite.
\item The mean $\tfrac12(\nabla^L+\nabla^R)$ is the Levi-Civita connection of the round metric $g$, with
curvature $R(X,Y)Z=-\tfrac14[[X,Y],Z]$ and sectional curvature $1$.
\item $s_x$ is the geodesic symmetry at $x$, and $(M,s)$ is a symmetric space in the sense of Loos:
$s_xx=x$, $s_xs_x=\mathrm{id}$, $s_x(s_yz)=s_{s_xy}(s_xz)$, and $x$ is an isolated fixed point of $s_x$.
The Lie triple system of Loos's symmetric space at a point is the tangent space with $[X,Y,Z]=\tfrac14[[X,Y],Z]$, matching
$R=-\tfrac14[[X,Y],Z]$; it is isomorphic to $(T,[[X,Y],Z])$ via $X\mapsto X/2$.
\item The form $\theta=\bar x\,dx$, with values in $\im\HH$, satisfies
$d\theta+\tfrac12[\theta\wedge\theta]=0$; the torsion form $\eta(X,Y,Z)=g(X,[Y,Z])$ is bi-invariant, parallel and closed, equals twice the volume form,
and $\int_M\eta=4\pi^2$.
\end{enumerate}
\end{proposition}
\begin{proof}
(i), (ii), (iv) are the standard facts on the $(\pm)$-connections of a Lie group with bi-invariant metric
\cite{CS,Milnor}; $[X,Y]=2X\times Y$ gives sectional curvature $\tfrac14\lvert[X,Y]\rvert^2=1$ on orthonormal
pairs, and $\eta=2\,\mathrm{vol}$ with $\mathrm{vol}(S^3)=2\pi^2$.  (iii) $s_x(x\exp(t\xi))=x\exp(-t\xi)$, so
$s_x$ fixes $x$ and has differential $-1$ there; the identities are checked from $s_x(w)=x\bar wx$, and
\cite{Loos} identifies the Lie triple system.
\end{proof}

The ternary operation thus carries the same information as the pair $(\nabla^L,\nabla^R)$: transport is
commutative exactly when they agree, and with $D:=\nabla^R-\nabla^L$, $D(X,Y)=[X,Y]$, the curvature of
their mean is $R(X,Y)Z=-\tfrac14D(D(X,Y),Z)$.  In the right trivialisation the two transports and the
Levi-Civita transport along a minimal arc take a particularly simple form.

\begin{proposition}\label{prop:LCarc}
Let $u,w\in M$, $w\ne-u$, $g:=w\bar u$ (so $\ell_{w,u}=g\,\cdot$), and $h:=g^{1/2}=(u+w)\bar u/\lVert u+w\rVert$,
the square root with $\re h>0$.  For $\xi\in\im\HH$:
\begin{enumerate}[label=(\roman*)]
\item $\nabla^R$-transport from $u$ to $w$ maps $\xi u\mapsto\xi w$, and $\nabla^L$-transport maps
$\xi u\mapsto(g\xi\bar g)\,w$;
\item Levi-Civita transport along the minimal arc from $u$ to $w$ maps $\xi u\mapsto(h\xi\bar h)\,w$;
\item this transport equals $-ds_m$ restricted to $T_uM$, where $m=hu=m(u,w)$ is the midpoint of the arc.
\end{enumerate}
\end{proposition}
\begin{proof}
(i) $r_{u,w}(z)=z\bar uw$ maps $\xi u$ to $\xi w$, and $\ell_{w,u}(z)=gz$ maps $\xi u$ to
$g\xi u=(g\xi\bar g)gu=(g\xi\bar g)w$.  (iii) $s_m(z)=m\bar zm$ has differential $v\mapsto m\bar vm$; for
$v=\xi u$, $\bar v=-\bar u\xi$, so $-ds_m(v)=m\bar u\xi m=h\xi hu=(h\xi\bar h)h^2u=(h\xi\bar h)w$, using $m=hu$
and $h^2u=w$.  (ii) By Proposition~\ref{prop:linear}(iii), $s_m$ is the restriction of the linear reflection
$v\mapsto2\ip mvm-v$, so $-ds_m(v)=v-2\ip mvm$ is the reflection of $\R^4$ in the hyperplane $m^\perp$.  Levi-Civita
transport on the round sphere along the great circle through $u,w$ fixes vectors orthogonal to the plane
$\operatorname{span}(u,w)$ and rotates that plane with the arc, carrying the unit tangent at $u$ to the unit tangent at
$w$.  The reflection in $m^\perp$ fixes $\operatorname{span}(u,w)^\perp$ (which lies in $m^\perp$) and acts on
$\operatorname{span}(u,w)$ as the reflection across the line $m^\perp\cap\operatorname{span}(u,w)$.  Since $2\ip mum=u+w$, it maps
$u$ to $-w$, hence $T_uM=u^\perp$ onto $w^\perp=T_wM$, and it carries the unit tangent $t_u$ at $u$ pointing towards $w$ to the
unit tangent $t_w$ at $w$ pointing away from $u$ (in the plane, with $m$ at angle $0$ and $u,w$ at $\mp\varepsilon/2$, it is
$\theta\mapsto\pi-\theta$).  So both maps agree on $T_uM$, and (ii) follows from (iii).  Equivalently, transport is the
differential of the transvection $s_ms_u$ \cite[Ch.~IV, \S3]{Helgason}.
\end{proof}

In the right trivialisation, then, $\nabla^R$ transports by $\Ad(1)$, $\nabla^L$ by $\Ad(g)$, and the
Levi-Civita connection by $\Ad(g^{1/2})$.  The element $h=(1+w\bar u)/\lvert1+w\bar u\rvert$ is the familiar half-angle rotor of
geometric algebra and of rotation computations; what the proposition adds is its identification with Levi-Civita transport
on the heap and with $-ds_m$.  Around a closed polygon of minimal arcs the displacements $g_k$
multiply to $1$ ($\nabla^L$ is flat), and their square roots in general do not; this defect is computed in
Section~\ref{sec:stokes}.

\begin{theorem}\label{thm:obstruction}
If a sphere $S^n$, $n\ge1$, carries a continuous ternary operation satisfying $(\Omega1)$, then
$n\in\{1,3\}$.
\end{theorem}
\begin{proof}
By Lemmas~\ref{lem:certaine} and \ref{lem:lie}, $S^n$ would be a compact connected Lie group.  If abelian it
is a torus, so $n=1$.  Otherwise the bi-invariant $3$-form $g(X,[Y,Z])$ of a bi-invariant metric is nonzero,
and bi-invariant forms inject into de~Rham cohomology \cite{CE}, so $H^3(S^n;\R)\ne0$ and $n=3$.
\end{proof}

\begin{proposition}[weakened associativity]\label{prop:moufang}
Replace the third identity of $(\Omega1)$ by the requirement that it hold in every subsystem generated by
two elements.
\begin{enumerate}[label=(\roman*)]
\item The unit octonions $S^7$ with $\hp xyz=(x\bar y)z$ and $\kappa=\lVert\tfrac12(x+y)\rVert^2$ satisfy the
weakened $(\Omega1)$, $(\Omega2)$, $(\Omega2\mathrm e)$ (for the standard smooth structure) and $(\Omega3)$, and violate
the third identity.
\item If $S^n$, $n\ge1$, carries a continuous ternary operation with the first two identities, then
$n\in\{1,3,7\}$ \cite{Adams}.
\end{enumerate}
\end{proposition}
\begin{proof}
(i) The subsystem generated by two unit octonions lies in the subalgebra they generate (note $\bar x=2\re x-x$), which is
associative by Artin's theorem, so the third identity holds there.  Left and right multiplications by unit octonions are
isometries, and alternativity with $x\bar w+w\bar x=2\ip xw$ gives $(x\bar w)x=2\ip xwx-w$, which yields $(\Omega3)$ as in
Example~\ref{ex:models}.  (ii) $x\cdot y:=\hp xoy$ has the two-sided unit $o$, so $S^n$ is an $H$-space.
\end{proof}

The weakened axiom of Proposition~\ref{prop:moufang} does not by itself make $x\cdot y=\hp xoy$ di-associative, nor give
$\hp xyz=(x\cdot y^{-1})\cdot z$ (Proposition~\ref{prop:weaknot} below): a subsystem generated by two elements is a line, while
di-associativity concerns subsystems generated by three elements $o,x,y$.  For the unit octonions both properties hold by
Artin's theorem, since $x,y$ generate an associative subalgebra.  In the following they are therefore hypotheses.
Under these hypotheses the classification can be pushed through the tangent algebra.  A loop on a real-analytic manifold is
\emph{real-analytic} if its multiplication and both divisions are analytic.  We use the following facts, collected as (L).  The
local statements are Mal'cev's \cite[Thms.~1,~2]{Malcev55}: a local analytic alternative loop (products of powers of two elements
associate near $e$, which di-associativity gives) has canonical coordinates of the first kind, its tangent algebra is binary-Lie, and in
these coordinates the product is the Campbell--Hausdorff series; the tangent algebra determines the local loop up to local
isomorphism, and that of a local analytic Moufang loop is a Malcev algebra \cite[Thm.~4]{Malcev55}.  The local correspondence between
Malcev algebras and local analytic Moufang loops is Kuz'min's \cite{Kuzmin71}; the global theory of analytic Moufang loops is Kerdman's
\cite{Kerdman} (see \cite{NS02}).

\smallskip\noindent(L) \emph{Let $L$ be a connected real-analytic loop with identity $e$ that is di-associative (every
subloop generated by two elements is a group).  Its tangent space $\mathfrak m$ at $e$, with the bracket $[X,Y]$ defined by
the second-order term of the commutator, is a binary-Lie algebra (every two-generated subalgebra is a Lie algebra).
There is an analytic map $\exp\colon\mathfrak m\to L$, a diffeomorphism from a ball $B_{2r}$ onto a neighbourhood of $e$, such
that $t\mapsto\exp(tX)$ is a homomorphism $\R\to L$ for every $X$; in these coordinates the product is given by the
Campbell--Hausdorff series of $\mathfrak m$, so $\exp X\cdot\exp Y=\exp\bigl(X+Y+\tfrac12[X,Y]+\dots\bigr)$, and the commutator
$(x\cdot y)/(y\cdot x)$ of $x=\exp sX$, $y=\exp tY$ is $\exp\bigl(st[X,Y]+O((\lvert s\rvert+\lvert t\rvert)^3)\bigr)$.}

\smallskip
Statement (L) holds for every connected real-analytic di-associative loop.  Everything near $e$ is Mal'cev's theorem, with a local
exponential $\exp_0$ on a ball $B_{2r}$.  The global part is elementary: for $X\in\mathfrak m$ choose $n$ with $X/n\in B_r$ and put
$\exp X:=\exp_0(X/n)^n$; powers are unambiguous because an element generates a group, and the value does not depend on $n$ because
$\exp_0(X/(nm))^m=\exp_0(X/n)$.  For real $s,t$ and large $N$ the elements $\exp_0(sX/N)$ and $\exp_0(tX/N)$ commute and multiply to
$\exp_0((s+t)X/N)$, so in the group they generate $\exp((s+t)X)=\exp(sX)\exp(tX)$; and on $B_{nr}$, $\exp$ is a composite of $\exp_0$ and
loop multiplications, hence analytic.

\begin{proposition}\label{prop:moufangclass}
Let $M$ be a connected real-analytic manifold with an analytic ternary operation and $o\in M$ such that $x\cdot y:=\hp xoy$ is
a di-associative loop with $\hp xyz=(x\cdot y^{-1})\cdot z$ (it is then a real-analytic loop, since $x\backslash z=\hp oxz$ and $z/y=\hp zyo$;
these hypotheses cannot be replaced by the weakened
axiom, Proposition~\ref{prop:weaknot}), and let $\kappa$ satisfy $(\Omega2)$ (with loop
translations in (c)) and $(\Omega3)$.  We use the following two facts; (K1) is a theorem of Grishkov, and (K2) is
proved after the proof of the proposition from published classifications:
\begin{enumerate}[label=(K\arabic*),leftmargin=2.6em]
\item a finite-dimensional semisimple binary-Lie algebra over a field of characteristic $0$ is a Malcev algebra
(Grishkov's theorem \cite{Grishkov80}; a simple algebra with nonzero product is semisimple);
\item a simple real Malcev algebra that is not a Lie algebra and carries a definite invariant form is isomorphic to
$\im\mathbb O$ with $[x,y]=xy-yx$ (stated in \cite[\S2]{Loginov01}).
\end{enumerate}
Then:
\begin{enumerate}[label=(\roman*)]
\item the tangent algebra $\mathfrak m$ is $\R$, $\su(2)$ or $\im\mathbb O$, and a neighbourhood of $o$ is isomorphic, as a local loop, to a
neighbourhood of $1$ in $S^1$, $S^3$ or the unit octonions $S^7$;
\item $x\cdot y$ is a Moufang loop (Proposition~\ref{prop:mfauto} below; this was a hypothesis in an earlier version), and, by the
global correspondence for analytic Moufang loops \cite{Kerdman}, $M$ is isomorphic to $S^1$, $S^3$ or $S^7$ with
$\kappa=\lVert(x+y)/2\rVert^2$.
\end{enumerate}
\end{proposition}

The proof uses the following lemma; it shows that commuting directions span abelian subgroups, which was an additional hypothesis in an earlier version of this proposition.

\begin{lemma}\label{lem:K3}
Let $L$ be a connected real-analytic di-associative loop, with (L), and let $X,Y\in\mathfrak m$ with $[X,Y]=0$.  Then $\exp(sX)$ and $\exp(tY)$ commute for all $s,t$, the
closure $H$ of the subloop they generate is an abelian Lie group, and $\exp(sX+tY)=\exp(sX)\cdot\exp(tY)$.
\end{lemma}
\begin{proof}
(1) Put $a_k=\exp(2^{-k}X)$, $b_k=\exp(2^{-k}Y)$ and let $\Gamma_k$ be the subloop generated by $a_k,b_k$, a group by
di-associativity.  Since $a_k=a_{k+1}^2$ and $b_k=b_{k+1}^2$, $\Gamma_k\subset\Gamma_{k+1}$, so $\Gamma=\bigcup_k\Gamma_k$ is a subloop in
which any three elements lie in a common $\Gamma_k$; it is a group containing $\exp(qX)$, $\exp(qY)$ for all dyadic $q$.

(2) Its closure $H$ is a subloop (multiplication and the two divisions are continuous) and associative (a closed
condition), hence a closed subgroup of $L$ containing $\exp(\R X)$ and $\exp(\R Y)$.

(3) $L$ has no small subgroups: if $e\ne g=\exp Z$ with $Z\in B_r$, then $g^n=\exp(nZ)$, and for the $n$ with
$n\lvert Z\rvert\in[r,2r)$ injectivity on $B_{2r}$ gives $g^n\notin\exp(B_r)$.  So $H$ is a locally compact group without small
subgroups, hence a Lie group \cite{Gleason52,Yamabe53,MZ}.

(4) Every continuous homomorphism $\phi\colon\R\to L$ is $t\mapsto\exp(tW)$: for small $t$, $\phi(t)=\exp W(t)$ with $W(t)\in B_r$;
$\exp(2W(t/2))=\phi(t/2)^2=\phi(t)$ and injectivity on $B_{2r}$ give $W(t/2)=W(t)/2$.  Hence, for fixed small $t_0$ and all integers $k$ and $j\ge0$,
$\phi(kt_02^{-j})=\exp(kW(t_0)2^{-j})$, and continuity gives $\phi(t)=\exp(tW(t_0)/t_0)$.  Applying this to $t\mapsto\exp_H(tZ)$ defines $\delta\colon\mathfrak h\to\mathfrak m$ with
$\exp_H Z=\exp(\delta Z)$.  It is injective.  It is additive: by the Lie product formula in $H$,
\[
\exp_H(Z_1+Z_2)=\lim_{n\to\infty}\bigl(\exp_H(Z_1/n)\cdot\exp_H(Z_2/n)\bigr)^n,
\]
and by (L) the element $g_n=\exp(\delta Z_1/n)\cdot\exp(\delta Z_2/n)$ is
$\exp C_n$ with $nC_n\to\delta Z_1+\delta Z_2$, so $g_n^n=\exp(nC_n)\to\exp(\delta Z_1+\delta Z_2)$; applied to $tZ_1,tZ_2$ for small $t$,
this gives $\delta(Z_1+Z_2)=\delta Z_1+\delta Z_2$.  It preserves brackets: the commutator of $\exp_H(sZ_1)$ and $\exp_H(tZ_2)$ is
the same element of $H\subset L$ whether computed in $H$ or in $L$, and comparing the coefficients of $st$ in (L) and in the
group $H$ gives $\delta[Z_1,Z_2]_{\mathfrak h}=[\delta Z_1,\delta Z_2]$.

(5) Every continuous one-parameter subgroup of the Lie group $H$ is $t\mapsto\exp_H(tZ)$ for some $Z\in\mathfrak h$; applied to
$t\mapsto\exp(tX)$ and $t\mapsto\exp(tY)$ this gives $X_H,Y_H\in\mathfrak h$ with $X=\delta X_H$, $Y=\delta Y_H$ (by (4)).  Then $\delta[X_H,Y_H]=[X,Y]=0$ gives $[X_H,Y_H]=0$, so the one-parameter subgroups commute, $\Gamma$ and $H$
are abelian, and $u\mapsto\exp(usX)\cdot\exp(utY)$ is a continuous homomorphism $\R\to L$ with velocity $sX+tY$ at $u=0$ (by (L));
by (4) it equals $u\mapsto\exp(u(sX+tY))$, and $u=1$ gives the last claim.
\end{proof}

\begin{proof}[Proof of Proposition~\ref{prop:moufangclass}]
In a di-associative loop every element $x$ has a two-sided inverse $x^{-1}$, and the inverse properties
$x^{-1}\cdot(x\cdot y)=y=(y\cdot x)\cdot x^{-1}$ hold, because $x,y$ generate a group.  Loop translations preserve $\kappa$, so the left
translation by $x^{-1}$ gives $\kappa(x,y)=\varphi(x^{-1}\cdot y)$ with $\varphi=\kappa(o,\cdot)$, and $\varphi$ is invariant under
$c_a\colon g\mapsto(a\cdot g)\cdot a^{-1}$, a composite of two translations with $c_a(o)=o$.  For $a=\exp tX$ and $g=\exp sY$,
by (L) the element $(a\cdot g)\cdot a^{-1}$ is given by the Campbell--Hausdorff series in the Lie subalgebra generated by $X$ and
$Y$, so $c_a(g)=\exp\bigl(sY+st[X,Y]+O((\lvert s\rvert+\lvert t\rvert)^3)\bigr)$; so $\frac{d}{dt}\big|_{t=0}dc_{\exp tX}=\ad X$ at $o$.  Hence
$h=-\mathrm{Hess}_o\varphi$ satisfies $h([X,Y],Y)=0$ for all $X,Y$, and by polarisation $h([X,Y],Z)=-h(Y,[X,Z])=h(X,[Y,Z])$,
using anticommutativity.  One-parameter subloops are groups and $\kappa(o,\hp gog)=\varphi(g\cdot g)$, so Lemmas~\ref{lem:dup} and
\ref{lem:closing} hold verbatim: $h$ is positive definite and one-parameter subloops close at $h$-radius $\sqrt2\pi$.  For an
anticommutative algebra with a definite invariant form, orthogonal complements of ideals are ideals, and an ideal $I$
with $[I,I]=0$ satisfies $h([\mathfrak m,I],I)=h(\mathfrak m,[I,I])=0$, hence is central; so $\mathfrak m$ is the orthogonal sum of its
centre and of simple ideals.  By Lemma~\ref{lem:K3} the argument of Lemma~\ref{lem:rankone} applies to abelian
subalgebras: they are lines.  This excludes a centre together with a simple ideal and two simple ideals, so $\mathfrak m$ is
$\R$ or simple.  If $\mathfrak m=\R X$, the one-parameter subloop $S=\exp(\R X)$ contains a neighbourhood $V$ of $o$; it is open,
since $y\cdot V\subset S$ for $y\in S$, and closed, since $z\cdot v\in S$ with $v\in V$ forces $z=(z\cdot v)/v\in S$; so $M=S$ is a
one-dimensional connected Lie group, $\R$ or $S^1$, and $\R$ is excluded because one-parameter subloops close.  A simple $\mathfrak m$ is a Malcev algebra by (K1), hence either a compact simple Lie algebra of rank one,
$\su(2)$, or, by (K2), $\im\mathbb O$, whose commuting pairs are dependent.  In canonical coordinates the product is the
Campbell--Hausdorff series of $\mathfrak m$ (L), as it is for $S^1$, $S^3$, $S^7$, which gives (i).  For (ii), $\mathfrak m$ is a Malcev algebra
(Lie algebras are Malcev algebras, and so is $\im\mathbb O$), so $x\cdot y$ is Moufang by Proposition~\ref{prop:mfauto}; more directly, (i) gives a
local isomorphism with the Moufang loop $S^1$, $S^3$ or $S^7$, so the Moufang identity holds near $(o,o,o)$, and the identity-theorem argument
in the proof of Proposition~\ref{prop:mfauto} extends it to $M^3$, without Kuz'min's theorem.  A connected
analytic Moufang loop is a quotient of the simply connected one with the same tangent algebra by a discrete central
normal subloop \cite{Kerdman}; this leaves $\R$, $S^1$, $S^3$, $\SO(3)$, $S^7$ and $S^7/\{\pm1\}$; $\R$ is excluded by the closing law,
and on $\SO(3)$ and $S^7/\{\pm1\}$ the closing law makes $\varphi$ vanish on all elements of order two, which $(\Omega2\mathrm d)$
excludes.
\end{proof}

\begin{proof}[Proof of (K2)]
Let $\mathfrak m$ be a simple real non-Lie Malcev algebra with a positive definite invariant form $h$.  (a) The centroid $\Gamma$
(endomorphisms commuting with all multiplications) of a simple algebra with nonzero product is a field \cite[Ch.~II]{Schafer66},
here $\R$ or $\C$.  Every invariant bilinear form is $h(T\cdot,\cdot)$ with $T\in\Gamma$: for an invariant $B$ define $T$ by
$B(X,Y)=h(TX,Y)$; then $h(T[a,b],c)=B([a,b],c)=B(a,[b,c])=h(Ta,[b,c])=h([Ta,b],c)$, so $T[a,b]=[Ta,b]$.  If $\Gamma=\C$, then $\mathfrak m$
is a complex simple non-Lie Malcev algebra, hence the seven-dimensional one \cite{Sagle62,Kuzmin68}, which carries a
nondegenerate $\C$-bilinear invariant form $K$ (from the bilinear form of the complex octonion norm); by the above,
$h=\re(\lambda K)$ for some $\lambda\in\C$, and $h(iX,iX)=-h(X,X)$ contradicts definiteness.  So $\Gamma=\R$, $\mathfrak m$ is central simple,
and invariant forms on $\mathfrak m$ are unique up to a scalar.  (b) By Filippov's theorem \cite{Filippov76} (see
\cite[\S2]{Loginov01}), $\mathfrak m$ is the traceless part $C_0$ of an octonion algebra $C$ over $\R$ with $[x,y]=xy-yx$, and $C$ is
the division algebra $\mathbb O$ or the split octonions.  (c) On $C_0$ the bilinear form $t(xy)$, $t$ the trace, is invariant,
because $t((xy)z)=t(x(yz))$ and $t(xy)=t(yx)$ in alternative algebras \cite{Schafer66}; it is a multiple of the norm form.
(d) For the split octonions the norm has signature $(4,4)$ and $1$ is anisotropic, so on $C_0=1^\perp$ it has signature $(3,4)$,
which no multiple makes definite.  Hence $C=\mathbb O$ and $\mathfrak m\cong\im\mathbb O$.
\end{proof}

The Moufang identity in (ii) is a consequence of results of Mal'cev and Kuz'min together with the identity theorem for analytic
maps; we record it separately, since it answers a question posed in an earlier version of this paper.

\begin{proposition}[Mal'cev, Kuz'min]\label{prop:mfauto}
Let $L$ be a connected real-analytic di-associative loop whose tangent algebra $\mathfrak m$ is a Malcev algebra.  Then $L$
is a Moufang loop: $(x\cdot y)\cdot(z\cdot x)=(x\cdot(y\cdot z))\cdot x$ for all $x,y,z\in L$.
\end{proposition}
\begin{proof}
Let $H(X,Y)$ be the product in the canonical coordinates of (L), an analytic map near $(0,0)$ whose Taylor series is the
Campbell--Hausdorff series of $\mathfrak m$ (by (L), that is, \cite[Thm.~2]{Malcev55}, also stated in \cite[\S2]{Kuzmin71}).  By Kuz'min's Theorem~1 \cite{Kuzmin71}, for a finite-dimensional Malcev algebra over a field
of characteristic $0$ the formal loop defined by the Campbell--Hausdorff series is a formal Moufang loop, that is,
$H(H(X,Y),H(Z,X))=H(H(X,H(Y,Z)),X)$ as formal power series.  Both sides are analytic near $0$ (compositions of analytic maps
vanishing at $0$), and the Taylor series of a composite is the composite of the Taylor series; so the two sides have the same Taylor
series at $0$ and agree on a connected neighbourhood of $0$.  Hence $F(x,y,z)=(x\cdot y)\cdot(z\cdot x)$ and $G(x,y,z)=(x\cdot(y\cdot z))\cdot x$,
analytic maps $L^3\to L$, agree on a neighbourhood of $(e,e,e)$.  Let $W$ be the interior of $\{F=G\}$; it is open and nonempty.  It is
closed: if $p\in\overline W$, then $F(p)=G(p)$ by continuity; take a chart around $F(p)$ and a connected chart neighbourhood $U$ of $p$
mapped into it by $F$ and $G$; the coordinate functions of $F$ and $G$ are analytic on $U$ and agree on the nonempty open set $U\cap W$,
so they agree on $U$, and $p\in W$.  Since $L^3$ is connected, $W=L^3$.
\end{proof}

So the Moufang identity is automatic for analytic di-associative loops with Malcev tangent algebra.  Proposition~\ref{prop:moufangclass}(ii)
is thus a special case of the description of connected analytic Moufang loops through their Malcev algebras \cite{Kerdman} (see
\cite{NS02}); what the comparison $\kappa$ adds is the reduction to the rank-one algebras $\R$, $\su(2)$, $\im\mathbb O$ and the exclusion of
the quotients by $\{\pm1\}$.  On this Lie-theoretic route analyticity and di-associativity are needed: di-associativity cannot be
replaced by the two-generator axiom (Proposition~\ref{prop:weaknot}), and we do not know a reference showing that a di-associative loop on
a topological manifold with continuous operations is analytic (Section~\ref{sec:open}, Q1).  Under invariance of $\kappa$ under the loop
translations, Theorem~\ref{thm:round} below gives the conclusion without them (Remark~\ref{rem:roundrel}(ii)).

\begin{proposition}[the weakened axiom does not give di-associativity]\label{prop:weaknot}
The Mal'cev identities together with the third identity in every subsystem generated by two elements do not imply that
$x\cdot y=\hp xoy$ is di-associative, nor that $\hp xyz=(x\cdot y^{-1})\cdot z$:
\begin{enumerate}[label=(\Alph*)]
\item On the seven points of the Fano plane, the nonzero vectors of $\mathbb F_2^3$, put $\hp xxz=z$, $\hp xzz=x$, $\hp xyx=x+y$, and for
pairwise distinct $x,y,z$: $\hp xyz=y$ if $x+y+z=0$ (a line) and $\hp xyz=x+y+z$ otherwise.
\item On $S^3$ put $\hp abx=a\bar ba\bar bx\bar ab$, i.e.\ $b\,u^2w\bar u$ with $u=\bar ba$, $w=\bar bx$.  This operation is real-analytic,
the standard comparison $\kappa=\tfrac12(1+\ip xy)$ satisfies $(\Omega2\mathrm{a,b,d,e})$, $(\Omega3)$ and is invariant under all left
transports $x\mapsto\hp abx$, but not under the right transports.
\end{enumerate}
\end{proposition}
\begin{proof}
(A) The Mal'cev identities hold by definition.  For distinct $p,q$ the line $L=\{p,q,p+q\}$ is closed under the operation, and the
bijection $p\mapsto0$, $q\mapsto1$, $p+q\mapsto2$ carries it to the heap of $\Z/3$: for three distinct indices $i-j+k\equiv j$, and
$2i-j$ is the third index; so the subsystem generated by $p,q$ is $L$, on which the third identity holds.  For a base $o$, the
lines through $o$ are subloops of order $3$.  If the loop were di-associative, two points $x,y$ not on a common line through $o$
would generate a group containing two distinct subgroups of order $3$, hence of order at least $9>7$.  An exhaustive
computation also shows that the third identity fails for $7560$ of the $7^5$ quintuples and that $\hp xyz\ne(x\cdot y^{-1})\cdot z$
for $192$ of the $343$ triples ($o=1$).

(B) If $a,b,x$ lie on a great circle, then $u,w$ commute and $\hp abx=bu\,w=a\bar bx$; so every subsystem generated by two
elements lies on the great circle through them, where the operation is the standard heap, and the Mal'cev identities hold
($\hp aax=x$, $\hp abb=a\bar ba\bar a b=a$).  The round trip $\hp xwx=x\bar wx\bar w x\bar xw=x\bar wx$ is the standard one, so $(\Omega3)$
holds, and $x\mapsto(a\bar ba\bar b)\,x\,(\bar ab)$ is an isometry.  At $o=1$ the product is $u*w=u^2w\bar u$, and for
$u=\cos t+j\sin t$ one computes $u*i=i\cos3t-k\sin3t$; so $1*i=i=u'*i$ for $u'=-\tfrac12+\tfrac{\sqrt3}2j$.  Hence $(S^3,*)$ is not even a
loop, so it is not di-associative, and by Lemma~\ref{lem:certaine} the third identity cannot hold globally.  Explicitly, for
$b=1$, $a=x=\tfrac12(1+i+j+k)$, $y=\tfrac12(1+i+j-k)$, $w=\tfrac12(1+i-j+k)$ the two sides $\hp{\hp abx}yw$ and $\hp ab{\hp xyw}$ are
antipodal, so the third identity fails by $2$, the largest possible distance in $\R^4$; and the right transport $a\mapsto\hp a1w$ with
$w=(i-j)/\sqrt2$ takes the same value at $u=\tfrac12(1+i+j+k)$ and at $\bar u$, so it changes $\kappa(u,\bar u)=\tfrac14$ to $1$ (numerically the largest
change of $\kappa$ under right transports is $4/(3\sqrt3)\approx0.770$; we do not prove this).
\end{proof}

Example~(A) is the Mal'cev term $p(x,y,z)=y\circ(x\circ z)$ of the Steiner quasigroup of the Fano plane, in which $x\circ x=x$
and $x\circ y$ is the third point of the line through $x\ne y$ (in the encoding above, $x\circ y=x+y$): since $x\circ(x\circ y)=y$ and
$x\circ y=y\circ x$, one has $p(x,x,z)=z$ and $p(x,z,z)=x$.  Steiner quasigroups form a congruence-permutable variety
\cite{Quackenbush76}.  Example~(A) is finite, so it bears only on the algebraic implication.  Example~(B) satisfies everything
except invariance under right transports.  On $S^3$, two-sided invariance and continuity force $a\bar bx$ or its opposite
(Theorem~\ref{thm:rigid}); on $S^7$ they do not (Proposition~\ref{prop:S7}).

Let $A$ be $\HH$ or the octonions $\mathbb O$, with the Euclidean norm, unit $1$ and conjugation $u\mapsto\bar u$, let $S(A)$ be its
unit sphere and $\kappa$ the standard comparison on it.  A \emph{unital composition product} on $A$ is a bilinear map
$\mu\colon A\times A\to A$ with $\lvert\mu(u,v)\rvert=\lvert u\rvert\,\lvert v\rvert$ and $\mu(1,v)=v=\mu(v,1)$.  We use two classical facts.
(H1) In a unital composition algebra every element satisfies $\mu(v,v)=2\ip v1\,v-\lvert v\rvert^21$ \cite[Ch.~1]{SV00}.
(H2) \emph{Hurwitz's theorem} \cite{Hurwitz1898}, \cite[Ch.~1]{SV00}, \cite{CS03}: a unital composition algebra over $\R$ with positive
definite norm is isomorphic to $\R$, $\C$, $\HH$ or $\mathbb O$.

\begin{lemma}[classical, {\cite[Ch.~1]{SV00}}]\label{lem:hurwitz}
Let $\mu$ be a unital composition product on $A$.  Then $\mu(u,v)=\phi(\phi^{-1}u\cdot\phi^{-1}v)$ for an orthogonal map $\phi$ of $A$
fixing $1$, that is, $\phi\in\OO(\im A)$, unique up to right composition with $\Aut(A)$.  For $A=\HH$: $\mu(u,v)=uv$ if
$\det\phi=1$, and $\mu(u,v)=vu$ if $\det\phi=-1$.
\end{lemma}
\begin{proof}
By (H2) and $\dim A$ there is an algebra isomorphism $\psi\colon(A,\mu)\to A$.  It maps $1$ to $1$ and is isometric: by (H1) in both
algebras, $\psi(v)^2=2\ip v1\,\psi(v)-\lvert v\rvert^2$ and $\psi(v)^2=2\ip{\psi(v)}1\,\psi(v)-\lvert\psi(v)\rvert^2$, and for $v\notin\R1$ the
elements $1,\psi(v)$ are independent, so $\lvert\psi(v)\rvert=\lvert v\rvert$.  Put $\phi=\psi^{-1}$; two such maps differ by an
automorphism.  For $A=\HH$ every rotation of $\im\HH$ is $\Ad(q)$ for a unit $q$, an automorphism, which gives $uv$.  If
$\det\phi=-1$, then $\phi=\Ad(q)\circ c$ with $c(u)=\bar u$ (on $\im\HH$, $c=-\mathrm{id}$ has determinant $-1$), and since $c$ reverses
products, $\mu(u,v)=vu$.
\end{proof}

\begin{lemma}[normal form]\label{lem:normal}
Let $T\colon S(A)^3\to S(A)$ be any map.  Then $T$ satisfies the Mal'cev identities $T(a,a,x)=x$, $T(a,x,x)=a$, and $\kappa$ is
invariant under all left transports $x\mapsto T(a,b,x)$ and all right transports $x\mapsto T(x,a,b)$, if and only if
\[
T(a,b,x)=b\,\mu_b(\bar ba,\bar bx)\qquad(a,b,x\in S(A))
\]
for a family $(\mu_b)_{b\in S(A)}$ of unital composition products on $A$.  The family is determined by $T$, through
$\mu_b(u,v)=\bar b\,T(bu,b,bv)$ for unit $u,v$, and $T$ is continuous if and only if $b\mapsto\mu_b$ is continuous.  Every such $T$
satisfies $T(x,w,x)=2\ip xw\,x-w$, hence $(\Omega3)$, and on every great circle $C$ it is the heap of the circle group:
$T(ce^{rX},ce^{sX},ce^{tX})=ce^{(r-s+t)X}$ for $C=\{ce^{tX}\}$; so the third identity holds on every subsystem generated by two
elements.
\end{lemma}
\begin{proof}
A map $f$ of $S(A)$ that preserves $\kappa$ preserves inner products, so it is the restriction of an orthogonal map: for an orthonormal
basis $(e_i)$ the $f(e_i)$ are orthonormal and $\ip{f(x)}{f(e_i)}=\ip x{e_i}$, so $f(x)=\sum_i\ip x{e_i}f(e_i)$.  Left multiplication
$\ell_b(y)=by$ by a unit $b$ is orthogonal with inverse $\ell_{\bar b}$, since $\bar b(by)=(\bar bb)y$ by alternativity.  Fix $b$.  For a unit $u$
the map $v\mapsto\bar b\,T(bu,b,bv)$ of $S(A)$ is $\ell_{\bar b}\circ T(bu,b,\cdot)\circ\ell_b$ and preserves $\kappa$, so it is the restriction of
an orthogonal map $L_u$; likewise $u\mapsto\bar b\,T(bu,b,bv)$ is the restriction of an orthogonal map $R_v$, by invariance under the
right transport $T(\cdot,b,bv)$.  Put $\mu_b(u,v)=\lvert u\rvert\lvert v\rvert\,\bar b\,T(b\hat u,b,b\hat v)$ with $\hat u=u/\lvert u\rvert$, and
$\mu_b(u,v)=0$ if $u=0$ or $v=0$.  Then $\mu_b(u,v)=\lvert u\rvert L_{\hat u}(v)=\lvert v\rvert R_{\hat v}(u)$, so $\mu_b$ is linear in each
variable; it is norm-multiplicative; and $\mu_b(1,v)=\bar b\,T(b,b,bv)=v$, $\mu_b(u,1)=\bar b\,T(bu,b,b)=u$ by the Mal'cev identities.
Finally $b\,\mu_b(\bar ba,\bar bx)=b(\bar b\,T(a,b,x))=T(a,b,x)$.  Conversely, the formula gives the Mal'cev identities, and each transport is a
composite of $\ell_{\bar b}$, a map $v\mapsto\mu_b(u,v)$ or $u\mapsto\mu_b(u,v)$ with a fixed unit argument (an isometry) and $\ell_b$.  A
bilinear map is determined, uniformly, by its values on pairs of unit vectors, so $b\mapsto\mu_b$ is continuous if and only if
$(b,u,v)\mapsto\bar b\,T(bu,b,bv)$ is, that is, if and only if $T$ is.  By (H1), for unit $v$, $\mu_w(v,v)=2\ip v1\,v-1$; with $v=\bar wx$ and
$\ip{\bar wx}1=\ip xw$ this gives $T(x,w,x)=w\bigl(2\ip xw\,\bar wx-1\bigr)=2\ip xw\,x-w$, and $(\Omega3)$ follows as in
Example~\ref{ex:models}.  On a great circle write $\mu_c=\mu^\phi$ by Lemma~\ref{lem:hurwitz}; $\phi$ maps the circle group $\{e^{tY}\}$,
$Y=\phi^{-1}X$, onto $\{e^{tX}\}$ by $e^{tY}\mapsto e^{tX}$, so $\mu_c(e^{rX},e^{tX})=e^{(r+t)X}$, which gives the formula for base
point $c$; any other point $ce^{sX}$ of $C$ may serve as $c$, since $(ce^{sX})e^{uX}=ce^{(s+u)X}$ by Artin's theorem.  The subsystem generated by two points lies on a great circle through them.
\end{proof}

The first step of this proof works on a round sphere of any dimension and gives a dimension theorem.

\begin{proposition}[dimension; an elementary consequence of Hurwitz's theorem]\label{prop:dim}
Let $V$ be a Euclidean space of dimension $m\ge1$, $S(V)$ its unit sphere with $\kappa(x,y)=\lVert\tfrac12(x+y)\rVert^2$, and
$T\colon S(V)^3\to S(V)$ any map, not assumed continuous, satisfying the Mal'cev identities $T(a,a,x)=x$, $T(a,x,x)=a$, such that $\kappa$
is invariant under all left transports $x\mapsto T(a,b,x)$ and all right transports $x\mapsto T(x,a,b)$.  Then $m\in\{1,2,4,8\}$, that is,
$S(V)$ is $S^0$, $S^1$, $S^3$ or $S^7$.  For every $b\in S(V)$ the map $B_b(a,x)=\lvert a\rvert\lvert x\rvert\,T(\hat a,b,\hat x)$ (and $0$ if $a=0$ or
$x=0$) is a unital composition product on $V$ with unit $b$, and $T(a,b,x)=B_b(a,x)$.
\end{proposition}
\begin{proof}
Fix $b$.  For unit $a$ the map $T(a,b,\cdot)$ preserves $\kappa$, hence inner products, so it is the restriction of an orthogonal map
$L_a$ (first paragraph of the proof of Lemma~\ref{lem:normal}; the argument uses only an orthonormal basis); likewise $T(\cdot,b,x)$ is the
restriction of an orthogonal map $R_x$ for unit $x$.  Then $B_b(a,x)=\lvert a\rvert L_{\hat a}(x)=\lvert x\rvert R_{\hat x}(a)$ is linear in each variable,
$\lvert B_b(a,x)\rvert=\lvert a\rvert\lvert x\rvert$, and $B_b(b,x)=T(b,b,x)=x$, $B_b(a,b)=T(a,b,b)=a$.  So $(V,B_b)$ is a unital composition algebra
with positive definite norm, and Hurwitz's theorem (H2) gives $m\in\{1,2,4,8\}$.
\end{proof}

\begin{remark}\label{rem:dim}
(i) Compared with Proposition~\ref{prop:moufang}(ii), which needs a continuous operation and Adams' theorem, this uses no topology,
at the price of the metric hypotheses (the round sphere, the standard $\kappa$, invariance under left and right transports).  Neither
$(\Omega3)$ nor any form of the third identity is assumed; for $m=4,8$ they follow from Lemma~\ref{lem:normal} (after an isometry $V\cong\HH$
or $\mathbb O$): $(\Omega3)$ holds and the third identity holds on subsystems generated by two elements.  (ii) For $m=2$ the product $B_b$ is
unique: if $f\perp b$ is a unit vector, (H1) gives $B_b(f,f)=-b$, and bilinearity with the unit $b$ determines $B_b$; it is complex
multiplication with unit $b$.  Identifying $V$ isometrically with $\C$ (for either orientation, by commutativity),
$B_b(a,x)=a\bar bx$.  So on $S^1$ the operation is the heap $a\bar bx$, without any continuity hypothesis.  For $m=1$, bilinearity and
$B_b(b,b)=b$ give $B_b(a,x)=abx$, the heap of $S^0=\{\pm1\}$.  For $m=4$ and $m=8$ the products $B_b$ vary with $b$ as described in
Lemma~\ref{lem:normal} and Proposition~\ref{prop:S7} ($B_b(a,x)=b\,\mu_b(\bar ba,\bar bx)$).  (iii) Summary for round spheres with the standard
$\kappa$, the Mal'cev identities and invariance under left and right transports:
\begin{center}\small
\begin{tabular}{@{}lll@{}}
\toprule
sphere & solutions & selected by\\
\midrule
$S^0$, $S^1$ & only the heap $a\bar bx$ & no further hypothesis\\
$S^3$ & $a\bar bx$ or $x\bar ba$ (isomorphic) & continuity, or $(\mathrm C_o)$ at one point (Thm.~\ref{thm:rigid}, Prop.~\ref{prop:C})\\
$S^7$ & continuous maps $S^7\to\OO(7)/G_2$ & $(\mathrm C_o)$ selects $(a\bar b)x$ up to isometry (Props.~\ref{prop:S7}, \ref{prop:C})\\
other & none & (Prop.~\ref{prop:dim})\\
\bottomrule
\end{tabular}
\end{center}
Without continuity, the solutions on $S^3$ and $S^7$ are the maps $b\mapsto\mu_b$ into the same sets, chosen arbitrarily.
\end{remark}

\begin{theorem}[rigidity on $S^3$; an elementary consequence of Hurwitz's theorem]\label{thm:rigid}
Let $T\colon(S^3)^3\to S^3$ be continuous and satisfy the Mal'cev identities, and suppose that the standard comparison $\kappa$ is invariant
under all left transports $x\mapsto T(a,b,x)$ and all right transports $x\mapsto T(x,a,b)$.  Then either $T(a,b,x)=a\bar bx$ for all
$a,b,x$, or $T(a,b,x)=x\bar ba$ for all $a,b,x$.  In both cases $(S^3,T,\kappa)$ is isomorphic to the structure of
Example~\ref{ex:models} (in the second case via $x\mapsto\bar x$).
\end{theorem}
\begin{proof}
By Lemmas~\ref{lem:normal} and \ref{lem:hurwitz}, for each $b$ either $\mu_b(u,v)=uv$, and then $T(a,b,x)=b(\bar ba)(\bar bx)=a\bar bx$ for
all $a,x$, or $\mu_b(u,v)=vu$, and then $T(a,b,x)=x\bar ba$ for all $a,x$.  The two sets of $b$ are disjoint ($\HH$ is not commutative),
closed ($b\mapsto\mu_b$ is continuous) and cover the connected space $S^3$, so one of them is empty.  Finally $\overline{\bar x\,b\,\bar a}=a\bar bx$,
so $x\mapsto\bar x$, which preserves $\kappa$, carries $x\bar ba$ to $a\bar bx$.
\end{proof}

The hypotheses of the earlier version of Theorem~\ref{thm:rigid} ($C^1$, $(\Omega3)$, the third identity on two-generated
subsystems) are not needed: by Lemma~\ref{lem:normal} the last two hold for every $T$ of the theorem.  The earlier proof
differentiated the two normal forms $a\,e^{t\,\mathrm{Rot}_{A,\vartheta}X}$ and $e^{s\,\mathrm{Rot}_{X,\phi}A}e^{tX}$ of $T(a,1,x)$,
$a=e^{sA}$, $x=e^{tX}$, $A\perp X$, which agree exactly for $(\vartheta,\phi)=(0,0)$ and $(-2s,2t)$, the two branches $ax$ and $xa$;
this is kept as a numerical cross-check in the accompanying script.  Continuity cannot be dropped.  For any set $E\subset S^3$ let
$T_E(a,b,x)=a\bar bx$ for $b\in E$ and $T_E(a,b,x)=x\bar ba$ for $b\notin E$.  This is of the form of Lemma~\ref{lem:normal} with
$\mu_b(u,v)=uv$ or $vu$, so it satisfies all algebraic hypotheses, $(\Omega3)$ and the third identity on two-generated subsystems; for
$\emptyset\ne E\ne S^3$ it is discontinuous at every boundary point $b$ of $E$, where the two branches differ by
$\max_{u,v}\lvert uv-vu\rvert=2$.

\begin{proposition}[$S^7$]\label{prop:S7}
Let $\mathcal P$ be the set of unital composition products on $\mathbb O$, a subset of the space of bilinear maps.
\begin{enumerate}[label=(\roman*)]
\item The continuous maps $T\colon(S^7)^3\to S^7$ that satisfy the Mal'cev identities and leave $\kappa$ invariant under all left and
right transports are exactly $T(a,b,x)=b\,\mu_b(\bar ba,\bar bx)$ with $b\mapsto\mu_b$ a continuous map $S^7\to\mathcal P$.  All of them
satisfy $(\Omega3)$ and the third identity on subsystems generated by two elements.
\item $\phi\mapsto\mu^\phi$, $\mu^\phi(u,v)=\phi(\phi^{-1}u\,\phi^{-1}v)$, identifies $\OO(7)/G_2$ with $\mathcal P$, where $G_2=\Aut(\mathbb O)\subset\SO(7)$
is connected of dimension $14$.  So $\mathcal P$ is a compact manifold of dimension $7$ with two components $\{\mu^\phi:\det\phi=\pm1\}$,
each diffeomorphic to $\SO(7)/G_2\cong\R P^7$, exchanged by $\mu\mapsto\mu^{\mathrm{op}}$, $\mu^{\mathrm{op}}(u,v)=\mu(v,u)$.
\item Let $R_\theta$ be the rotation by $\theta$ in a $2$-plane $P\subset\im\mathbb O$, fixing $P^\perp$.  Then $\mu^{R_\theta}=\mu^{R_{\theta'}}$ if and only
if $\theta-\theta'\in2\pi\Z$.  Hence for every continuous $\theta\colon S^7\to\R$ the map $T_\theta(a,b,x)=b\,\mu^{R_{\theta(b)}}(\bar ba,\bar bx)$ is
a solution, and $T_\theta=T_{\theta'}$ only if $\theta-\theta'\in2\pi\Z$ pointwise: the solutions form an infinite-dimensional family.
\item The octonionic operation $(a\bar b)x$ is the solution with $\mu_b(u,v)=\bar b\bigl((bu\bar b)(bv)\bigr)$.
\end{enumerate}
\end{proposition}
\begin{proof}
(i) is Lemma~\ref{lem:normal}.  (ii) Surjectivity is Lemma~\ref{lem:hurwitz}, and $\mu^\phi=\mu^{\phi'}$ if and only if $\phi^{-1}\phi'$ is an
automorphism; automorphisms fix $1$ and are isometric (the argument in the proof of Lemma~\ref{lem:hurwitz}), so
$\Aut(\mathbb O)\subset\OO(7)$.  It acts simply transitively on basic triples, the triples $(e_1,e_2,e_3)$ of unit imaginary octonions with
$e_2\perp\operatorname{span}(1,e_1)$ and $e_3\perp\operatorname{span}(1,e_1,e_2,e_1e_2)$ \cite[\S4.1]{Baez02}; these form a connected space (choose
$e_1$ in a $6$-sphere, then $e_2$ in a $5$-sphere, then $e_3$ in a $3$-sphere), so $G_2$ is homeomorphic to it, connected, of dimension
$6+5+3=14$, hence contained in $\SO(7)$.  So $\OO(7)/G_2$ has dimension $21-14=7$ and two components.  Conjugation $c$ is
$-\mathrm{id}$ on $\im\mathbb O$, of determinant $-1$, and reverses products, so $\mu^{\phi\circ c}=(\mu^\phi)^{\mathrm{op}}$.  Finally
$\mathrm{Spin}(7)/G_2=S^7$ \cite[\S4.1]{Baez02}, and $-1\in\mathrm{Spin}(7)$ acts as $-1$ on $S^7$, so $\SO(7)/G_2=S^7/\{\pm1\}=\R P^7$.  (iii) $\mu^{R_\theta}=\mu^{R_{\theta'}}$
if and only if $R_{\theta'-\theta}\in G_2$.  An automorphism fixing the five-dimensional space $P^\perp$ pointwise fixes the subalgebra generated
by $1$ and $P^\perp$; a subalgebra of $\mathbb O$ containing $1$ is a composition algebra under the restricted norm, so by (H2) its dimension
is $1$, $2$, $4$ or $8$, here $8$, and the automorphism is the identity.  The rest follows from (i), since $T$ determines $\mu_b$.
(iv) $\bar b\,T(bu,b,bv)=\bar b\bigl(((bu)\bar b)(bv)\bigr)$, and $(bu)\bar b=b(u\bar b)$ by flexibility.
\end{proof}

The products $\mu_b$ of (iv) depend on $b$: numerically, $\sup_{\lvert u\rvert=\lvert v\rvert=1}\lvert\mu_b(u,v)-uv\rvert=4\lvert\re b\rvert\,\lvert\im b\rvert$
(checked to rounding error in the accompanying script; we have not proved it), which vanishes exactly for real and for purely imaginary $b$.  So the question left open in the earlier version, whether the solutions on $S^7$ are the operations
obtained from $(a\bar b)x$ by symmetries, has a negative answer.  The family of Proposition~\ref{prop:S7}(iii) satisfies the round-trip
law and the third identity on two-generated subsystems, so neither selects $(a\bar b)x$.  Two further natural conditions do not select
it either.  Every base loop $x\cdot_oy=T(x,o,y)$ of every solution is isomorphic to the Moufang loop $S^7$ (by
$\psi(p)=o\,\phi(p)$ for $\mu_o=\mu^\phi$, since $\psi^{-1}(x\cdot_oy)=\psi^{-1}(x)\,\psi^{-1}(y)$), so the Moufang identities for all base loops hold
throughout the family; and the left transports $x\mapsto(a\bar b)x$ are not automorphisms of $T(a,b,x)=(a\bar b)x$, whereas they are
automorphisms of a heap: in $\mathbb O=\HH\oplus\HH\ell$ with $(a+b\ell)(c+d\ell)=(ac-\bar db)+(da+b\bar c)\ell$, for $g=(1+i)/\sqrt2$ the points $T(g,gj,g\ell)$ and $g\,T(1,j,\ell)$ are
antipodal.  The following identity does select it.

\begin{proposition}[the identity that selects $(a\bar b)x$]\label{prop:C}
Let $T\colon S(A)^3\to S(A)$ satisfy the Mal'cev identities and two-sided invariance of $\kappa$; no continuity is assumed.  For
$o\in S(A)$ consider
\[
(\mathrm C_o)\qquad T(x,y,z)=T\bigl(T(x,o,T(o,y,o)),o,z\bigr)\qquad\text{for all }x,y,z.
\]
The following are equivalent: (i) $(\mathrm C_o)$ holds for some $o$; (ii) $(\mathrm C_o)$ holds for every $o$; (iii) there is an orthogonal
map $\psi$ of $A$ with $T(\psi p,\psi q,\psi r)=\psi\bigl((p\bar q)r\bigr)$ for all $p,q,r$, that is, $(S(A),T,\kappa)$ is isomorphic to
$(S(A),(a\bar b)x,\kappa)$.  For $A=\HH$ they are also equivalent to: $T(a,b,x)=a\bar bx$ for all $a,b,x$, or $T(a,b,x)=x\bar ba$ for all
$a,b,x$.
\end{proposition}
\begin{proof}
By Lemma~\ref{lem:normal}, $T(o,y,o)=2\ip oy\,o-y$, so $\bar o\,T(o,y,o)=\overline{\bar oy}$, and $\bar o\,T(x,o,w)=\mu_o(\bar ox,\bar ow)$.  Hence
$(\mathrm C_o)$ says $T(x,y,z)=o\,\mu_o\bigl(\mu_o(\bar ox,\overline{\bar oy}),\bar oz\bigr)$.  Write $\mu_o=\mu^\phi$ (Lemma~\ref{lem:hurwitz}); $\phi$ fixes $1$
and preserves $\im A$, so it commutes with conjugation.  With $\psi(p)=o\,\phi(p)$, an orthogonal map, and $p=\psi^{-1}(x)=\phi^{-1}(\bar ox)$,
$q=\psi^{-1}(y)$, $r=\psi^{-1}(z)$, this reads $T(x,y,z)=o\,\phi\bigl((p\bar q)r\bigr)=\psi\bigl((p\bar q)r\bigr)$; so (i) implies (iii).  For (iii)
$\Rightarrow$ (ii) it suffices, $(\mathrm C_o)$ being an identity in $T$, that $T_0(a,b,x)=(a\bar b)x$ satisfies $(\mathrm C_o)$ for every $o$.  Here
$T_0(o,y,o)=(o\bar y)o$, and the right Moufang identity $((ab)c)b=a(b(cb))$ of alternative algebras \cite{Schafer66} with $a=x\bar o$, $b=o$,
$c=\bar y$ gives $(x\bar o)\bigl((o\bar y)o\bigr)=(((x\bar o)o)\bar y)o=(x\bar y)o$, using $(w\bar o)o=w=(wo)\bar o$ for unit $o$; hence
$T_0(T_0(x,o,T_0(o,y,o)),o,z)=(((x\bar y)o)\bar o)z=(x\bar y)z$.  (ii) $\Rightarrow$ (i) is trivial.  For $A=\HH$, $\mu_o$ is $uv$ or $vu$, and
the displayed form of $(\mathrm C_o)$ gives $o(\bar ox)(\bar yo)(\bar oz)=x\bar yz$ or $o(\bar oz)(\bar yo)(\bar ox)=z\bar yx$ for all $x,y,z$;
both satisfy $(\mathrm C_o)$.
\end{proof}

For a Moufang loop the ternary operation $(xy^{-1})z$ is a \emph{ternary Moufang loop}, and the ternary operation is recovered from the
loop at every origin $e$ with $y^{-1}=(eye)$ \cite[Lemma~3.4 and Remark~(3)]{BK12}; that is $(\mathrm C_o)$ for all $o$, and Proposition~\ref{prop:C}
says that under the metric hypotheses one origin suffices and only the unit octonions occur.  $(\mathrm C_o)$ is the identity
$\hp xyz=(x\cdot y^{-1})\cdot z$ assumed in Proposition~\ref{prop:moufangclass}, for the base loop
$x\cdot y=T(x,o,y)$, with the inverse $y^{-1}=T(o,y,o)$ (a two-sided inverse in that loop, since $\mu^\phi(v,\bar v)=\lvert v\rvert^2$).  On $S^3$,
Proposition~\ref{prop:C} shows that $(\mathrm C_o)$ at a single point replaces continuity in Theorem~\ref{thm:rigid}; the maps $T_E$ above
violate it.  On $S^7$, $(\mathrm C_o)$ for all $o$ is an origin-free identity which, together with the Mal'cev identities and two-sided
invariance, determines $(a\bar b)x$ up to isometry.


\subsection{Roundness without associativity}\label{sec:round}
The results above assumed the round sphere and the standard comparison.  These are consequences of the comparison axioms
alone, with no form of the third heap identity: it suffices that a binary operation with a two-sided unit at one point has
$\kappa$-preserving translations on both sides and satisfies the round-trip law at that point.  The argument is metric and
never uses a group structure on $M$; it uses the Lie theory of the group of $\kappa$-preserving homeomorphisms
(Lemmas~\ref{lem:roundG} and \ref{lem:roundLie}).  Remark~\ref{rem:roundaudit} records where each hypothesis enters.

\begin{theorem}[roundness]\label{thm:round}
Let $M$ be a connected locally Euclidean space of dimension $n\ge1$ without boundary (it is then Hausdorff and second countable,
Lemma~\ref{lem:roundG}), and let
$\kappa\colon M\times M\to[0,1]$ be continuous and symmetric, with $\kappa(x,y)=1$ if and only if $x=y$.  Let $w\in M$ and let
$B\colon M\times M\to M$ satisfy $B(w,x)=x=B(x,w)$ for all $x$, such that for every $a$ the maps $L_a=B(a,\cdot)$ and
$R_a=B(\cdot,a)$ are continuous and preserve $\kappa$.  Suppose
\[
\kappa\bigl(w,B(x,x)\bigr)=\bigl(2\kappa(w,x)-1\bigr)^2\qquad(x\in M),
\]
and that at most one $z\in M$ satisfies $\kappa(w,z)=0$.  Then $M$ is compact, $n\in\{1,3,7\}$, and there is a homeomorphism
$F\colon S^n\to M$ with $\kappa(Fu,Fv)=\lVert\tfrac12(u+v)\rVert^2$ for all $u,v\in S^n$.  In particular $c=2\kappa-1$ is a positive
semidefinite kernel of rank $n+1$, $\arccos c$ is a metric, and every point has exactly one zero partner.
\end{theorem}

\begin{corollary}\label{cor:round}
Let $M$ be a connected locally Euclidean space of dimension $\ge1$ without boundary, $T\colon M^3\to M$ a map with the Mal'cev identities
$T(a,a,x)=x$, $T(a,x,x)=a$, continuous in each of its outer arguments, and $\kappa$ continuous with $(\Omega2\mathrm{a,b,c,d})$ and
$(\Omega3)$.  No form of the third identity is assumed.  Then, after the homeomorphism $F$ of Theorem~\ref{thm:round}, $M$ is the
round $S^1$, $S^3$ or $S^7$ with the standard $\kappa$, so Lemma~\ref{lem:normal} applies: the third identity holds on every subsystem
generated by two elements; on $S^1$, $T$ is the heap $a\bar bx$ (Remark~\ref{rem:dim}(ii)); on $S^3$ it is $a\bar bx$ or $x\bar ba$ if $T$ is continuous
(Theorem~\ref{thm:rigid}) or if $(\mathrm C_o)$ holds at one point (Proposition~\ref{prop:C}); on $S^7$ the operations are described by
Lemma~\ref{lem:normal} and Proposition~\ref{prop:S7}, and $(\mathrm C_o)$ selects $(a\bar b)x$ up to isometry.
\end{corollary}
\begin{proof}
Apply Theorem~\ref{thm:round} with any $w$ and $B(a,x)=T(a,w,x)$; the transports $L_a$, $R_x$ are left and right transports of $T$, and
$B(x,x)=T(x,w,x)$.  Transport $T$ and $\kappa$ to $S^n$ by $F$ and apply the results cited.
\end{proof}

Throughout the proof of Theorem~\ref{thm:round} write $c=2\kappa-1$ and $q=1-c$; $q$ is continuous, symmetric, and vanishes exactly on the
diagonal.

\begin{lemma}\label{lem:roundG}
Under the hypotheses of Theorem~\ref{thm:round}:
\begin{enumerate}[label=(\roman*)]
\item every continuous $\kappa$-preserving injection $f\colon M\to M$ is a homeomorphism of $M$; in particular all $L_a$, $R_a$ are;
\item the group $G$ of $\kappa$-preserving homeomorphisms of $M$, with the compact-open topology, is a locally compact Hausdorff
topological group acting continuously and transitively on $M$; the stabiliser $H=G_w$ is compact; the maps $a\mapsto L_a$ and
$a\mapsto R_a$ from $M$ to $G$ are continuous; and $G/H\to M$, $gH\mapsto gw$, is a homeomorphism.
\end{enumerate}
\end{lemma}
\begin{proof}
$M$ is Hausdorff: in a locally Euclidean space sequences determine the topology, and if $z_j\to x$ and $z_j\to y$, then
$\kappa(x,y)=\lim\kappa(z_j,z_j)=1$, so $x=y$.  Second countability is not used below before it is proved at the end of (ii).

(i) $f$ is injective, and open by invariance of domain \cite[Thm.~2B.3]{Hatcher}.  Let $y$ be in the closure of $f(M)$.  Fix a coordinate ball $U$
about $w$ with compact closure in the chart and put $\epsilon=\min_{z\in\partial U}q(w,z)>0$.  For $x\in M$ the map $j=f\circ R_x$ is a continuous
injection with $j(w)=f(x)$; it maps $\bar U$ homeomorphically onto a compact set and $U$ onto an open set with boundary $j(\partial U)$, on which
$q(f(x),\cdot)=q(w,\cdot)\ge\epsilon$.  Choose a path-connected neighbourhood $V$ of $y$ with $q<\epsilon$ on $V\times V$ and $x$ with $f(x)\in V$.  A path in
$V$ from $f(x)$ to $y$ cannot meet $j(\partial U)$, so it stays in $j(U)\subset f(M)$, and $y\in f(M)$.  Thus $f(M)$ is open, closed and nonempty.

(ii) Transitivity: $R_x(w)=x$.  Put $\delta(x,y)=\inf_\sigma\max_{s,t}q(\sigma(s),\sigma(t))$ over paths $\sigma$ from $x$ to $y$; $\delta$ is symmetric and
$G$-invariant, and (a) $x_j,y_j\to x$ implies $\delta(x_j,y_j)\to0$ (paths in small coordinate balls), (b) $u_j\to u$ and $\delta(u_j,v_j)\to0$
imply $v_j\to u$ (for a small coordinate ball $V$ about $u$, $q(u_j,\cdot)$ is bounded below on $\partial V$ uniformly in large $j$, and a path of
smaller $q$-diameter cannot leave $V$).  Let $f_0\colon M\to[0,1]$ be continuous with compact support $K_0$ and unique maximum $f_0(w)=1$, and put
$\rho(x,y)=\sup_{h\in G}\lvert f_0(hx)-f_0(hy)\rvert$.  This is a $G$-invariant metric: for $x\ne y$ take $h$ with $hx=w$.  It induces the topology: if
$y_j\to x$ and $\rho(x,y_j)\ge\eta>0$ along a subsequence, choose $h_j$ realising this up to $\eta/2$; one of $h_jx$, $h_jy_j$ lies in $K_0$ and
converges along a further subsequence, the other converges to the same point by (a) and (b), contradicting continuity of $f_0$; conversely
$\rho(x,y_j)\to0$ gives $f_0(hy_j)\to1$ for $hx=w$, so $hy_j\to w$ and $y_j\to x$.  Since $\rho(x,y)\le\frac12$ forces $f_0(hy)\ge\frac12$, closed
$\rho$-balls of radius $\frac12$ are compact, and $\rho$ is complete.  So $M$ is second countable: starting from $\{w\}$, the
$\frac14$-neighbourhood of a compact set lies in finitely many closed $\frac12$-balls, and since every point is reached by a chain of steps
below $\frac14$, $M$ is a countable union of compact sets, hence a separable metric space.  Elements of $G$ are $\rho$-isometries, so on $G$ the compact-open topology is
that of uniform convergence on compact sets, and pointwise convergence on a countable dense set implies it; composition and inversion
are continuous ($\rho(h_j^{-1}x,h^{-1}x)=\rho(x,h_jh^{-1}x)$), and so is the action.  The set $N=\{h:\rho(hw,w)\le\frac12\}$ is a neighbourhood of
the identity.  It is compact: for $h_j\in N$, $h_jw$ and $h_j^{-1}w$ lie in a compact ball; joining $w$ to any $x$ by a chain with steps below
$\frac14$ and covering at each step the $\frac14$-neighbourhood of a compact set by finitely many closed $\frac12$-balls shows that $h_jx$ and
$h_j^{-1}x$ stay in compact sets; a diagonal subsequence converges on a countable dense set, the limits extend to isometries $h,k$ with
$h_j\to h$, $h_j^{-1}\to k$ uniformly on compact sets, $hk=kh=\mathrm{id}$, and $h$ preserves $\kappa$ by continuity.  So $G$ is locally
compact, and $H\subset N$ is closed, hence compact.  The maps $a\mapsto L_a$ converge pointwise as $a\to a_0$ ($L_a(x)=R_x(a)$ is
continuous in $a$), hence in $G$; likewise for $R$.  As the action of $G$ is continuous, $(a,x)\mapsto L_a(x)=B(a,x)$ is jointly continuous.
Finally $a\mapsto L_aH$ is a continuous inverse of $G/H\to M$, because
$L_a(w)=a$.
\end{proof}

\begin{lemma}\label{lem:roundLie}
$G$ is a Lie group.  $M$ carries a smooth structure compatible with its topology for which $G$ acts smoothly, a $G$-invariant
Riemannian metric $g$, and the maps $a\mapsto L_a$, $a\mapsto R_a$ from $M$ to $G$ are smooth.
\end{lemma}
\begin{proof}
$G$ is locally compact, $H$ closed, $G/H\cong M$ connected and locally contractible, and $G$ acts faithfully; by Hofmann--Kramer
\cite[Thm.~5.4]{HK15} (whose erratum \cite{HK15e} corrects the proof of a lemma, not the statement) $G$ is a Lie group.  Then $H$ is a compact
Lie subgroup, $G/H$ is a smooth manifold on which $G$ acts smoothly and $G\to G/H$ has smooth local sections \cite[Ch.~21]{Lee13}, and
averaging an inner product on $T_wM$ over $H$ gives a $G$-invariant metric $g$.  An element of the Lie algebra whose fundamental field
vanishes at every point generates a one-parameter subgroup acting trivially, hence is $0$ by faithfulness; so there are finitely many
points $x_1,\dots,x_N$ for which $E(h)=(hx_1,\dots,hx_N)$ has injective differential at $e$, hence, by equivariance, everywhere.  Near
$L_{a_0}$, $E$ is an embedding; $a\mapsto L_a$ is continuous, and $E(L_a)=(R_{x_1}(a),\dots,R_{x_N}(a))$ is smooth because each $R_{x_i}\in G$
acts smoothly; so $a\mapsto L_a$ is smooth.  The same argument with $E(R_x)=(L_{x_i}(x))_i$ applies to $R$.
\end{proof}

\begin{lemma}\label{lem:roundgeo}
For $v\in T_wM$ let $X_v(x)=d(R_x)_wv$ and $Y_v(x)=d(L_x)_wv$, and for a unit vector $v$ let $\gamma_v$ be the geodesic with $\gamma_v(0)=w$,
$\dot\gamma_v(0)=v$.  Then:
\begin{enumerate}[label=(\roman*)]
\item $X_v$, $Y_v$ are Killing fields of constant length $\lvert v\rvert$, their integral curves are geodesics, $g$ is complete, and $\gamma_v$
is the integral curve of both $X_v$ and $Y_v$ through $w$;
\item $B(\gamma_v(s),\gamma_v(t))=\gamma_v(s+t)$ for all $s,t$;
\item there is $a>0$, independent of $v$, with $c(w,\gamma_v(t))=\cos(\pi t/a)$ for all $t$; the points $\gamma_v(a)$ coincide, and this point $w^\perp$
is the unique zero partner of $w$;
\item $M$ is compact, every segment $\gamma_v|_{[0,a]}$ is minimising, and $c(x,y)=\cos(k\,d_g(x,y))$ for all $x,y$, with $k=\pi/a$.
\end{enumerate}
\end{lemma}
\begin{proof}
(i) By Lemma~\ref{lem:roundLie}, $s\mapsto L_{\alpha(s)}$ for a curve $\alpha$ with $\alpha(0)=w$, $\dot\alpha(0)=v$ is a smooth curve in $G$ through
$L_w=\mathrm{id}$; its velocity is a Lie-algebra element whose fundamental field is $x\mapsto\frac d{ds}L_{\alpha(s)}(x)=\frac d{ds}R_x(\alpha(s))=X_v(x)$.
So $X_v$ is a complete Killing field, and $\lvert X_v(x)\rvert=\lvert v\rvert$ because $R_x$ is an isometry; likewise for $Y_v$.  For a Killing field $X$ of
constant length, $\langle\nabla_XX,Z\rangle=-\langle\nabla_ZX,X\rangle=-\frac12Z\lvert X\rvert^2=0$, so its integral curves are geodesics.  The values
$X_v(x)$, $v\in T_wM$, fill $T_xM$, so every geodesic is an integral curve of some $X_v$ and is defined for all time.  At $w$, $X_v(w)=Y_v(w)=v$.
(ii) $t\mapsto L_{\gamma_v(s)}\gamma_v(t)$ is a geodesic starting at $\gamma_v(s)$ with velocity $Y_v(\gamma_v(s))=\dot\gamma_v(s)$.
(iii) Put $p(t)=c(w,\gamma_v(t))$.  Applying $L_{\gamma_v(-t)}$ and symmetry gives $p(-t)=p(t)$; the round-trip law with $B(x,x)=\gamma_v(2t)$ for
$x=\gamma_v(t)$ gives $p(2t)=2p(t)^2-1$.  By (ii) the return times $\{t:\gamma_v(t)=w\}$ form a closed subgroup of $\R$, discrete since $v\ne0$.
If it were $\{0\}$, then $p<1$ on $(0,\infty)$ by faithfulness, contradicting Lemma~\ref{lem:noline}.  So it is $P\Z$ with $P>0$; $p$ has
period $P$ and $p=1$ exactly on $P\Z$, and Lemma~\ref{lem:circle} applied to $\theta\mapsto p(P\theta/2\pi)$ gives $p(t)=\cos(2\pi t/P)$.
Put $a(v)=P/2$; then $c(w,\gamma_v(a(v)))=-1$, and by the hypothesis on zero partners all $\gamma_v(a(v))$ coincide.  The function $a$ is
continuous ($a(v)=\pi\eta/\arccos c(w,\gamma_v(\eta))$ for a fixed $\eta$ below the injectivity radius at $w$; indeed
$\exp_w(a(v)v)=w^\perp=\exp_w(-a(v)v)$, so $a(v)$ is at least the injectivity radius, and $\pi\eta/a(v)\in(0,\pi)$) and, by the implicit function
theorem applied to a function vanishing at $w^\perp$ with nonzero derivative along $\dot\gamma_{v_0}(a(v_0))$, smooth.  For a curve $v(s)$ of
unit vectors, $J=\partial_s\gamma_{v(s)}(t)$ satisfies $\partial_t\langle J,\dot\gamma\rangle=\frac12\partial_s\lvert\dot\gamma_{v(s)}\rvert^2=0$ and $J(0)=0$, so
$J\perp\dot\gamma$; differentiating $\gamma_{v(s)}(a(v(s)))=w^\perp$ gives $J(a)+a'(0)\dot\gamma(a)=0$, hence $a'(0)=0$.  The unit sphere of $T_wM$ is
connected for $n\ge2$, and for $n=1$, $\gamma_{-v}(t)=\gamma_v(-t)$; so $a$ is constant.  Elements of $G$ carry this to every base point.
(iv) By the Hopf--Rinow theorem \cite[Ch.~6]{Lee18} any two points are joined by a minimising unit-speed geodesic.  By periodicity and
$\gamma_{-v}(t)=\gamma_v(-t)$ every point $\gamma_v(t)$ is also $\gamma_{\pm v}(t')$ with $0\le t'\le a$, so all distances are at most $a$; a minimising
geodesic from $w$ to $w^\perp$ has length $t$ with $\cos(\pi t/a)=-1$, i.e.\ $t=a$.  Hence each $\gamma_v|_{[0,a]}$ minimises, $d_g(w,\gamma_v(t))=t$
for $0\le t\le a$, $M=\exp_w(\{\lvert u\rvert\le a\})$ is compact, and $c(w,x)=\cos(kd_g(w,x))$; transitivity of $G$ gives the formula at every point.
\end{proof}

\begin{lemma}\label{lem:roundcurv}
If $n\ge2$, $(M,g)$ has constant sectional curvature $k^2$.
\end{lemma}
\begin{proof}
Fix a unit $v\in T_wM$, put $X=X_v$, $\gamma=\gamma_v$, let $\Phi_t$ be the flow of $X$ (isometries), $F_t=d(\Phi_t)_w$ and $A_0Z=\nabla_ZX$ at $w$.  By
the Killing equation $A_0$ is skew-adjoint, and $A_0v=\nabla_XX(w)=0$.  Since $\Phi_t$ preserves $X$ and $\nabla$, $\nabla X$ at $\gamma(t)$ is $F_tA_0F_t^{-1}$.
For a curve $\sigma$ through $w$ with velocity $Z_0$, $(s,t)\mapsto\Phi_t(\sigma(s))$ is a variation through integral curves of $X$, i.e.\ through
geodesics, with variation field $Z(t)=F_tZ_0$, a Jacobi field; and $\nabla_tZ=\nabla_ZX=F_tA_0Z_0$, so $\nabla_t^2Z=F_tA_0^2Z_0$.  The Jacobi equation
$\nabla_t^2Z+R(Z,\dot\gamma)\dot\gamma=0$ gives $R(\cdot,\dot\gamma)\dot\gamma=-F_tA_0^2F_t^{-1}$ along $\gamma$, and $P_t=F_te^{-tA_0}$ is parallel transport
(differentiate), so in a parallel frame this operator is the constant $-A_0^2\ge0$ (compare \cite[Prop.~1, Thm.~5]{BN09}).  Let $z\perp v$ be a unit
eigenvector of $-A_0^2$ with eigenvalue $\lambda^2$.  The variation $\gamma_{v(s)}$ with $v(s)$ unit, $v'(0)=z$, has the Jacobi field $J(t)=j(t)P_tz$
with $j''=-\lambda^2j$, $j(0)=0$, $j'(0)=1$, and $J(a)=0$ because all $\gamma_{v(s)}(a)=w^\perp$; so $\lambda=m\pi/a$ with an integer $m\ge1$.  The
segment $\gamma|_{[0,a]}$ minimises length and has constant speed, so it minimises energy among curves with the same ends, and the second
variation \cite[Ch.~10]{Lee18} along $W(t)=\sin(\pi t/a)P_tz$ is non-negative:
\[
0\le\int_0^a\bigl(\lvert\nabla_tW\rvert^2-\langle R(W,\dot\gamma)\dot\gamma,W\rangle\bigr)dt=\frac a2\Bigl(\frac{\pi^2}{a^2}-\lambda^2\Bigr).
\]
Hence $m=1$, $-A_0^2=k^2$ on $v^\perp$, and every plane containing $v$ has curvature $k^2$.  As $v$ is arbitrary and $G$ is transitive, the
curvature is $k^2$ everywhere.
\end{proof}

\begin{proof}[Proof of Theorem~\ref{thm:round}]
For $n=1$, $M$ is a compact connected $1$-manifold, a circle, and $\gamma_v$ is injective on $[0,2a)$ with period $2a$, so $(M,g)$ is a
circle of length $2a$.  For $n\ge2$, by Lemma~\ref{lem:roundcurv} and the Killing--Hopf theorem \cite[Ch.~12]{Lee18}, $(M,g)$ is
isometric to $S^n(1/k)/\Gamma$ for a finite group $\Gamma$ of isometries acting freely.  The distance from $[x]$ to $[y]$ is $\min_{\phi\in\Gamma}d(x,\phi y)$,
and $d_g(w,w^\perp)=a=\pi/k$ forces $\phi y=-x$ for all $\phi\in\Gamma$, hence $\Gamma=\{1\}$.  In both cases rescaling gives an isometry
$F\colon S^n\to(M,k^2g)$, and $c(Fu,Fv)=\cos d_{S^n}(u,v)=\ip uv$ by Lemma~\ref{lem:roundgeo}(iv), i.e.\ $\kappa(Fu,Fv)=\lVert\frac12(u+v)\rVert^2$.
The operation $F^{-1}B(F\cdot,F\cdot)$ on $S^n$ has two-sided unit $F^{-1}w$ and inner-product-preserving translations, so the proof of
Proposition~\ref{prop:dim} gives $n+1\in\{2,4,8\}$.  Finally $\ip uv$ is a Gram kernel of rank $n+1$ on $S^n$.
\end{proof}

\begin{remark}[the hypotheses]\label{rem:roundaudit}
(i) \emph{Unique zero partner.}  $\SO(3)$ with its group heap and $\kappa=\cos^2(\alpha/2)$, $\alpha$ the angle of the relative rotation,
satisfies every other hypothesis (Remark~\ref{rem:quotient}); it is $\R P^3$, the zero set of each point is a projective plane, and $c$ is not
positive semidefinite (the identity and the three coordinate half-turns have $c$-Gram matrix $2I-\mathbf 1\mathbf 1^{\mathsf T}$, with an
eigenvalue $-2$).  (ii) \emph{Round-trip law.}  On the torus $(\R/2\pi\Z)^m$, $m\ge2$, with $B(a,x)=a+x$ and $c(x,y)=\frac1m\sum_i\cos(y_i-x_i)$,
and on $S^1\times\R$ with the group operation and $c((z,s),(z',s'))=e^{-(s'-s)^2}\re(\bar zz')$, all hypotheses except the round-trip law
hold, with a unique zero partner; both fail the law (for the torus at the difference $(\pi,0,\dots,0)$, for the cylinder at a vertical
difference $1$, where the two sides are $e^{-4}$ and $2e^{-2}-1$).  (iii) \emph{Faithfulness.}  $\kappa\equiv1$ on $\R$ with $B(a,x)=a+x$ satisfies
all hypotheses except faithfulness (it has no zero partner, which the hypothesis ``at most one'' allows).  (iv) \emph{Continuity.}  Only
continuity of $B$ in each argument separately is used, at one unit $w$; in Corollary~\ref{cor:round} no continuity in the middle argument
of $T$ is needed: the maps $T_E$ after Theorem~\ref{thm:rigid} are discontinuous in the middle slot and satisfy the conclusion.
(v) \emph{One base point.}  The round-trip law and the zero-partner condition are used only at $w$, and invariance only for the
translations of the single binary operation $B$.  (vi) \emph{Right invariance.}  It cannot be dropped: Proposition~\ref{prop:onesided}
satisfies every other hypothesis, with smooth $B$ and diffeomorphic right translations, and is not round.  Its comparison is continuous
but not $C^2$ on the diagonal, so whether right invariance can be replaced by left invariance together with $(\Omega2\mathrm e)$ is open.
\end{remark}

\begin{remark}[relation to the earlier results]\label{rem:roundrel}
(i) For a comparison structure, Theorem~\ref{thm:round} with $B(a,x)=\hp awx$ gives a second proof of the geometric part of
Theorem~\ref{thm:main}: $M$ is the round $S^1$, $S^3$ or $S^7$ with the standard $\kappa$; the heap identity then excludes $S^7$ (with both
middle arguments equal to $1$ it makes the unital composition product $\mu_1$ of Lemma~\ref{lem:normal} associative, and no isotope of
$\mathbb O$ is associative) and fixes the operation (Theorem~\ref{thm:rigid} and the remark after it).  This route uses no group structure
on $M$ and no regularity of $\kappa$, but uses Hofmann--Kramer, Hurwitz and Killing--Hopf; Theorem~\ref{thm:noreg} remains the result
without $(\Omega2\mathrm d)$, where $\SO(3)$ survives.  (ii) In Proposition~\ref{prop:moufangclass} the comparison is invariant under the loop
translations at $o$ and satisfies $(\Omega2)$, $(\Omega3)$; Theorem~\ref{thm:round} with $B(x,y)=\hp xoy$ already makes $M$ a round sphere, and
the hypothesis $\hp xyz=(x\cdot y^{-1})\cdot z$ is $(\mathrm C_o)$, the loop inverse being $T(o,y,o)$ there.  Proposition~\ref{prop:C} also needs
the Mal'cev identities and invariance under all transports; both hold: $T(a,b,\cdot)$ is the loop translation by $a\cdot b^{-1}$ and $T(\cdot,a,b)$
is the composite $x\mapsto(x\cdot a^{-1})\cdot b$ of two loop translations, and $T(a,a,x)=x$, $T(a,x,x)=a$ are the inverse properties of the
loop, which on the round sphere is isomorphic to the unit sphere of $\C$, $\HH$ or $\mathbb O$ (proof of Lemma~\ref{lem:normal}).  So
Proposition~\ref{prop:C} gives the conclusion of Proposition~\ref{prop:moufangclass} with continuity of the loop translations in place of
analyticity, and without di-associativity.  That proposition and
Proposition~\ref{prop:mfauto} remain as the Lie-theoretic route and as statements about analytic loops.  (iii) Positive semidefiniteness of
$\kappa$ and its rank, derived in Theorem~\ref{thm:main} from the classification, follow here from the metric alone.
\end{remark}


\subsection{One-sided invariance}\label{sec:onesided}
Theorem~\ref{thm:round} assumes that the left and the right translations preserve $\kappa$.  Invariance under the left translations
alone does not suffice for a continuous comparison, even for smooth operations with diffeomorphic right translations.  The example fails
the third identity on a two-generated subsystem, and its comparison is not $C^2$ on the diagonal; so the variants that keep that identity,
or require $(\Omega2\mathrm e)$, remain open (Section~\ref{sec:open}, Q3).

\begin{proposition}[one-sided invariance does not suffice]\label{prop:onesided}
On $S^3\subset\HH$ write $x=s+it+vj$ with $s,t\in\R$, $v\in\C=\operatorname{span}(1,i)$; let $\beta(s)=e^{1-1/s}$ for $s>0$ and $\beta(s)=0$ for $s\le0$,
$\phi(x)=\frac14t\beta(s)$, $h_x=e^{i\phi(x)/2}$, and
\[
B(x,y)=x\,h_x\,y\,h_x^{-1},\qquad S(x)=B(x,x).
\]
Let $\theta$ be the unique continuous fixed point, among continuous $u\colon S^3\to[0,\pi]$ with $u(1)=0$, $u(-1)=\pi$, of
\[
(\mathcal Tu)(x)=\tfrac12u(S(x))\ (s\ge0),\qquad (\mathcal Tu)(x)=\pi-\tfrac12u(S(x))\ (s\le0),
\]
and put $c(a,b)=\cos\theta(a^{-1}b)$, $\kappa=\frac12(1+c)$.  Then:
\begin{enumerate}[label=(\roman*)]
\item $B$ is smooth with two-sided unit $1$; every left translation $L_a$ preserves $\kappa$; every right translation is a diffeomorphism
of $S^3$; $\kappa$ is continuous and symmetric, $\kappa(a,b)=1$ exactly for $a=b$, each point $a$ has exactly one zero partner $-a$, and
$\kappa(1,S(x))=(2\kappa(1,x)-1)^2$.  So all hypotheses of Theorem~\ref{thm:round} hold except invariance under the right translations.
\item $(S^3,\kappa)$ is not round: for $x=(\frac1{\sqrt2},\frac12,\frac12,0)$ one has $\theta(x)<\pi/4$ and $\theta(-x)=3\pi/4$, so the Gram matrix of
$c$ on $(1,-1,x)$ has determinant $-(c(1,x)+c(-1,x))^2<0$, and $\arccos c$ violates the triangle inequality on $1,x,-1$.  No
$\kappa$-preserving map fixing $1$ carries $x$ to $y=e^{i\theta(x)}$, although $\kappa(1,x)=\kappa(1,y)$.
\item The ternary operation $T(a,b,x)=b\,B(\bar ba,\bar bx)$ satisfies the Mal'cev identities, its left transports preserve $\kappa$, and
$(\Omega3)$ holds at every base point; but for $p=(\frac12,\frac12,\frac1{\sqrt2},0)$, $T(T(p,1,p),1,p)\ne T(p,1,T(p,1,p))$, so the third identity
fails on the subsystem generated by $1$ and $p$, and $B$ is not power-associative.
\item $c$ is invariant under the maps $x\mapsto qxe^{i\alpha}$ ($q\in S^3$, $\alpha\in\R$), a group isomorphic to $\mathrm U(2)$ generated by the left
transports and a right circle (a clock in the sense of the companion paper \cite{Companion}), but not under right multiplication by
$e^{j\alpha}$.  $\kappa$ is continuous but not $C^2$ on the diagonal: a second-order expansion at $1$ would force $\theta(e^{\epsilon u})/\epsilon\to1$
for every unit imaginary $u$ (see Remark~\ref{rem:onesidedreg}).
\end{enumerate}
\end{proposition}
\begin{proof}
Write $u=s+it$ and $C(x)=h_xxh_x^{-1}=u+e^{i\phi(x)}vj$ (conjugation by $e^{i\alpha}$ fixes $u$ and multiplies $v$ by $e^{2i\alpha}$, since $jz=\bar zj$).
\emph{Units and translations.}  $h_1=1$ gives $B(1,y)=y$, and $B(x,1)=x$.  Each $L_x$ is a conjugation composed with a left
multiplication.  For $R_y(x)=xh_xyh_x^{-1}$ the derivative in a tangent direction $V$ is $Va+\frac12d\phi(V)\,x[i,a]$ with $a=h_xyh_x^{-1}$, of norm at
least $(1-\lVert d\phi\rVert)\lvert V\rvert$; as $0\le\beta\le1$ and $0\le\beta'\le4/e$ (the maximum of $e^{1-1/s}s^{-2}$ is at $s=\frac12$),
$\lVert d\phi\rVert\le\frac14(1+4/e)<1$.  So $R_y$ is a local diffeomorphism of the compact simply connected $S^3$, hence a covering and a
diffeomorphism.  \emph{The square.}  $S(x)=xC(x)$ has complex component $u^2-\lvert v\rvert^2e^{-i\phi(x)}$, so its real part is
\[
s'=2s^2-1+\lvert v\rvert^2\bigl(1-\cos\phi(x)\bigr).
\]
Since $\lvert\phi\rvert\le\frac14$, $s'\le1-\frac32(1-s^2)$, and $s'\ge2s^2-1$; hence $S^{-1}(1)=\{\pm1\}$ and $S^{-1}(-1)=\{s=0\}$ (on $s=0$, $\phi=0$ and $S(x)=x^2=-1$);
and for $0\le s<1$, $s-s'\ge(1-s)(1+2s)-\frac12(1-s^2)\ge\frac12(1-s^2)$.  \emph{The fixed point.}  $\mathcal T$ is well defined (both branches give $\pi/2$
on $s=0$, where $S=-1$), preserves the endpoint conditions ($S(\pm1)=1$), and is a contraction with factor $\frac12$ for the sup norm, so
$\theta$ exists, is unique and continuous, and $\lVert\mathcal T^Nu_0-\theta\rVert_\infty\le\pi2^{-N}$.  It gives $\cos\theta(S(x))=2\cos^2\theta(x)-1$.  If
$\theta(x)=0$ then $s>0$ and $\theta(S(x))=0$; iterating, all $S^n(x)$ would have positive real part decreasing by at least $\frac12(1-s^2)$ at each
step unless $s=1$; so $\theta=0$ only at $1$, and then $\theta=\pi$ only at $-1$; moreover $\theta<\pi/2$ for $s>0$, $\theta>\pi/2$ for $s<0$.
\emph{Invariances.}  For $g$ in the circle $D$ of conjugations by $e^{i\alpha}$, $S(gxg^{-1})=gS(x)g^{-1}$, and $S(x^{-1})=h_x^{-1}S(x)^{-1}h_x$ (as
$h_{x^{-1}}=h_x^{-1}$); both maps preserve $s$.  So $\mathcal T$ preserves invariance under $D$ and under inversion, $u_0=\arccos s$ has both, and so
has $\theta$.  Hence $c$ is symmetric, and $(L_xa)^{-1}L_xb=h_x(a^{-1}b)h_x^{-1}$ gives $c(L_xa,L_xb)=c(a,b)$.  The values of $\theta$ give faithfulness and
the zero partners, and the fixed-point equation gives the round-trip law at $1$.  This proves (i).

(ii) At $x$, $s=\frac1{\sqrt2}$, so $s'=\frac14(1-\cos\phi(x))>0$, $\theta(S(x))<\pi/2$ and $\theta(x)<\pi/4$; at $-x$, $\phi=0$ and $S(-x)=x^2$ has real part $0$,
so $\theta(-x)=\pi-\frac12\cdot\frac\pi2$.  With $c_1=c(1,x)>\frac1{\sqrt2}$ and $c_2=c(-1,x)=\cos\theta(-x)=-\frac1{\sqrt2}$, the Gram determinant on $(1,-1,x)$ is
$-(c_1+c_2)^2<0$, and $\arccos c_1+\arccos c_2<\frac\pi4+\frac{3\pi}4=\arccos c(1,-1)$.  On the circle $\{\cos r+i\sin r\}$ the twist acts trivially, $S(y)=y^2$, and
$\arccos\re$ restricted to the circle is the fixed point of the restricted contraction, so $c(1,y)=\cos r$ there.  For $y=e^{i\theta(x)}$,
$c(1,y)=c_1$ and $c(-1,y)=-c_1\ne c_2$; a $\kappa$-preserving map fixing $1$ fixes its unique zero partner $-1$.

(iii) $T(a,a,x)=a\,B(1,\bar ax)=x$ and $T(a,x,x)=x\,B(\bar xa,1)=a$; $x\mapsto T(a,b,x)$ is left multiplication by $b$ composed with $L_{\bar ba}$ and left
multiplication by $\bar b$, and left multiplications preserve $c$ ($c(ga,gb)=c(a,b)$); $c(b,T(x,b,x))=c(1,S(\bar bx))$ gives $(\Omega3)$ at $b$.  At
$p$, $\phi(p)=\frac18e^{-1}$, and $\re S(p)=-\frac12\cos\phi(p)<0$, so $h_{S(p)}=1$.  Write $C_p(y)=h_pyh_p^{-1}$.  Then
$T(T(p,1,p),1,p)=B(S(p),p)=S(p)\,p=p\,C_p(p)\,p$, while $T(p,1,T(p,1,p))=B(p,S(p))=p\,h_pS(p)h_p^{-1}=p\,C_p(p)\,C_p^2(p)$; they differ by
$\lvert p-C_p^2(p)\rvert=2\lvert v\rvert\sin\phi(p)>0$.
(iv) $c(qae^{i\alpha},qbe^{i\alpha})=\cos\theta(e^{-i\alpha}a^{-1}be^{i\alpha})=c(a,b)$ by $D$-invariance, and the left transports
$L_x(y)=(xh_x)\,y\,e^{-i\phi(x)/2}$ lie in this group.  If $\kappa(1,\cdot)$ had a second-order expansion at $1$, its
Hessian form $Q$ would be invariant under $D$ and inversion, hence $Q(ai+bj+ck)=Aa^2+C(b^2+c^2)$; along the circles $e^{ri}$ and $e^{rj}$ the
twist acts trivially and $\theta(e^{ri})=\theta(e^{rj})=r$, so $A=C$ and $\theta(e^{\epsilon u})/\epsilon\to1$ for all unit imaginary $u$.  Remark~\ref{rem:onesidedreg}
shows that this fails.
\end{proof}

The numerical values, computed with $N=48$ steps of the contraction (error at most $\pi2^{-N}$ in $\theta$), are $\pi/4-\theta(x)=4.2584\cdot10^{-4}$, smallest
eigenvalue $-9.067\cdot10^{-8}$ of the Gram matrix on $(1,-1,x)$, and $-1.1125\cdot10^{-4}$ on $(1,x,S(x),-1)$.


\begin{proposition}[what one-sided invariance gives]\label{prop:onesidedpos}
Assume the hypotheses of Theorem~\ref{thm:round} except that the right translations $R_a$ are only continuous (not necessarily
$\kappa$-preserving).
\begin{enumerate}[label=(\roman*)]
\item Lemma~\ref{lem:roundG} holds, with the statements on $R_a$ deleted, and $G$ is a Lie group acting transitively, with compact
stabiliser $H$, on $M\cong G/H$; $B$ is jointly continuous; $M$ carries a complete $G$-invariant Riemannian metric.  The data are
equivalent to a continuous section $\ell\colon M\to G$, $\ell(x)=L_x$, together with the function $f(g)=c(w,gw)$ on $G$, which satisfies
$f(h_1gh_2)=f(g)=f(g^{-1})$ for $h_1,h_2\in H$, $f^{-1}(1)=H$, $f(\ell(x)^2)=2f(\ell(x))^2-1$, and $f=-1$ on at most one coset $gH$; conversely a
continuous section $\ell$ with $\ell(w)=e$ and a continuous $f$ with these properties define $B(x,y)=\ell(x)y$ and $c(gw,kw)=f(g^{-1}k)$ satisfying the
hypotheses.
\item The following are equivalent: (a) $(M,\kappa)$ is a round sphere with the standard comparison (as in Theorem~\ref{thm:round});
(b) every Gram matrix of $c$ is positive semidefinite; (c) the Gram matrix of $c$ on $(w,x,S(x),y)$, $S(x)=B(x,x)$, is positive
semidefinite for all $x,y$.  In that case $n\in\{1,3,7\}$, by Adams' theorem \cite{Adams} applied to the $H$-space $(M,B)$.
\item On a round sphere with the standard comparison the operations are exactly $B(a,x)=O_ax$ with $a\mapsto O_a\in\OO(n+1)$ continuous,
$O_w=I$, $O_aw=a$ and $O_aa=2\ip aw\,a-w$; on each great circle through $w$ such an operation is the circle product, so it is
power-associative.
\end{enumerate}
\end{proposition}
\begin{proof}
(i) In the proof of Lemma~\ref{lem:roundG} replace $R_x$ by $L_x$ in (i) and in the transitivity ($L_x(w)=x$); the continuity of $a\mapsto L_a$
uses only continuity of the $R_x$.  The Lie group, the smooth structure and the invariant metric come from the first part of the proof of
Lemma~\ref{lem:roundLie}; a homogeneous Riemannian metric is complete.  For the reduction: $f(h_1gh_2)=c(h_1^{-1}w,gw)=f(g)$, $f(g^{-1})=c(gw,w)=f(g)$ by
symmetry, $f(g)=1$ iff $gw=w$, and $f(L_x^2)=c(w,S(x))$; the converse is a direct check, the double invariance making $c$ well defined.

(ii) (a)$\Rightarrow$(b)$\Rightarrow$(c) are clear.  Assume (c).  With $r=c(w,x)$, $c(x,S(x))=c(L_xw,L_xx)=r$ and $c(w,S(x))=2r^2-1$, so $(1,-2r,1,0)$ has
quadratic form $0$, hence lies in the kernel, and the last row gives $c(w,y)+c(S(x),y)=2c(w,x)c(x,y)$.  Transporting by $L_p$ gives, with
$r_x(p)=L_pS(L_p^{-1}x)$,
\begin{equation}\label{eq:reflection}
c(r_x(p),y)=2c(p,x)c(x,y)-c(p,y)\qquad(p,x,y\in M).
\end{equation}
With $y=x$, $c(r_x(p),x)=c(p,x)$; applying \eqref{eq:reflection} twice gives $c(r_x^2(p),y)=c(p,y)$ for all $y$, so $r_x^2=\mathrm{id}$ by faithfulness; and
$c(r_x(p),r_x(y))=2c(p,x)c(x,y)-c(r_x(y),p)=c(p,y)$.  The map $r_x$ is continuous (part (i)), so $r_x\in G$; its fixed points are the $p$ with
$c(p,x)^2=1$, i.e.\ $x$ and its unique zero partner $x^\perp$ (indeed $c(r_x(x^\perp),x)=c(x^\perp,x)=-1$, so $r_x(x^\perp)$ is a zero partner of $x$
and equals $x^\perp$).  For an invariant metric $g$, $r_x$ is an isometric involution with isolated fixed point $x$,
so $d(r_x)_x=-I$ and $r_x$ reverses every geodesic through $x$.  Along a geodesic $\gamma$ from $w$, $r_{\gamma(t)}w=\gamma(2t)$, so \eqref{eq:reflection}
with $p=y=w$ gives the duplication law for $p(t)=c(w,\gamma(t))$; $r_w$ gives evenness, and $r_{\gamma(s/2)}r_w$ shifts the parameter by $s$, so the
return times form a closed discrete subgroup.  The proof of Lemma~\ref{lem:roundgeo}(iii),(iv) now applies verbatim: $c(w,\gamma_v(t))=\cos(kt)$ with
$k$ independent of $v$, $M$ is compact, and $c(x,y)=\cos(k\,d_g(x,y))$.  Fix $q\in M$ and put $u(t)=c(\gamma_v(t),q)$.  Since $r_{\gamma(t)}\gamma(t-s)=\gamma(t+s)$ and
$c(\gamma(t-s),\gamma(t))=\cos(ks)$, \eqref{eq:reflection} gives $u(t+s)+u(t-s)=2\cos(ks)u(t)$.  Integrating over $s\in[0,\epsilon]$ shows that $u$ is a
multiple of a sliding average of itself, hence smooth, and $u''=-k^2u$; so $u(t)=\cos(kt)c(w,q)+\sin(kt)c(U_v,q)$ with $U_v=\gamma_v(\pi/2k)$.
Applying this in both arguments and using $c(w,U_v)=0$ gives $c(\gamma_v(t),\gamma_u(s))=\cos kt\cos ks+\sin kt\sin ks\,c(U_v,U_u)$; at $s=t\to0$,
$1-\cos(kd_g(\gamma_v(t),\gamma_u(t)))\sim\frac12k^2t^2\lvert v-u\rvert^2$ gives $c(U_v,U_u)=\ip vu$.  So $(\cos kt,\sin kt\,v)\mapsto\gamma_v(t)$, $0\le t\le\pi/k$, is a
well-defined continuous bijection from the unit sphere of $\R\oplus T_wM$ onto $M$ that carries the inner product to $c$; this is (a).  A
continuous multiplication with two-sided unit on $S^n$ exists only for $n\in\{0,1,3,7\}$ \cite{Adams}.

(iii) Each $L_a$ preserves inner products and is the restriction of an orthogonal $O_a$ (proof of Lemma~\ref{lem:normal}); $O_w=I$ and $O_aw=a$ are the
unit laws, and with $r=\ip aw$ the vectors $w,a,O_aa$ have pairwise inner products $r$, $r$, $2r^2-1$, so $\lvert w+O_aa-2ra\rvert^2=0$.  Conversely
these conditions give the units, invariance and the round-trip law.  For $a=\cos t\,w+\sin t\,v$, $0<t<\pi$, the two conditions determine $O_a$
on $\operatorname{span}(w,v)$ as the rotation by $t$; by continuity $O_{-w}=-I$ there; so the great circle through $w$ in the direction $v$ is
closed under $B$, which is its circle product.
\end{proof}

\begin{remark}\label{rem:onesidedlit}
Kinyon \cite{Kinyon00} constructed on the unit sphere of every real Hilbert space a binary operation $x\odot y=(2P_{x^{1/2}}-I)Jy$
($x\ne-e_0$), $-e_0\odot y=-y$, with two-sided identity $e_0$, orthogonal left translations and the circle product on every plane through
$e_0$: an operation of the form of Proposition~\ref{prop:onesidedpos}(iii) on every round sphere.  His Theorem~5.1 states that each right
translation $R_y$ has a nonremovable discontinuity at $-e_0$; his argument \cite[(5.1)--(5.3)]{Kinyon00} gives this when $x,y,e_0$ are not coplanar
for some $x$, that is, in dimension at least $3$ with $y\ne\pm e_0$; for $y=\pm e_0$, $R_y=\pm\mathrm{id}$, and on $S^1$ the operation is complex
multiplication \cite[(3.6)]{Kinyon00}.  By Proposition~\ref{prop:onesidedpos}(ii) and Adams' theorem, continuity of all right translations is
impossible outside dimensions $1,3,7$.  We did not find the twisted operations $xh_xyh_x^{-1}$ of
Proposition~\ref{prop:onesided} in the literature we searched.
\end{remark}

\begin{remark}[regularity of the example on the diagonal]\label{rem:onesidedreg}
Near the unit, $S(e^X)=e^{2X+\delta(X)}$ with $\delta(X)=\Ad(h_{e^X})X-X+O(\lvert X\rvert^3)=O(\lvert X\rvert^2)$: the twist rotates the $(j,k)$-part of $X$ by
$\phi\approx\frac14t$, which is of first order in $\lvert X\rvert$ because $\beta(1)=1$.  For a unit imaginary $u$ and small $\epsilon_0>0$ put
$W_k=S^k(e^{2^{-k}\epsilon_0u})$.  Since $S(e^{2^{-k-1}\epsilon_0u})=e^{2^{-k}\epsilon_0u}$ up to $O(4^{-k})$ and $S^k$ is Lipschitz with constant $O(2^k)$ near
the unit, $(W_k)$ is a Cauchy sequence; its limit $W(\epsilon_0)$ satisfies $S(W(\epsilon_0))=W(2\epsilon_0)$, and, all points lying in the positive
hemisphere, $\theta(e^{2^{-k}\epsilon_0u})/(2^{-k}\epsilon_0)=\theta(W_k)/\epsilon_0\to\Lambda_u(\epsilon_0):=\theta(W(\epsilon_0))/\epsilon_0$ with
$\Lambda_u(2\epsilon_0)=\Lambda_u(\epsilon_0)$.  So $\theta(e^{\epsilon u})/\epsilon$ converges along every geometric sequence, to a log-periodic function of $\epsilon$.
Along $i$ and $j$ it is identically $1$.  For $u=(i+j)/\sqrt2$ the computed values are $\Lambda_u(2^{-k})=0.999071031057$ and
$\Lambda_u(1.5\cdot2^{-k})=0.999169235236$ (stable for $14\le k\le24$ in double precision with the error bound above; an independent
computation at $120$ digits agrees), so $\Lambda_u$ is neither constant nor $1$, and $\kappa$ has no second-order expansion at the diagonal.  The
non-constancy is a numerical, not an interval-certified, statement.
\end{remark}


% ---- insertion point: Part "Beyond the axiom" (modules m2, m3, d4)
\paperpart{part:arc}{The calculus of minimal arcs}

\section{Minimal arcs and the geodesic recurrence}\label{sec:recur}

\begin{definition}\label{def:arc}
For $x,y\in M$ with $y\ne-x$, the \emph{minimal arc} from $x$ to $y$ is the shortest great-circle arc,
of length $d(x,y)=2\arccos\sqrt{\kappa(x,y)}$.  The \emph{sagitta} of the step $(x,y)$ is
\[
v(x,y):=1-c(x,y)=1-\cos d(x,y)=2\bigl(1-\kappa(x,y)\bigr)=\tfrac12\lVert y-x\rVert^2 .
\]
\end{definition}

By Proposition~\ref{prop:linear}(iii), $\tfrac12\bigl(y+s_x(y)\bigr)=c(x,y)\,x=(1-v(x,y))\,x$: the chord from
$y$ to its mirror image $s_xy$ has its midpoint on the ray through $x$, at height $v(x,y)$ below $x$.  So
$v(x,y)$ is the sagitta of the arc of length $2d(x,y)$ from $y$ through $x$ to $s_xy$.  The sagitta of the
minimal arc from $x$ to $y$ itself is $1-\sqrt{\kappa(x,y)}=1-\cos(d/2)$, the distance from the midpoint
$m(x,y)$ to $\tfrac12(x+y)=\sqrt\kappa\,m$.

\begin{theorem}[classical in substance; cf.\ \cite{Loos,MH}]\label{thm:recur}
For a sequence $(x_n)_{n\ge0}$ in $M$ the following are equivalent:
\begin{enumerate}[label=(\alph*)]
\item $x_{n+1}=\hp{x_n}{x_{n-1}}{x_n}$ for all $n\ge1$;
\item $x_n=\ell^n(x_0)$ for the left transport $\ell=\ell_{x_1,x_0}$, i.e.\ $x_n=g^nx_0$ with $g=x_1\bar x_0$;
\item $x_n=r^n(x_0)$ for the right transport $r=r_{x_0,x_1}$;
\item $x_{n+1}=2c\,x_n-x_{n-1}$ for all $n\ge1$, where $c:=c(x_0,x_1)$.
\end{enumerate}
In this case:
\begin{enumerate}[label=(\roman*)]
\item if $x_1\ne\pm x_0$, then $x_n=\cos(n\varepsilon)x_0+\sin(n\varepsilon)\tau$ with
$\varepsilon=d(x_0,x_1)$ and $\tau=(x_1-\cos\varepsilon\,x_0)/\sin\varepsilon$: the points are the samples at
integer times of the constant-speed geodesic through $x_0,x_1$;
\item for $n\ge1$, $x_n=U_{n-1}(c)\,x_1-U_{n-2}(c)\,x_0$ (and for $n=0$ with the convention $U_{-2}=-1$), where $U_k$ are the
Chebyshev polynomials of the second kind
($U_{-1}=0$, $U_0=1$); hence $c(y,x_{n+1})=2c\,c(y,x_n)-c(y,x_{n-1})$ for every $y\in M$, and
$c(x_0,x_n)=T_n(c)$;
\item $x_{n+1}-2x_n+x_{n-1}=-2v\,x_n$ and $\tfrac12(x_{n-1}+x_{n+1})=(1-v)\,x_n$, with $v=v(x_0,x_1)$;
\item $x_{n+1}\bar x_n$, $\bar x_nx_{n+1}$, the bivector $x_n\wedge x_{n+1}\in\Lambda^2\R^4$ and
$c(x_n,x_{n+1})$ do not depend on $n$.
\end{enumerate}
\end{theorem}
\begin{proof}
(a)$\Leftrightarrow$(b): (a) says $x_{n+1}\bar x_n=x_n\bar x_{n-1}$.  (b)$\Leftrightarrow$(c):
$g^nx_0=x_0(\bar x_0gx_0)^n=x_0(\bar x_0x_1)^n$.  (a)$\Leftrightarrow$(d): by
Proposition~\ref{prop:linear}(iii), (a) reads $x_{n+1}=2c(x_n,x_{n-1})x_n-x_{n-1}$.  Under (a),
$c(x_{n+1},x_n)=c(s_{x_n}x_{n-1},s_{x_n}x_n)=c(x_{n-1},x_n)$, so the coefficient is constant.  Under (d),
$c(x_{n+1},x_n)=2c-c(x_{n-1},x_n)$, and induction from $c(x_1,x_0)=c$ shows the coefficient is $c(x_n,x_{n-1})$,
so (d) implies (a).  (i), (ii): $\cos(n\varepsilon)$, $\sin(n\varepsilon)$, $U_{n-1}(c)$ satisfy the linear
recurrence (d) with the right initial values, and $T_n=cU_{n-1}-U_{n-2}$.  (iii) is (d) with $v=1-c$.  (iv)
follows from (b), (c) and $x_n\wedge(2cx_n-x_{n-1})=x_{n-1}\wedge x_n$.
\end{proof}

The recurrence uses three multiplications and no transcendental function; given two consecutive samples,
every further sample is a polynomial with integer coefficients in their coordinates.  Statement (iii) is the
exact form of the geodesic equation $\ddot x=-\lvert\dot x\rvert^2x$: the discrete acceleration is normal,
of size twice the sagitta.  The quantities in (iv) are conserved because the discrete action of Theorem~\ref{thm:action} is invariant under left and right
transports (a discrete Noether theorem in the sense of \cite{MW}).

\begin{remark}[floating point]
Off the sphere the cubic form $x_n\bar x_{n-1}x_n$ multiplies norms: if $\lvert x_k\rvert=e^{\delta_k}$ then
$\delta_{n+1}=2\delta_n+\delta_{n-1}$, so rounding errors in the norm grow like $(1+\sqrt2)^n$.  The linear
form (d) with $c$ fixed has a companion matrix with eigenvalues $e^{\pm i\varepsilon}$, and its errors grow at
most linearly.  Over $\Q$ both forms are exact.  Both behaviours are confirmed numerically in
\texttt{check\_arc\_calculus.py}.
\end{remark}

\begin{theorem}[discrete least action]\label{thm:action}
Let $p,q\in M$ with $\Theta:=d(p,q)<\pi$, and $N\ge1$.  For $x_0=p$, $x_N=q$ and arbitrary
$x_1,\dots,x_{N-1}\in M$,
\[
S(x):=\sum_{k=0}^{N-1}v(x_k,x_{k+1})\ \ge\ N\bigl(1-\cos(\Theta/N)\bigr),
\]
with equality if and only if $x_k=\gamma(k\Theta/N)$ for the unit-speed minimal geodesic $\gamma$ from $p$ to
$q$; this minimiser is the solution of the recurrence of Theorem~\ref{thm:recur} with $x_0=p$,
$x_1=\gamma(\Theta/N)$.  Moreover, a sequence with fixed ends whose steps all satisfy $d(x_k,x_{k+1})<\pi/2$ is
a critical point of $S$ on $M^{N-1}$ if and only if it satisfies the recurrence.
\end{theorem}
\begin{proof}
Put $f(t)=1-\cos t$ and $d_k=d(x_k,x_{k+1})\in[0,\pi]$.  If $\Theta=0$ the bound is $0$ and equality forces all
$d_k=0$, i.e.\ the constant sequence; assume $\Theta>0$.  By the triangle inequality $\sum d_k\ge\Theta>0$.
With $d_k':=d_k\Theta/\sum_jd_j\le d_k$ and $f$ increasing on $[0,\pi]$, $\sum f(d_k)\ge\sum f(d_k')$.  On the
compact simplex $\{t_k\ge0,\ \sum t_k=\Theta\}$ the function $\sum f(t_k)$ attains its minimum.  For
$s,t\ge0$ with $s+t\le\Theta<\pi$,
\[
f(s)+f(t)=2-2\cos\tfrac{s+t}2\cos\tfrac{s-t}2\ \ge\ 2-2\cos\tfrac{s+t}2=2f\bigl(\tfrac{s+t}2\bigr),
\]
with equality iff $s=t$, because $\cos\frac{s+t}2>0$.  Replacing two unequal entries by their mean therefore
strictly decreases the sum, so the minimum is attained only at $t_k=\Theta/N$ and equals $Nf(\Theta/N)$.
Equality in the theorem forces $d_k=d_k'$ for all $k$ (as $f$ is strictly increasing), i.e.\
$\sum d_k=\Theta$, and $d_k=\Theta/N$.  Then
$d(x_0,x_N)=d(x_0,x_k)+d(x_k,x_N)$ with $d(x_0,x_k)=k\Theta/N$ for every $k$ (both triangle inequalities
$d(x_0,x_k)\le\sum_{i<k}d_i$, $d(x_k,x_N)\le\sum_{i\ge k}d_i$ must be equalities).  A broken geodesic of length $d(p,q)$ is a
minimising geodesic, and since $\Theta<\pi$ the minimising geodesic from $p$ to $q$ is unique; so it contains each $x_k$
at distance $k\Theta/N$ from $p$.  The
converse and the identification with the recurrence follow from Theorem~\ref{thm:recur}(i).

Critical points: the ambient gradient of $S$ in $x_k$ ($0<k<N$) is $-(x_{k-1}+x_{k+1})$, so $x$ is critical iff
$x_{k-1}+x_{k+1}=\mu_kx_k$ for some $\mu_k\in\R$.  Taking norms in $x_{k+1}=\mu_kx_k-x_{k-1}$ gives
$\mu_k^2=2\mu_kc(x_k,x_{k-1})$, so $\mu_k=0$ or $\mu_k=2c(x_k,x_{k-1})$.  If $\mu_k=0$ then $x_{k+1}=-x_{k-1}$
and $d(x_k,x_{k+1})=\pi-d(x_k,x_{k-1})$, impossible when both are $<\pi/2$.  So $\mu_k=2c(x_k,x_{k-1})$, which is
the recurrence by Proposition~\ref{prop:linear}(iii); the converse is the same computation.
\end{proof}

The Euler--Lagrange half is an instance of the discrete variational principle of
\cite{Veselov88,MV,MW}: up to an additive constant, $S$ is $-\sum\ip{x_k}{Jx_{k+1}}$ on the unit sphere with
$J=1$, the Lagrangian of the discrete Neumann system of \cite{MV} in its degenerate case.  The global inequality is elementary; it shows that the heap
recurrence is the exact discrete minimiser, not merely a stationary point.  The proof shows more: a critical sequence satisfies the recurrence at every interior $k$ with
$d(x_{k-1},x_k)+d(x_k,x_{k+1})\ne\pi$, because the alternative $\mu_k=0$ forces $x_{k+1}=-x_{k-1}$ and hence
$d(x_k,x_{k+1})=\pi-d(x_{k-1},x_k)$; so it suffices that no two consecutive steps sum to $\pi$.  The constant $\pi/2$ is sharp: for
$0<\varepsilon<\pi/2$ and $d(x_0,x_1)=\pi/2-\varepsilon$, the sequence $(x_0,x_1,-x_0,-x_1)$ has steps
$\pi/2-\varepsilon,\pi/2+\varepsilon,\pi/2-\varepsilon$, is critical (both $x_0+(-x_0)$ and $x_1+(-x_1)$ vanish), and violates the recurrence,
since $s_{x_1}(x_0)=2\sin\varepsilon\,x_1-x_0\ne-x_0$ (defect $2\sin\varepsilon$).

\begin{theorem}[consistency]\label{thm:consist}
Let $x\colon(a,b)\to M$ be of class $C^4$ and $\nabla_t\dot x=\ddot x+\lvert\dot x\rvert^2x$ its covariant
acceleration.  Uniformly for $t$ in compact subintervals, as $h\to0$,
\begin{align*}
x(t+h)-\hp{x(t)}{x(t-h)}{x(t)}&=h^2\,\nabla_t\dot x(t)-\tfrac{h^3}2\,\tfrac{d}{dt}\lvert\dot x\rvert^2(t)\;x(t)+O(h^4),\\
\Delta_h^2x(t):=x(t+h)+x(t-h)-(c_++c_-)\,x(t)&=h^2\,\nabla_t\dot x(t)+O(h^4),
\end{align*}
where $c_\pm=c(x(t),x(t\pm h))$, and $\Delta_h^2x(t)$ is tangent to $M$ at $x(t)$ for every $h$.  The
residual of the recurrence vanishes for all $t$ and all small $h$ iff $x$ is a constant-speed geodesic.
\end{theorem}
\begin{proof}
Write $x=x(t)$.  Differentiating $\lvert x\rvert^2=1$ gives $\ip x{\dot x}=0$, $\ip x{\ddot x}=-\lvert\dot x\rvert^2$,
$\ip x{x'''}=-3\ip{\dot x}{\ddot x}=-\tfrac32(\lvert\dot x\rvert^2)'$.  Taylor's formula with remainder gives
$x(t+h)+x(t-h)=2x+h^2\ddot x+O(h^4)$ and
$c_\pm=1-\tfrac{h^2}2\lvert\dot x\rvert^2\mp\tfrac{h^3}4(\lvert\dot x\rvert^2)'+O(h^4)$.  By
Proposition~\ref{prop:linear}(iii) the residual is $x(t+h)+x(t-h)-2c_-x$, which gives the first formula.  For
the second, $c_++c_-=2-h^2\lvert\dot x\rvert^2+O(h^4)$, and $\ip x{\Delta_h^2x}=c_++c_--(c_++c_-)=0$.  If the
residual vanishes identically, the $h^2$ term gives $\nabla_t\dot x=0$; conversely,
Theorem~\ref{thm:recur}(i) applies to equally spaced samples of a constant-speed geodesic.
\end{proof}

So the heap gives a second difference that is intrinsic, exactly tangent in its symmetric form, and of
the same order as the linear second difference $x(t+h)-2x(t)+x(t-h)=h^2\ddot x+O(h^4)$, whose normal part
$-h^2\lvert\dot x\rvert^2x$ is the sagitta term $-2v\,x$ of Theorem~\ref{thm:recur}(iii).

\section{Sagittas and quadratic variation}

\begin{proposition}[exact split]\label{prop:split}
For $x,y\in M$,
\[
y-x=\im(y\bar x)\,x-v(x,y)\,x,\qquad \im(y\bar x)\,x=P_x(y-x),\qquad \lvert\im(y\bar x)\rvert=\sin d(x,y),
\]
where $P_x$ is the orthogonal projection onto $T_xM$.  Consequently, for points $x_0,\dots,x_n$,
\[
x_n-x_0=\sum_{k<n}P_{x_k}(x_{k+1}-x_k)-\sum_{k<n}v(x_k,x_{k+1})\,x_k .
\]
\end{proposition}
\begin{proof}
$\im(y\bar x)x=(y\bar x-c(x,y))x=y-c(x,y)x$, which is orthogonal to $x$ and differs from $y-x$ by
$(1-c)x=vx$; and $\ip x{y-x}=-v$.  Finally $\lvert y\bar x\rvert=1$ and $\re(y\bar x)=\cos d$.
\end{proof}

Each step is thus an exact sum of a rotation (the tangential part, the left-point Riemann increment of
$\int P_x\,dx$) and a radial sagitta.  The next theorem says when the sagittas survive in the limit.  Let
$x\colon[0,T]\to M$ be continuous, $\pi=\{0=t_0<\dots<t_n=T\}$ a partition with mesh $\lvert\pi\rvert$,
$x_k=x_{t_k}$, $d_k=d(x_k,x_{k+1})$, $v_k=v(x_k,x_{k+1})$, and for $f\in C(M,\R^m)$ put
$V_\pi(x;f):=\sum_kv_k\,f(x_k)$, $V_\pi(x):=V_\pi(x;1)$.  Following F\"ollmer \cite{Follmer}, $x$
\emph{has quadratic variation along} a sequence of partitions $(\pi_n)$ with $\lvert\pi_n\rvert\to0$ if the
measures $\mu_n:=\sum_k\lVert x_{k+1}-x_k\rVert^2\delta_{t_k}$ on $[0,T]$ converge weakly to a measure $\mu$
without atoms; then $[x]_t:=\mu([0,t])$.

\begin{theorem}[sagitta sums]\label{thm:sagitta}
\begin{enumerate}[label=(\alph*)]
\item For every partition, $V_\pi(x)\le\tfrac12\max_kd_k\sum_kd_k$.  Hence $V_\pi(x)\to0$ as
$\lvert\pi\rvert\to0$ if $x$ is rectifiable.  If $\lVert x_t-x_s\rVert\le H\lvert t-s\rvert^\alpha$ with
$\alpha>\tfrac12$, then $V_\pi(x)\le\tfrac12H^2T\lvert\pi\rvert^{2\alpha-1}$.
\item If $x$ has quadratic variation along $(\pi_n)$, then for every $f\in C(M,\R^m)$
\[
V_{\pi_n}(x;f)\to\tfrac12\int_0^Tf(x_t)\,d[x]_t,\qquad \sum_kd_k^2f(x_k)\to\int_0^Tf(x_t)\,d[x]_t,
\]
and the left Riemann sums $\sum_kP_{x_k}(x_{k+1}-x_k)$ converge to $x_T-x_0+\tfrac12\int_0^Tx_t\,d[x]_t$.
\item Let $x$ be Brownian motion on $S^3$ (generator $\tfrac12\Delta$, any initial point).  For every sequence of
deterministic partitions with $\lvert\pi_n\rvert\to0$ and every $f\in C(M,\R^m)$, in probability,
\[
V_{\pi_n}(x;f)\to\tfrac32\int_0^Tf(x_t)\,dt,\qquad\text{in particular}\quad V_{\pi_n}(x)\to\tfrac32T,\quad
\sum_kv_kx_k\to\tfrac32\int_0^Tx_t\,dt,
\]
and the tangential sums converge to the It\^o integral $\int_0^TP_{x_t}\,dx_t$.
\item If the partitions in (c) are nested, their union is dense, and $\sum_n\lvert\pi_n\rvert<\infty$ (for instance the
dyadic partitions), the convergence in (c) holds almost surely, simultaneously for all $f\in C(M,\R^m)$.
\end{enumerate}
\end{theorem}
\begin{proof}
(a) $v_k=1-\cos d_k\le d_k^2/2\le\tfrac12\max_jd_j\,d_k$.  If $x$ is rectifiable, $\sum d_k\le L(x)$ and
$\max d_k\to0$ by uniform continuity.  In the H\"older case $v_k=\tfrac12\lVert x_{k+1}-x_k\rVert^2\le
\tfrac12H^2(t_{k+1}-t_k)^{2\alpha}\le\tfrac12H^2\lvert\pi\rvert^{2\alpha-1}(t_{k+1}-t_k)$.

(b) $V_{\pi_n}(x;f)=\tfrac12\int f(x_t)\,d\mu_n(t)$ because $v_k=\tfrac12\lVert x_{k+1}-x_k\rVert^2$, and
$t\mapsto f(x_t)$ is bounded and continuous.  Since $2(1-\cos d)=d^2(1+\epsilon(d))$ with $\epsilon(d)\to0$ as
$d\to0$ and $\max_kd_k\to0$, the second sum has the same limit as $\int f(x_t)\,d\mu_n$.  The last claim is
Proposition~\ref{prop:split} combined with the first limit for $f(x)=x$.

(c) Brownian motion on $S^3$ is the solution of Stroock's equation $dx=P_x\,dB-\tfrac32x\,dt$ with $B$ a
Brownian motion in $\R^4$ \cite{Stroock,Hsu}; hence $d[x^i,x^j]_t=(\delta_{ij}-x^i_tx^j_t)\,dt$
and $\operatorname{tr}[x]_t=3t$.  For a continuous semimartingale the sums over the increments with $t_{k+1}\le t$ converge in probability to
$\operatorname{tr}[x]_t$ for each $t$ \cite[Ch.~IV, \S1]{RY}; $\mu_n([0,t])$ differs from them by at most one increment, which tends
to $0$ uniformly by continuity of the path, so $\mu_n([0,t])\to3t$ in probability.
Given a subsequence, extract (diagonally) a further subsequence along which this holds almost surely
simultaneously for all $t$ in the countable set $(\Q\cap[0,T])\cup\{T\}$; on that event the distribution functions
$t\mapsto\mu_n([0,t])$ are nondecreasing and converge on a dense set containing $T$ to the continuous function $3t$,
so $\mu_n\to3\,dt$ weakly, and (b) applies pathwise.  All sums involved are finite sums of random variables, hence
measurable.  As every subsequence has an almost surely convergent subsequence, the convergence holds in
probability.  Finally $\int P_x\,dx=\int P_x\,dB$ because $P_x^2=P_x$ and $P_xx=0$, and this equals
$x_T-x_0+\tfrac32\int x\,dt$, the limit in (b).

(d) Let $\pi$ be a partition of $[0,t]$ with steps $h_k$ and mesh $\lvert\pi\rvert\le1$.  Write
$\Delta x_k=M_k+A_k$ with $M_k=\int_{t_k}^{t_{k+1}}P_{x_s}\,dB_s$ and $A_k=-\tfrac32\int_{t_k}^{t_{k+1}}x_s\,ds$, so
$\lvert A_k\rvert\le\tfrac32h_k$.  For $N_s=\int_{t_k}^sP_x\,dB$, It\^o's formula and $\operatorname{tr}P_x=3$ give
$\lvert M_k\rvert^2-3h_k=\zeta_k:=2\int_{t_k}^{t_{k+1}}\ip{N_s}{P_{x_s}dB_s}$.  The $\zeta_k$ are orthogonal martingale increments
with $\mathbb E\zeta_k^2\le4\int_{t_k}^{t_{k+1}}\mathbb E\lvert N_s\rvert^2ds=6h_k^2$.  Moreover
$(\sum_kh_k\lvert M_k\rvert)^2\le t\sum_kh_k\lvert M_k\rvert^2$, whose expectation is at most $3t^2\lvert\pi\rvert$, and
$\sum_k\lvert A_k\rvert^2\le\tfrac94t\lvert\pi\rvert$.  Since
$\sum_k\lvert\Delta x_k\rvert^2-3t=\sum\zeta_k+\sum2\ip{M_k}{A_k}+\sum\lvert A_k\rvert^2$,
\[
\mathbb E\Bigl[\Bigl(\sum_k\lvert\Delta x_k\rvert^2-3t\Bigr)^2\Bigr]\le3\bigl(6t+27t^2+\tfrac{81}{16}t^2\bigr)\lvert\pi\rvert .
\]
For $t$ in the countable dense set $\bigcup_n\pi_n$, $t\in\pi_n$ for large $n$, and applying the bound to $\pi_n\cap[0,t]$ gives
$\mathbb E\sum_n(S_n(t)-3t)^2<\infty$ for the sums $S_n(t)$ over the increments in $[0,t]$, so $S_n(t)\to3t$ almost surely; as in
(c), $\mu_n([0,t])$ differs from $S_n(t)$ by one increment, which tends to $0$, so $\mu_n([0,t])\to3t$ almost surely.  On the intersection of these countably many
events the argument of (c) gives $\mu_n\to3\,dt$ weakly, and (b) applies for every $f$.  This is the argument of L\'evy's
theorem for Brownian motion in $\R^d$.
\end{proof}

\begin{remark}
(1) The intrinsic and ambient quadratic variations coincide: $\sum d_k^2$ and $2V_\pi$ have the same limit.
(2) In (c) the sagitta sum is the It\^o drift: in $x_T-x_0=\sum P_{x_k}\Delta x_k-\sum v_kx_k$ the first sum
tends to the martingale part and the second to $\tfrac32\int x\,dt$.  For smooth paths the second sum
vanishes and the first tends to $\int\dot x\,dt$; the sagitta is the term of second order in the step
that the linear difference $\Delta x$ hides in its normal component.
(3) The \emph{geodesic random walk} $x_{k+1}=\exp(\xi_k)x_k$ with $\xi_k$ independent $N(0,hI_3)$ consists of
minimal arcs; $\mathbb E\,v_k=\mathbb E(1-\cos\lvert\xi_k\rvert)=\tfrac32h+O(h^2)$, so $\sum_{k<T/h}v_k\to\tfrac32T$
in $L^2$.  That the walk converges to Brownian motion on $S^3$ is a classical theorem
\cite{Jorgensen}; piecewise-geodesic interpolations of the sample points produce Stratonovich limits by
Wong--Zakai type theorems \cite{WZ}, \cite[Ch.~VI]{IW}.
(4) For deterministic equal steps of angle $\sigma\lvert\pi\rvert^\alpha$ one obtains $V_\pi\to0$,
$\tfrac12\sigma^2T$, $\infty$ for $\alpha>\tfrac12$, $=\tfrac12$, $<\tfrac12$, by direct summation.  Theorem~\ref{thm:sagitta} replaces it by statements about arbitrary partitions and paths.
\end{remark}

\begin{proposition}[fluctuations; Jacod]\label{prop:clt}
Let $x$ be Brownian motion on $S^3$ and $\pi_n$ the partition of $[0,T]$ into $N_n$ equal steps, $N_n\to\infty$.  Then
$\sqrt{N_n}\,\bigl(V_{\pi_n}(x)-\tfrac32T\bigr)\to N\bigl(0,\tfrac32T^2\bigr)$ stably in law (in particular in law); for the
dyadic partitions, $2^{n/2}\bigl(V_{\pi_n}(x)-\tfrac32T\bigr)\to N(0,\tfrac32T^2)$.  Moreover, with $h=T/N_n$,
\[
\mathbb E\,V_{\pi_n}=N_n\bigl(1-e^{-3h/2}\bigr),\qquad
\operatorname{Var}\bigl(\sqrt{N_n}\,V_{\pi_n}\bigr)=N_n^2\Bigl[\tfrac{1+3e^{-4h}}4-e^{-3h}\Bigr]\to\tfrac32T^2,
\]
and $\sqrt{N_n}(\mathbb E\,V_{\pi_n}-\tfrac32T)=-\tfrac{9T^2}{8\sqrt{N_n}}+O(N_n^{-3/2})$.
\end{proposition}
\begin{proof}[Proof by Jacod's theorem]
By Stroock's equation, $x$ is a continuous It\^o semimartingale with drift $-\tfrac32x$ and diffusion matrix $c=P_xP_x^{\top}=P_x$,
both bounded.  With $\Delta_n=T/N_n$, the theorem of Jacod \cite[Thm.~5.4.2]{JP12} gives stable convergence in law of
$\Delta_n^{-1/2}\bigl(\sum_k\Delta x_k\Delta x_k^{\top}-[x,x]_T\bigr)$ to a centred variable $Z$, conditionally Gaussian given the
path, with $\mathbb E(Z^{ij}Z^{kl}\mid\mathcal F)=\int_0^T(c^{ik}c^{jl}+c^{il}c^{jk})\,ds$.  For the trace,
\[
\operatorname{Var}(\operatorname{tr}Z\mid\mathcal F)=2\int_0^T\operatorname{tr}(c^2)\,ds=2\int_0^T\operatorname{tr}P_{x_s}\,ds=6T,
\]
which is deterministic, so
$\Delta_n^{-1/2}(\sum_k\lvert\Delta x_k\rvert^2-3T)\to N(0,6T)$ in law.  Since $V_{\pi_n}=\tfrac12\sum_k\lvert\Delta x_k\rvert^2$ and
$\sqrt{N_n}=\sqrt T\,\Delta_n^{-1/2}$, the limit of $\sqrt{N_n}(V_{\pi_n}-\tfrac32T)$ is $\tfrac12\sqrt T\,N(0,6T)=N(0,\tfrac32T^2)$.
\end{proof}
\begin{proof}[Elementary proof]
Let $d_k=d(x_{t_k},x_{t_{k+1}})$, so $V_{\pi_n}=\sum_k(1-\cos d_k)$.  By the Markov property and the invariance of the law of
Brownian motion under the isometries of $S^3$, which act transitively, the conditional law of $d_k$ given $\mathcal F_{t_k}$ is the
law of $d(x_0,x_h)$ for Brownian motion started anywhere; so the $d_k$ are independent and identically distributed.  The
function $y\mapsto\ip{x_0}y$ is a spherical harmonic of degree $1$, with $\tfrac12\Delta=-\tfrac32$ on it, and
$\ip{x_0}y^2-\tfrac14$ one of degree $2$, with $\tfrac12\Delta=-4$; hence $\mathbb E\cos d_k=e^{-3h/2}$ and
$\mathbb E\cos^2d_k=\tfrac14+\tfrac34e^{-4h}$, which give the stated mean and variance and their expansions.  Since $1-\cos d_k\le d_k^2/2\le\tfrac{\pi^2}8\lvert\Delta x_k\rvert^2$ and
$\mathbb E\lvert\Delta x_k\rvert^8\le Ch^4$ (by Burkholder's inequality for the martingale part, and $\lvert A_k\rvert\le\tfrac32h$), the summands $\sqrt{N_n}(1-\cos d_k-\mathbb E(1-\cos d_k))$ satisfy Lyapunov's
condition, $N_n\cdot N_n^2\cdot O(h^4)=O(N_n^{-1})\to0$, and the central limit theorem for triangular arrays applies.
\end{proof}

\section{Product integrals and the fundamental theorem}

For a $C^1$ curve $x$ in $M$, the \emph{left} and \emph{right logarithmic derivatives} are
$\lambda_x:=\dot x\bar x$ and $\rho_x:=\bar x\dot x$, curves in $\im\HH\cong\su(2)$.  In heap terms,
$\ell_{x(t+h),x(t)}=\exp(h\lambda_x(t)+o(h))\cdot$ and $r_{x(t),x(t+h)}=\cdot\,\exp(h\rho_x(t)+o(h))$: the
logarithmic derivatives are the generators of the transports between nearby points, the replacement of the
difference quotient.  The discrete statement is exact and purely heap-theoretic: by para-associativity,
$\hp ab{\hp bcz}=\hp{\hp abb}cz=\hp acz$, so
\begin{equation}\label{eq:telescope}
\ell_{x_n,x_{n-1}}\circ\dots\circ\ell_{x_1,x_0}=\ell_{x_n,x_0}\qquad\text{for all }x_0,\dots,x_n\in M .
\end{equation}

\begin{theorem}[classical \cite{Volterra,DF,Slavik,CG,IMNZ,Iserles01,DLP,Sharpe}]\label{thm:ftc}
\begin{enumerate}[label=(\roman*)]
\item For $A\in C([0,T],\im\HH)$ and $x_0\in M$ the problem $\dot x=Ax$, $x(0)=x_0$ has a unique solution, and
it stays in $M$.  For partitions with tags $\tau_k\in[t_k,t_{k+1}]$ the ordered products
\[
E\bigl(\Delta t_{n-1}A(\tau_{n-1})\bigr)\cdots E\bigl(\Delta t_0A(\tau_0)\bigr)\,x_0
\]
converge to $x(T)$, both for $E=\exp$ and for $E(\xi)=\Cay(\xi/2)$ (Section~\ref{sec:rational}).  For uniform steps $h$ the error is $O(h)$ if $A\in C^1$,
and $O(h^2)$ with midpoint tags if $A\in C^2$.
\item For $x\in C^1$, $\lambda_x$ is the unique $A$ in (i) reproducing $x$; $\lambda_{xb}=\lambda_x$ and
$\lambda_{ax}=a\lambda_x\bar a$ for $a,b\in M$.
\item \emph{(Maurer--Cartan.)} For $U\subset\R^2$ open, connected and simply connected, and $A_1,A_2\in C^1(U,\im\HH)$ there is
$F\colon U\to M$ with $\partial_iF=A_iF$ iff $\partial_1A_2-\partial_2A_1=[A_1,A_2]$; then $F$ is unique up
to $F\mapsto Fb$, i.e.\ up to a right transport.
\end{enumerate}
\end{theorem}
\begin{proof}[Comments]
(i) is the existence theory of linear equations together with $\ip x{Ax}=0$; convergence of the product is
Volterra's product integral \cite{Volterra,DF,Slavik}; the orders are those of the Lie--Euler and
exponential/Cayley midpoint methods \cite{CG,IMNZ,Iserles01,DLP}.  (ii) is direct.  (iii) is the
fundamental theorem of calculus for Lie-group-valued maps \cite[Ch.~3]{Sharpe}; the sign follows
from $\partial_1\partial_2F=(\partial_1A_2+A_2A_1)F=(\partial_2A_1+A_1A_2)F$.
\end{proof}

The discrete counterparts of (ii) and (iii), exact at every mesh size, are \eqref{eq:telescope} and
Corollary~\ref{cor:poincare} below.

\section{Finite Stokes theorem and holonomy of minimal arcs}\label{sec:stokes}

\subsection{Cochains and holonomy}
Let $K$ be a finite planar polygonal complex whose underlying space is a closed disk, with the
counterclockwise orientation.  An \emph{$\SU(2)$-cochain} $U$ assigns to every oriented edge $p\to q$ an
element $U_{qp}\in M$ with $U_{pq}=U_{qp}^{-1}$.  For an edge path $\gamma=(p_0\to p_1\to\dots\to p_n)$ put
$\hol(\gamma):=U_{p_np_{n-1}}\cdots U_{p_1p_0}$, so $\hol(\gamma*\gamma')=\hol(\gamma')\hol(\gamma)$ for
concatenation.  For a closed path we speak of its holonomy at the base point.  If $\gamma=\tau*\rho$ is a closed
path at $p$ and $\tau$ ends at $p'$, then $\rho*\tau$ is the same loop based at $p'$ and
\begin{equation}\label{eq:rebase}
\hol(\rho*\tau)=\hol(\tau)\,\hol(\tau*\rho)\,\hol(\tau)^{-1}.
\end{equation}

\begin{lemma}[cutting]\label{lem:cut}
Let $\gamma=\alpha*\beta$ be a closed edge path at $p$, with $\alpha$ from $p$ to $q$, and let $\delta$ be an edge
path from $q$ to $p$.  Then $\hol(\gamma)=\hol(\bar\delta*\beta)\,\hol(\alpha*\delta)$, both loops based at $p$.
\end{lemma}
\begin{proof}
$\hol(\alpha*\delta*\bar\delta*\beta)=\hol(\beta)\hol(\delta)^{-1}\hol(\delta)\hol(\alpha)=\hol(\gamma)$.
\end{proof}

If $\delta$ is a chain of interior edges cutting $K$ into two disks, then $\alpha*\delta$ and $\bar\delta*\beta$ are
their counterclockwise boundaries: holonomy is multiplicative under cutting, with the base point at an end
of the cut.  Iterating this along a shelling gives the finite Stokes theorem.

\begin{definition}
Let $K$ be a finite planar complex whose underlying space is a closed disk, each closed face being a closed disk
bounded by an edge cycle; all boundary cycles are oriented counterclockwise.  A \emph{shelling} of $K$ is an
enumeration $f_1,\dots,f_N$ of its faces such that for $2\le j\le N$ the union $D_{j-1}:=f_1\cup\dots\cup f_{j-1}$ is a
closed disk and $f_j\cap D_{j-1}$ is an arc $\alpha_j$ formed by at least one edge (a simple edge path with distinct
endpoints).
\end{definition}

\begin{lemma}[gluing two disks]\label{lem:glue}
Let $D,f\subset\R^2$ be closed disks such that $D\cap f=\alpha$ is an arc contained in $\partial D\cap\partial f$ and different
from $\partial D$ and from $\partial f$; let $a,b$ be its endpoints, $\beta$ the closure of $\partial D\setminus\alpha$ and $\rho$
that of $\partial f\setminus\alpha$.  Then $D\cup f$ is a closed disk with boundary $\beta\cup\rho$, and the counterclockwise
orientations of $\partial D$ and $\partial f$ traverse $\alpha$ in opposite directions.
\end{lemma}
\begin{proof}
The arcs $\alpha,\beta,\rho$ join $a$ to $b$ and meet pairwise only in $\{a,b\}$, since $\beta\cap\rho\subset D\cap f=\alpha$.  By the
theta-curve lemma \cite[\S2.1]{MT}, $\R^2\setminus(\alpha\cup\beta\cup\rho)$ has exactly three components, bounded by $\alpha\cup\beta$,
$\alpha\cup\rho$ and $\beta\cup\rho$, and each arc lies on the boundary of exactly two of them.  The interior of $D$ is connected,
disjoint from $\alpha\cup\beta\cup\rho$ and has boundary $\alpha\cup\beta$, so it is the first component; likewise
$\operatorname{int}f$ is the second.  Both are bounded, so the third is the unbounded component, i.e.\ the exterior of the
Jordan curve $\beta\cup\rho$.  Hence the interior of $\beta\cup\rho$ is $\operatorname{int}D\cup\operatorname{int}f\cup(\alpha\setminus\{a,b\})$,
its closure $D\cup f$ is a closed disk by the Jordan--Schoenflies theorem \cite{Thomassen}, and near an interior point
of $\alpha$ the disk $D$ lies on one side and $f$ on the other.  A counterclockwise Jordan curve has its interior on the
left, so the two orientations traverse $\alpha$ oppositely.
\end{proof}

By Lemma~\ref{lem:glue} (with $D=D_{j-1}$, $f=f_j$), $D_j$ is a closed disk whose boundary cycle is obtained from that of
$D_{j-1}$ by replacing $\alpha_j$ with the complementary arc $\rho_j$ of $\partial f_j$, and the counterclockwise boundaries
of $f_j$ and $D_{j-1}$ traverse $\alpha_j$ in opposite directions.  The hypotheses of the lemma hold: an arc is never a whole boundary cycle, and $\alpha_j\subset\partial f_j\cap\partial D_{j-1}$.
Indeed, faces have disjoint interiors and each face is the closure of its interior, so an open set meeting $f_i$ meets
$\operatorname{int}f_i$; hence $\operatorname{int}f_j\cap D_{j-1}=\emptyset$ and $\alpha_j\subset\partial f_j$; and a point of $\alpha_j$ in
$\operatorname{int}D_{j-1}$ would have a neighbourhood inside $D_{j-1}$, which meets $\operatorname{int}f_j$ because the point lies in
$f_j$.  Let $a_j$ and $b_j$ be the initial and final vertex of
$\alpha_j$ for the counterclockwise orientation of $\partial D_{j-1}$, and let $a_1$ be any vertex of $f_1$.

\begin{example}\label{ex:shellings}
(1) The grid on $[0,m]\times[0,n]$ with its unit squares taken row by row, from left to right.  (2) The
\emph{triangulated grid}, each square $[i,i+1]\times[j,j+1]$ divided by the diagonal from $(i,j)$ to $(i+1,j+1)$
into a lower triangle $T^-_{ij}$ and an upper triangle $T^+_{ij}$: take the squares row by row from left to right,
and in each square $T^-_{ij}$ first if $i=0$ and $T^+_{ij}$ first if $i>0$.  Then $T^-_{0j}$ is attached along its
bottom edge (or is $f_1$), $T^+_{0j}$ along the diagonal, $T^+_{ij}$ ($i>0$) along its left edge, and $T^-_{ij}$ ($i>0$) along
the diagonal together with its bottom edge when $j>0$; these intersections are arcs.  By a theorem of Danaraj and Klee,
every triangulation of a $2$-disk is shellable \cite{DK78}.  (3) A fan of triangles
$[o,p_k,p_{k+1}]$ around an interior vertex $o$, attached in cyclic order, the last one along two edges.  All
three are checked combinatorially in \texttt{check\_arc\_calculus.py}.
\end{example}

\begin{theorem}[finite non-abelian Stokes theorem; lasso decomposition \cite{Arefeva,Bralic80,Levy03,Levy10}]\label{thm:stokes}
Let $f_1,\dots,f_N$ be a shelling of $K$, $U$ an $\SU(2)$-cochain, and $F_j:=\hol(\partial f_j)$ at $a_j$.  For every
vertex $p$ of $\partial K$ there are edge paths $\tau_j$ from $p$ to $a_j$ in the $1$-skeleton of $K$ such that the
holonomy of $\partial K$ at $p$ is
\[
\hol(\partial K)=L_1L_2\cdots L_N,\qquad L_j:=\hol(\tau_j)^{-1}F_j\,\hol(\tau_j).
\]
Each face contributes exactly one factor, in the order of the shelling from left to right.  The identity is
exact for every cochain; no limit is involved.
\end{theorem}
\begin{proof}
We prove by induction on $j$ the statement $(S_j)$: for every vertex $q$ of $\partial D_j$ there are edge paths
$\tau_{j,i}$ ($i\le j$) from $q$ to $a_i$ in the $1$-skeleton of $D_j$ with
$\hol(\partial D_j\text{ at }q)=\prod_{i=1}^{j}\hol(\tau_{j,i})^{-1}F_i\hol(\tau_{j,i})$ (increasing $i$ from left to
right).  The theorem is $(S_N)$.

\emph{Change of base point.}  If $(S_j)$ holds at one vertex $a$ of $\partial D_j$, it holds at every vertex $q$ of
$\partial D_j$: let $\sigma'$ be the counterclockwise boundary path from $a$ to $q$; by \eqref{eq:rebase},
$\hol(\partial D_j\text{ at }q)=\hol(\sigma')\hol(\partial D_j\text{ at }a)\hol(\sigma')^{-1}$, and conjugating every
factor by $\hol(\sigma')$ replaces each tail $\tau$ by $\bar\sigma'*\tau$, because
$\hol(\bar\sigma'*\tau)=\hol(\tau)\hol(\sigma')^{-1}$.

$(S_1)$ holds at $a_1$ with the trivial tail, hence everywhere.  Assume $(S_{j-1})$ and write $a=a_j$,
$b=b_j$, $\alpha=\alpha_j$, $\rho=\rho_j$.  Based at $a$, the counterclockwise boundaries are
$\partial D_{j-1}=\alpha*\beta$ with $\beta$ from $b$ to $a$, $\partial f_j=\rho*\bar\alpha$ with $\rho$ from $a$ to $b$,
and $\partial D_j=\rho*\beta$.  Lemma~\ref{lem:cut}, applied to $\gamma=\rho*\beta$ with the cut $\delta=\bar\alpha$
from $b$ to $a$, gives
\[
\hol(\partial D_j\text{ at }a)=\hol(\alpha*\beta)\,\hol(\rho*\bar\alpha)=\hol(\partial D_{j-1}\text{ at }a)\,F_j .
\]
Since $a$ is a vertex of $\partial D_{j-1}$, $(S_{j-1})$ at $a$ expresses the first factor with tails in $D_{j-1}$;
take $\tau_{j,j}$ trivial.  This is $(S_j)$ at $a$, and by the change of base point at every vertex of $\partial D_j$.
\end{proof}

For a rectangular grid one can also choose the tails along a comb and obtain a different explicit order; we
record it because it is convenient in computations.

\begin{proposition}\label{prop:grid}
Let $K$ be the grid on $[0,m]\times[0,n]$ with unit square faces $f_{ij}$ (lower left corner $(i,j)$), $H_{ij}$ the
holonomy of $\partial f_{ij}$ at $(i,j)$, and $\tau_{ij}$ the path from $(0,0)$ up to $(0,j)$ and then right to $(i,j)$.
Then, at $(0,0)$,
\[
\hol(\partial K)=\prod_{j=n-1}^{0}\ \prod_{i=0}^{m-1}\ \hol(\tau_{ij})^{-1}H_{ij}\,\hol(\tau_{ij}),
\]
both products written from left to right in the indicated order.
\end{proposition}
\begin{proof}
Let $P_j$ be the holonomy of the boundary of row $j$ at $(0,j)$, and $Q_j$ that of the boundary of rows
$j,\dots,n-1$ at $(0,j)$.  Based at $(0,j+1)$, the counterclockwise boundary of rows $j,\dots,n-1$ is $\alpha*\beta$, with
$\alpha$ running down to $(0,j)$, right to $(m,j)$ and up to $(m,j+1)$; with the cut $\delta$ from $(m,j+1)$ left to
$(0,j+1)$, $\alpha*\delta$ is the boundary of row $j$ and $\bar\delta*\beta$ that of rows $j+1,\dots,n-1$, both at $(0,j+1)$.
Lemma~\ref{lem:cut} and \eqref{eq:rebase} along the edge from $(0,j+1)$ down to $(0,j)$, which is the first edge of
both loops, give $Q_j=u_j^{-1}Q_{j+1}u_j\,P_j$ with $u_j=U_{(0,j+1)(0,j)}$.  Inside row $j$, let $R_i$ be the holonomy
at $(i,j)$ of the boundary of $[i,m]\times[j,j+1]$.  Cutting along $x=i+1$ from $(i+1,j+1)$ down to $(i+1,j)$ with
base $(i+1,j)$, and rebasing along the bottom edge $w$ from $(i,j)$ to $(i+1,j)$, which is the first edge of both
loops at $(i,j)$, gives $R_i=H_{ij}\,U_w^{-1}R_{i+1}U_w$.  Hence $P_j=R_0=\Lambda_{0j}\Lambda_{1j}\cdots\Lambda_{m-1,j}$ with
$\Lambda_{ij}=\hol(\rho_{ij})^{-1}H_{ij}\hol(\rho_{ij})$ and $\rho_{ij}$ the bottom path from $(0,j)$ to $(i,j)$.
Unwinding, $Q_0$ is the stated product, since $\hol(\tau_{ij})=\hol(\rho_{ij})\,u_{j-1}\cdots u_0$.
\end{proof}

In the abelian case the conjugations disappear and Theorem~\ref{thm:stokes} reduces to
$\sum_f\omega(\partial f)=\omega(\partial K)$.

\begin{corollary}[discrete fundamental theorem]\label{cor:poincare}
Let $K$ admit a shelling.  An $\SU(2)$-cochain $U$ has trivial holonomy around every face iff there is
$x\colon K^0\to M$ with $U_{qp}=x_q\bar x_p$ for every edge, i.e.\ $U_{qp}$ is the left transport $\ell_{x_q,x_p}$.
The map $x$ is unique up to a right transport $x\mapsto xb$.
\end{corollary}
\begin{proof}
If $U_{qp}=x_q\bar x_p$, every holonomy telescopes to $1$ by \eqref{eq:telescope}.  Conversely, define $x$ on the
vertices of $f_1$ by choosing $x_{a_1}$ and transporting along $\partial f_1$; the last edge is consistent because
$\hol(\partial f_1)=1$.  Suppose $x$ is defined on $D_{j-1}$ with $U_{qp}=x_q\bar x_p$ on its edges.  The vertices of
$f_j$ not in $D_{j-1}$ are the interior vertices of $\rho_j$; define $x$ there by transporting from $a_j$ along
$\rho_j$.  For the last edge of $\rho_j$, flatness of $f_j$ gives $\hol(\rho_j)=\hol(\alpha_j)$, and
$\hol(\alpha_j)x_{a_j}=x_{b_j}$ by the induction hypothesis, so $\hol(\rho_j)x_{a_j}=x_{b_j}$.  Every edge of $f_j$ lies in
$\alpha_j$ or $\rho_j$.  Uniqueness: if $x,x'$ are two solutions, $\bar x_px'_p=\bar x_qx'_q$ along every edge, and the
$1$-skeleton is connected.
\end{proof}

\begin{proposition}[plaquette expansion; standard]\label{prop:plaquette}
Let $A_1,A_2$ be $C^2$ maps from an open set $O\subset\R^2$ to $\im\HH$, and on the grid $h\Z^2$ put $U=\exp(hA_1)$ on
horizontal and $\exp(hA_2)$ on vertical edges, evaluated at edge midpoints.  The holonomy of the plaquette with lower left corner $(s,t)$, at
that corner, is $\exp\bigl(h^2F_{12}(s,t)+O(h^3)\bigr)$ with $F_{12}=\partial_1A_2-\partial_2A_1-[A_1,A_2]$.
\end{proposition}
\begin{proof}
The plaquette holonomy is $e^{X_4}e^{X_3}e^{X_2}e^{X_1}$ with $X_1=hA_1(s+\frac h2,t)$,
$X_2=hA_2(s+h,t+\frac h2)$, $X_3=-hA_1(s+\frac h2,t+h)$, $X_4=-hA_2(s,t+\frac h2)$.  By the
Baker--Campbell--Hausdorff formula its logarithm is $\sum X_i+\tfrac12\sum_{i>j}[X_i,X_j]+O(h^3)$.  Taylor
expansion gives $\sum X_i=h^2(\partial_1A_2-\partial_2A_1)(s,t)+O(h^3)$, and with $X_3=-X_1+O(h^2)$,
$X_4=-X_2+O(h^2)$ the commutator sum is $2[X_2,X_1]+O(h^3)=-2h^2[A_1,A_2]+O(h^3)$.
\end{proof}

So the face holonomies in Proposition~\ref{prop:grid} are $1+h^2F_{12}+O(h^3)$, and
Corollary~\ref{cor:poincare} is the discrete form of Theorem~\ref{thm:ftc}(iii).

\subsection{The half-displacement cochain}
Let $x\colon K^0\to M$ be a vertex map with $x_q\ne-x_p$ along edges.  Its displacement cochain
$g_{qp}=x_q\bar x_p$ is flat.  Its \emph{half-displacement cochain} is
\[
h_{qp}:=(x_q\bar x_p)^{1/2}=\frac{(x_p+x_q)\bar x_p}{\lVert x_p+x_q\rVert}=\frac{1+x_q\bar x_p}{2\sqrt{\kappa(x_p,x_q)}} ,
\]
an $\SU(2)$-cochain ($h_{pq}=h_{qp}^{-1}$).  By Proposition~\ref{prop:LCarc}, $\Ad(\hol_h(\gamma))$ is the
Levi-Civita parallel transport, in the right trivialisation, along the polygon of minimal arcs through the
vertices of $\gamma$.  We call $\hol_h(\gamma)\in M$ the \emph{spin holonomy} of the polygon and
$\Theta(a):=2d(1,a)=2\arccos\re a\in[0,2\pi]$ the \emph{spin angle} of $a\in M$; $\Ad(a)$ is the rotation by
$\Theta(a)$ about $\im a$.  Since $d$ is bi-invariant,
\begin{equation}\label{eq:subadd}
\Theta(ab)\le\Theta(a)+\Theta(b),\qquad\Theta(bab^{-1})=\Theta(a)=\Theta(a^{-1}).
\end{equation}

\begin{theorem}[minimal-arc triangles; cf.\ \cite{SB88,MS93}]\label{thm:triangle}
Let $x,y,z\in M$ be pairwise non-antipodal, $c_{xy}=c(x,y)$, $\kappa_{xy}=\kappa(x,y)$, etc., and let
$a=h_{xz}h_{zy}h_{yx}$ be the spin holonomy of the triangle $x\to y\to z\to x$ at $x$.  Then
\begin{equation}\label{eq:tri}
a=\frac{(1+c_{xy}+c_{yz}+c_{zx})+\im(y\bar x)\times\im(z\bar x)}{4\sqrt{\kappa_{xy}\kappa_{yz}\kappa_{zx}}},
\qquad
\re a=\frac{\kappa_{xy}+\kappa_{yz}+\kappa_{zx}-1}{2\sqrt{\kappa_{xy}\kappa_{yz}\kappa_{zx}}} .
\end{equation}
$a$ is unchanged if $x,y,z$ are replaced by $xb,yb,zb$.  If $x,y,z$ are linearly independent, let $E\in(0,2\pi)$
be the spherical excess (area) of the geodesic triangle $xyz$ in the great sphere $M\cap\operatorname{span}(x,y,z)$.
Then $\Theta(a)=E$,
\[
\tan\tfrac E2=\frac{\lVert x\wedge y\wedge z\rVert}{1+c_{xy}+c_{yz}+c_{zx}},\qquad
\cos E=\frac{(\kappa_{xy}+\kappa_{yz}+\kappa_{zx}-1)^2}{2\kappa_{xy}\kappa_{yz}\kappa_{zx}}-1,
\]
and $\im a$ is normal, in the right trivialisation at $x$, to the tangent plane of that great sphere.  If
$y=\exp(Y)x$, $z=\exp(Z)x$ with $\lvert Y\rvert,\lvert Z\rvert\le r$, then
$a=\exp\bigl(\tfrac14Y\times Z+O(r^4)\bigr)=\exp\bigl(\tfrac18[Y,Z]+O(r^4)\bigr)$.
\end{theorem}
\begin{proof}
Each factor is invariant under right translation: $(yb\,\overline{xb})^{1/2}=(y\bar x)^{1/2}$.  So we may take
$x=1$.  Then $h_{yx}=(1+y)/\lVert1+y\rVert$, $h_{zy}=(1+z\bar y)/\lVert y+z\rVert$, $h_{xz}=(1+\bar z)/\lVert1+z\rVert$,
and
\[
(1+\bar z)(1+z\bar y)(1+y)=2+(y+\bar y)+(z+\bar z)+(z\bar y+\bar zy).
\]
Writing $y=y_0+Y'$, $z=z_0+Z'$ with $Y',Z'\in\im\HH$, one has $z\bar y+\bar zy=2c_{yz}-2Z'\times Y'$, so the product
equals $2(1+c_{xy}+c_{yz}+c_{zx})+2\,Y'\times Z'$.  Since $\lVert u+w\rVert=2\sqrt{\kappa(u,w)}$, this is
\eqref{eq:tri}; the formula for $\re a$ uses $1+c=2\kappa$.

The Gram determinant of $x,y,z$ is $\lVert x\wedge y\wedge z\rVert^2=1+2c_{xy}c_{yz}c_{zx}-c_{xy}^2-c_{yz}^2-c_{zx}^2$, and
a direct expansion gives
\[
(1+c_{xy}+c_{yz}+c_{zx})^2+\lVert x\wedge y\wedge z\rVert^2=2(1+c_{xy})(1+c_{yz})(1+c_{zx})=16\,\kappa_{xy}\kappa_{yz}\kappa_{zx}.
\]
Hence $\lvert\im a\rvert=\lVert x\wedge y\wedge z\rVert/(4\sqrt{\kappa\kappa\kappa})$, and since $\lvert\im a\rvert\ge0$,
$\tfrac12\Theta(a)$ is the angle in $[0,\pi]$ with tangent $\lVert x\wedge y\wedge z\rVert/(1+c_{xy}+c_{yz}+c_{zx})$ and
the sign of $\cos$ given by the denominator.  In the three-dimensional space $\operatorname{span}(x,y,z)$,
$\lVert x\wedge y\wedge z\rVert=\lvert\det(x,y,z)\rvert$, and the formula of Van~Oosterom and Strackee
\cite{VOS,Eriksson} identifies this angle with half the solid angle of the cone over the triangle, which by
Girard's theorem is $E/2$.  The formula for $\cos E=2\cos^2(E/2)-1$ follows.  For $x=1$ the tangent plane of the
great sphere is spanned by $Y'$ and $Z'$, and $\im a$ is parallel to $Y'\times Z'$.

Finally, with $x=1$, $Y'=\sin\lvert Y\rvert\,\hat Y=Y+O(r^3)$, likewise for $Z'$, and the denominator in
\eqref{eq:tri} and $1+\sum c$ are $4+O(r^2)$; so $\im a=\tfrac14Y\times Z+O(r^4)$, and
$\log a=\im a\,(1+O(\lvert\im a\rvert^2))$.
\end{proof}

The formula for $\cos(E/2)$ in \eqref{eq:tri} is, in terms of side lengths $\alpha,\beta,\gamma$, the classical
$\cos\frac E2=(1+\cos\alpha+\cos\beta+\cos\gamma)/(4\cos\frac\alpha2\cos\frac\beta2\cos\frac\gamma2)$
\cite{Todhunter}.  The whole of \eqref{eq:tri} has the algebraic form of the Bargmann--Pancharatnam three-point invariant: for
unit vectors $n_1,n_2,n_3\in\R^3$ (Bloch vectors) and the corresponding projections,
$\operatorname{Tr}(P_1P_2P_3)=\tfrac14\bigl(1+\sum_{i<j}n_i\cdot n_j+i\,n_1\cdot(n_2\times n_3)\bigr)$, of modulus $\sqrt{\beta_{12}\beta_{23}\beta_{31}}$ with
$\beta_{ij}=\tfrac12(1+n_i\cdot n_j)$
and argument $+\Omega/2$, where $\Omega$ is the solid angle of the geodesic triangle signed by $\operatorname{sign}\det(n_1,n_2,n_3)$
(positive when $n_1\to n_2\to n_3$ runs counterclockwise seen from outside the sphere) \cite{SB88,MS93}; applied to $x,y,z$ in the three-dimensional space $\operatorname{span}(x,y,z)$, with $i$
replaced by the unit normal, it is \eqref{eq:tri} up to the factor $\sqrt{\kappa\kappa\kappa}$, with the same sign convention:
$a=\cos\frac E2+\sigma\sin\frac E2\,\nu\bar x$ with $\sigma=\operatorname{sign}\det(x,y,z,\nu)$ (Corollary~\ref{cor:excess}(ii)).  This is an identity of
algebraic form on $\operatorname{span}(x,y,z)$; it is not the Bargmann invariant of the Hopf images of $x,y,z$, which is a different number
(for $x=1$, $y=e^{i/2}$, $z=e^{j/2}$ one has $E=0.1302$, while $y$ commutes with $i$, so the Hopf images $xi\bar x$ and $yi\bar y$ coincide and their
three-point invariant is real).  What the heap formulation adds is
the reading of this invariant as the product of the three Levi-Civita half-displacements on $S^3$
(Proposition~\ref{prop:LCarc}), so that Theorem~\ref{thm:stokes} applies to it.  The small-triangle expansion agrees with the curvature $R(X,Y)Z=-\tfrac14[[X,Y],Z]$: $\Ad(a)$ is
the rotation by $\tfrac12\lvert Y\times Z\rvert+O(r^4)$, the area of the triangle to leading order.

\begin{lemma}[orientation]\label{lem:sign}
Let $W\subset\R^4$ be a $3$-dimensional subspace with unit normal $\nu$, and $x,y,z\in M\cap W$.  Then $\nu\bar x\in\im\HH$
and
\[
\im(y\bar x)\times\im(z\bar x)=\det(x,y,z,\nu)\;\nu\bar x ,
\]
the determinant taken in the basis $(1,i,j,k)$.  If $y\ne-x$, the half-displacement $h=(y\bar x)^{1/2}$ satisfies
$h(\nu\bar x)\bar h=\nu\bar y$.
\end{lemma}
\begin{proof}
$\re(\nu\bar x)=\ip\nu x=0$.  Right multiplication by $\bar x$ is in $\SO(4)$, so
$\det(x,y,z,\nu)=\det(1,y\bar x,z\bar x,\nu\bar x)=\det_{\im\HH}(Y',Z',N')=\ip{Y'\times Z'}{N'}$ with $Y'=\im(y\bar x)$,
$Z'=\im(z\bar x)$, $N'=\nu\bar x$.  The space $W\bar x$ contains $1,y\bar x,z\bar x$ and is orthogonal to $N'$, so $Y',Z'\perp N'$
and $Y'\times Z'$ is parallel to the unit vector $N'$, with coefficient $\ip{Y'\times Z'}{N'}$.  For the last claim,
the minimal arc $\gamma$ from $x$ to $y$ lies in $W$, so the constant vector $\nu$ is tangent to $M$ along $\gamma$, and its
covariant derivative along $\gamma$, the tangential part of its ordinary derivative, vanishes.  Hence $\nu$ is parallel,
and Proposition~\ref{prop:LCarc}(ii) with $\xi u=\nu$, $\xi=\nu\bar x$, gives $(h\nu\bar x\bar h)\,y=\nu$.
\end{proof}

\begin{corollary}[excess bound and Gauss--Bonnet]\label{cor:excess}
Let $K$ be a triangulated disk with a shelling (Example~\ref{ex:shellings}), and $x$ a
vertex map with non-antipodal neighbours.  Let $H_\partial$ be the spin holonomy of the boundary polygon of
minimal arcs, and $E_T:=\Theta(a_T)$, which is the excess of the image of $T$ when its vertices are linearly
independent (and $0$ or $2\pi$ otherwise).
\begin{enumerate}[label=(\roman*)]
\item $\Theta(H_\partial)\le\sum_TE_T$, and the rotation angle $\varrho(H_\partial):=\min(\Theta,2\pi-\Theta)(H_\partial)$ of the Levi-Civita
holonomy satisfies $\varrho(H_\partial)\le\sum_T\min(E_T,2\pi-E_T)\le\sum_T\operatorname{area}(T)$, where $\operatorname{area}(T)$ is the
area of the geodesic triangle (zero for linearly dependent vertices, where $E_T\in\{0,2\pi\}$).
\item If $x(K^0)$ lies in a great $2$-sphere $\Sigma=M\cap W$ with unit normal $\nu\perp W$, then
$H_\partial=\cos\bigl(\tfrac12\sum_T\sigma_TE_T\bigr)+\sin\bigl(\tfrac12\sum_T\sigma_TE_T\bigr)\,\nu\bar p$ at the base vertex
$p$, where $\sigma_T=\operatorname{sign}\det(x_T,y_T,z_T,\nu)\in\{-1,0,1\}$ for the counterclockwise vertex order of $T$.
If moreover $x$ extends to a homeomorphism of $K$ onto a region $R\subset\Sigma$ mapping each edge onto the minimal arc
and each triangle onto the geodesic triangle of its vertices, then $\re H_\partial=\cos\bigl(\tfrac12\operatorname{area}(R)\bigr)$.
\end{enumerate}
\end{corollary}
\begin{proof}
(i) By Theorem~\ref{thm:stokes}, applied to the half-displacement cochain, $H_\partial=L_1\cdots L_N$ where $L_j$ is a
conjugate of the holonomy $F_j$ of the triangle $f_j$ at one of its vertices.  By Theorem~\ref{thm:triangle} applied
with that vertex first (a cyclic relabelling, which keeps the counterclockwise order), $\Theta(F_j)=E_{f_j}$.
Apply \eqref{eq:subadd}.  For the second bound, $\varrho(a)=2\min\bigl(d(a,1),d(-a,1)\bigr)$ is twice the distance from $[a]$ to the
identity in the bi-invariant metric of $\SO(3)=S^3/\{\pm1\}$, so it is subadditive and conjugation invariant, and
$\varrho(F_j)=\min(E_{f_j},2\pi-E_{f_j})$.  This is at most $E_T=\operatorname{area}(T)$ for independent vertices and $0$ otherwise.

(ii) Let $T=f_j$ with vertices $x,y,z$ in counterclockwise order starting at $a_j=x$.  By \eqref{eq:tri} and
Lemma~\ref{lem:sign},
\[
F_j=\frac{(1+c_{xy}+c_{yz}+c_{zx})+\det(x,y,z,\nu)\,\nu\bar x}{4\sqrt{\kappa_{xy}\kappa_{yz}\kappa_{zx}}}
=\cos\tfrac{E_T}2+\sigma_T\sin\tfrac{E_T}2\;\nu\bar x ,
\]
using $\re F_j=\cos(E_T/2)$ and $\lvert\im F_j\rvert=\sin(E_T/2)\ge0$ (if $\sigma_T=0$, $F_j=\pm1$ and the formula holds
with $E_T\in\{0,2\pi\}$).  The tail $\tau_j$ runs from $p$ to $x$ along minimal arcs in $\Sigma$, so by the last claim
of Lemma~\ref{lem:sign}, applied edge by edge, $\Ad(\hol(\tau_j))(\nu\bar p)=\nu\bar x$.  Hence
$L_j=\cos\tfrac{E_T}2+\sigma_T\sin\tfrac{E_T}2\,\nu\bar p$.  All $L_j$ lie in the one-parameter group
$\{\cos\theta+\sin\theta\,\nu\bar p\}$, and their product is the stated element.  Under the additional hypothesis,
every triangle is mapped onto a geodesic triangle in $R$ with the orientation of $\Sigma$ induced by $K$; with the
orientation of $\Sigma$ for which $(t_1,t_2)$ is positive at $x$ iff $\det(x,t_1,t_2,\nu)>0$, a geodesic triangle traversed
counterclockwise has its second side to the left of the first, so $\det(x,y,z,\nu)=\det(x,P_xy,P_xz,\nu)$ has one sign
for all $T$.  Thus all $\sigma_T$ are equal, and $\sum_TE_T=\operatorname{area}(R)$ by Girard's theorem and additivity of
area, since the triangles cover $R$ with disjoint interiors; $\cos$ is even.
\end{proof}

Statement (ii) is the exact discrete Gauss--Bonnet theorem for geodesic polygons, at any mesh size; (i) is its
non-abelian replacement in $S^3$, where triangles in different great spheres contribute non-commuting factors.
We now pass to the continuum limit.

\begin{theorem}[continuum limit]\label{thm:continuum}
Let $\gamma\colon[0,L]\to M$ be continuous and of class $C^3$ on each interval of a partition $0=u_0<\dots<u_r=L$; let
$\lambda=\dot\gamma\bar\gamma$ (one-sided at the $u_j$) and let $U\colon[0,L]\to M$ solve $\dot U=\tfrac12\lambda U$, $U(0)=1$.
\begin{enumerate}[label=(\roman*)]
\item In the right trivialisation, Levi-Civita parallel transport along $\gamma$ from $\gamma(0)$ to $\gamma(u)$ is $\Ad(U(u))$;
$U(L)$ is the lift to $\SU(2)$ of the holonomy obtained by continuation from $U(0)=1$.
\item There are $\delta_0>0$ and $C$, depending only on $L$, $r$ and the $C^3$ norms of $\gamma$ on the pieces, such that for
every partition $\pi=\{t_0<\dots<t_n\}$ of $[0,L]$ with mesh $\delta\le\delta_0$ the spin holonomy $H_\pi=h_n\cdots h_1$,
$h_k=(\gamma(t_k)\bar\gamma(t_{k-1}))^{1/2}$, is defined and satisfies $\lvert H_\pi-U(L)\rvert\le C\delta^2$.  The partition need not
contain the break points $u_j$.
\end{enumerate}
\end{theorem}
\begin{proof}
(i) In the right trivialisation $\nabla^R$-parallel fields are constant, $\xi'=0$, and $\nabla^L$-parallel fields are
$\xi=\Ad(g)\xi_0$ with $g=\gamma\bar\gamma(0)$, $\dot g=\lambda g$, so $\xi'=[\lambda,\xi]$ (Proposition~\ref{prop:LCarc}(i)).  The
connection forms average, so Levi-Civita parallel fields satisfy $\xi'=\tfrac12[\lambda,\xi]$, which $\Ad(U)\xi_0$ solves.
(ii) Choose $\delta_0$ so that $d(\gamma(s),\gamma(t))<\pi/2$ for $\lvert s-t\rvert\le\delta_0$; then all $h_k$ are defined.  On a step
$[t_{k-1},t_k]$ of length $\Delta\le\delta$ inside one piece, $g_k:=\gamma(t_k)\bar\gamma(t_{k-1})$ is the value at $t_k$
of the solution of $\dot V=\lambda V$, $V(t_{k-1})=1$.  By the Magnus expansion \cite{Magnus,IMNZ},
$\log g_k=\int_{t_{k-1}}^{t_k}\lambda+\tfrac12\int\!\!\int_{s_2<s_1}[\lambda(s_1),\lambda(s_2)]+O(\Delta^3)=\Delta\lambda(m_k)+O(\Delta^3)$, $m_k$ the
midpoint: the midpoint rule has error $O(\Delta^3)$ and $\lvert[\lambda(s_1),\lambda(s_2)]\rvert=O(\lvert s_1-s_2\rvert)$.  So
$h_k=\exp(\tfrac12\log g_k)=\exp(\tfrac12\Delta\lambda(m_k)+O(\Delta^3))$, and the same expansion for $\tfrac12\lambda$ gives
$U(t_k)U(t_{k-1})^{-1}=\exp(\tfrac12\Delta\lambda(m_k)+O(\Delta^3))$.  For unit quaternions
$\lvert a_n\cdots a_1-b_n\cdots b_1\rvert\le\sum_k\lvert a_k-b_k\rvert$.  At most $r-1$ steps contain a break point in their
interior; on such a step $\lambda$ is bounded and piecewise continuous, so both $g_k$ and $U(t_k)U(t_{k-1})^{-1}$ (for $\tfrac12\lambda$)
equal $1+\int\lambda+O(\Delta^2)$, resp.\ $1+\tfrac12\int\lambda+O(\Delta^2)$, and $h_k=\exp(\tfrac12\log g_k)$ differs from the exact
step by $O(\Delta^2)$.  Hence $\lvert H_\pi-U(L)\rvert\le C\sum_k\Delta_k^3+C(r-1)\delta^2\le C'\delta^2$.
\end{proof}

\begin{corollary}[continuum excess bound]\label{cor:area}
Let $F\colon[0,1]^2\to M$ be of class $C^3$, $\gamma$ its boundary curve traversed counterclockwise, and $U_\gamma$ as in
Theorem~\ref{thm:continuum}.  Then
\[
\Theta(U_\gamma)\le\operatorname{area}(F):=\int_{[0,1]^2}\lvert\partial_1F\wedge\partial_2F\rvert ;
\]
in particular the Levi-Civita holonomy of $\gamma$ is a rotation by at most $\operatorname{area}(F)$.
\end{corollary}
\begin{proof}
Use the triangulated $n\times n$ grid (Example~\ref{ex:shellings}) and the vertex map $F$; for large $n$ neighbours are not
antipodal.  By Corollary~\ref{cor:excess}(i), $\Theta(H_\partial(n))\le\sum_TE_T$, and by Theorem~\ref{thm:continuum}(ii) and
continuity of $\Theta$, $\Theta(H_\partial(n))\to\Theta(U_\gamma)$.  For the lower triangle at $p$, with $\lambda_i=\partial_iF\,\bar F$,
$Y=\log(F(p+e_1/n)\bar F(p))=\lambda_1/n+O(n^{-2})$ and $Z=(\lambda_1+\lambda_2)/n+O(n^{-2})$, so by Theorem~\ref{thm:triangle}
($\re a>0$, $E=2\arcsin\lvert\im a\rvert$) $E_T=\tfrac12\lvert Y\times Z\rvert+O(n^{-4})=\tfrac1{2n^2}\lvert\lambda_1\times\lambda_2\rvert(p)+O(n^{-3})$;
the upper triangle gives the same.  Since right multiplication by $F$ is an isometry,
$\lvert\lambda_1\times\lambda_2\rvert=\lvert\partial_1F\wedge\partial_2F\rvert$, so $\sum_TE_T$ is a Riemann sum for
$\operatorname{area}(F)$ up to $O(1/n)$.  The rotation angle of $\Ad(U_\gamma)$ is $\min(\Theta,2\pi-\Theta)\le\Theta$.
\end{proof}

For a piecewise $C^1$ closed curve $\gamma$ the spin holonomy $U_\gamma$ is defined as in Theorem~\ref{thm:continuum}, by
$\dot U=\tfrac12\lambda U$, $U(0)=1$, $\lambda=\dot\gamma\bar\gamma$ piecewise continuous.  On a great $2$-sphere $\Sigma=S^3\cap\nu^\perp$ we use the
orientation in which a basis $(t_1,t_2)$ of $T_x\Sigma$ is positive iff $\det(x,t_1,t_2,\nu)>0$, as in Corollary~\ref{cor:excess}(ii).

\begin{theorem}[holonomy of a disk in a great $2$-sphere; classical]\label{thm:disk}
Let $R\subset\Sigma$ be a closed disk whose boundary is a piecewise $C^1$ Jordan curve $\gamma$, traversed with $R$ on its left
(positively) and based at $x_0$, and let $A=\operatorname{area}(R)$.  Then
\[
U_\gamma=\exp\bigl(\tfrac A2\,\nu\bar x_0\bigr)=\cos\tfrac A2+\sin\tfrac A2\,\nu\bar x_0 ;
\]
for the opposite orientation $U_\gamma=\exp(-\tfrac A2\,\nu\bar x_0)$.  In particular $\re U_\gamma=\cos(A/2)$, and $\Theta(U_\gamma)=A$ for
$A\le2\pi$, $\Theta(U_\gamma)=4\pi-A$ for $A\ge2\pi$.
\end{theorem}
\begin{proof}
(1) \emph{Connection and curvature.}  Write $\theta=dx\,\bar x$ for the right Maurer--Cartan form ($\im\HH$-valued) and $A=\tfrac12\theta$;
then $U_\gamma$ is the holonomy of $\dot U=A(\dot\gamma)U$, and by Proposition~\ref{prop:LCarc} and Theorem~\ref{thm:continuum}(i),
$\Ad(U)$ is Levi-Civita transport.  With the curvature $\Omega(X,Y)=X A(Y)-YA(X)-A([X,Y])-[A(X),A(Y)]$, the convention of
Proposition~\ref{prop:plaquette}, the Maurer--Cartan equation $d\theta(X,Y)=[\theta(X),\theta(Y)]$ (Theorem~\ref{thm:ftc}(iii)) gives
$\Omega=\tfrac12[\theta,\theta]-\tfrac14[\theta,\theta]$, i.e.\ $\Omega(X,Y)=\tfrac14[\theta(X),\theta(Y)]=\tfrac12\,\theta(X)\times\theta(Y)$.
(2) \emph{A parallel section.}  On $\Sigma$ put $n(x)=\nu\bar x$; it is imaginary because $\re(\nu\bar x)=\ip\nu x=0$.  Since
$d\bar x=-\bar x\,dx\,\bar x$, $dn=-n\theta$, and the covariant derivative for the adjoint action is
$Dn=dn-[A,n]=-n\theta-\tfrac12(\theta n-n\theta)=-\tfrac12(n\theta+\theta n)=\ip n\theta=\ip\nu{dx}$, which vanishes along $\Sigma$.  (This is
Lemma~\ref{lem:sign}.)  (3) \emph{The curvature on $\Sigma$.}  For a positive orthonormal basis $e_1,e_2$ of $T_x\Sigma$,
$\theta(e_i)=e_i\bar x$, and the computation of Lemma~\ref{lem:sign} gives $(e_1\bar x)\times(e_2\bar x)=\det(x,e_1,e_2,\nu)\,\nu\bar x=n(x)$;
so $\Omega=\tfrac12\,n\,dA_\Sigma$ with $dA_\Sigma$ the oriented area form.  (4) \emph{Abelian gauge.}  Choose $q\in\Sigma\setminus R$; the open set
$\Sigma'=\Sigma\setminus\{q\}$ is contractible, so the map $n\colon\Sigma'\to S^2\subset\im\HH$ lifts through the circle bundle
$S^3\to S^2$, $g\mapsto gn_0\bar g$: there is a smooth $g\colon\Sigma'\to S^3$ with $g\,n_0\,\bar g=n$ for a fixed unit $n_0$.  Along a curve in
$\Sigma'$ put $V=g^{-1}Ug(\gamma(0))$; then $\dot V=\hat A(\dot\gamma)V$ with $\hat A=g^{-1}Ag-g^{-1}dg$, whose curvature is
$g^{-1}\Omega g=\tfrac12n_0\,dA_\Sigma$.  Transporting the parallel field $n$ gives $[\hat A,n_0]=0$, since $g^{-1}ng=n_0$ is constant; so
$\hat A=a\,n_0$ with a real $1$-form $a$ on $\Sigma'$, and comparing curvatures, $da=\tfrac12dA_\Sigma$.  (5) \emph{Abelian Stokes.}
Along $\gamma$, $V(L)=\exp\bigl(n_0\oint_\gamma a\bigr)$, and by Green's theorem for the Jordan domain $R$ with piecewise $C^1$ boundary
(in a chart of $\Sigma'$, e.g.\ a stereographic one; a piecewise $C^1$ Jordan curve is rectifiable, and Green's theorem holds for
Jordan domains with rectifiable boundary \cite{Potts51}), $\oint_\gamma a=\int_Rda=A/2$.  Hence
$U_\gamma=g(x_0)V(L)g(x_0)^{-1}=\exp(\tfrac A2\,g(x_0)n_0\bar g(x_0))=\exp(\tfrac A2\,\nu\bar x_0)$.  Reversing $\gamma$ inverts $U_\gamma$.  The
last statement holds because $\Theta(a)=2\arccos\re a$.
\end{proof}

Theorem~\ref{thm:disk} is the classical fact that Levi-Civita holonomy of the round $2$-sphere around a region is the
rotation by its area (Gauss--Bonnet, \cite[\S4-5]{doCarmo76}), lifted to the spin connection, where the lift is fixed
because the abelian Stokes theorem is applied to the spin connection itself; it is the half-solid-angle law of the
spin-$\tfrac12$ geometric phase \cite{SB88,MS93}.  We include it because it closes the following equality case.  It agrees
with the discrete Corollary~\ref{cor:excess}(ii).

\begin{corollary}[equality for disks in great spheres]\label{cor:convex}
Let $R$ be a closed disk in a great $2$-sphere with piecewise $C^3$ Jordan boundary $\gamma$, traversed positively, and
$A=\operatorname{area}(R)\le2\pi$.  Then $\Theta(U_\gamma)=A$.  If moreover $R=F_0([0,1]^2)$ for a $C^3$ map $F_0$ injective on the interior
and parametrising $\gamma$ on the boundary, then $\Theta(U_\gamma)=A=\min_F\operatorname{area}(F)$ over $C^3$ maps $F$ of the square with
boundary curve $\gamma$.
\end{corollary}
\begin{proof}
The first claim is Theorem~\ref{thm:disk}; the second follows from Corollary~\ref{cor:area} and $\operatorname{area}(F_0)=A$.
\end{proof}

The converse of Corollary~\ref{cor:convex} rests on a continuum form of the comb order of Proposition~\ref{prop:grid}, which we
derive directly from the transport equation.  For a $C^3$ map $F\colon[0,1]^2\to M$ with coordinates $(u,v)$ put $\lambda_1=\partial_uF\,\bar F$,
$\lambda_2=\partial_vF\,\bar F$, $A_i=\tfrac12\lambda_i$ and $F_{12}=\partial_uA_2-\partial_vA_1-[A_1,A_2]$, as in Proposition~\ref{prop:plaquette}.  The spin
transport along a path $c$ is the solution of $\dot V=A(\dot c)V$, $V(0)=1$, with $A=\tfrac12\,dx\,\bar x$; it is multiplicative as in
Section~\ref{sec:stokes}.

\begin{lemma}[comb formula]\label{lem:comb}
Let $G(v,u)\in M$ be the spin transport along the comb path from $(0,0)$ up the left edge to $(0,v)$ and then along the row to
$(u,v)$: $\partial_vG(v,0)=A_2(0,v)G(v,0)$, $G(0,0)=1$, $\partial_uG(v,u)=A_1(u,v)G(v,u)$.  Let $Q(v)$ be the spin holonomy, based at
$F(0,0)$, of the counterclockwise boundary of $F|_{[0,1]\times[0,v]}$ (bottom edge, right edge up to height $v$, the row at height $v$
backwards, left edge down).  Then $Q(0)=1$, $Q(1)=U_\gamma$ for the counterclockwise boundary curve $\gamma$ of $F$ based at $F(0,0)$, and
\[
Q'(v)=W(v)\,Q(v),\qquad W(v)=\int_0^1\Ad\bigl(G(v,u)^{-1}\bigr)F_{12}(u,v)\,du .
\]
Moreover $F_{12}=\tfrac12\lambda_1\times\lambda_2$, so $\lvert F_{12}\rvert=\tfrac12\lvert\partial_uF\wedge\partial_vF\rvert$.
\end{lemma}
\begin{proof}
$A_1,A_2$ are $C^2$, so $G$ is $C^2$ by differentiable dependence on parameters.  Let $K(v)$ be the transport along the bottom edge and then
up the right edge to $(1,v)$, so $K'=A_2(1,v)K$ and $K(0)=G(0,1)$; by multiplicativity $Q(v)=G(v,1)^{-1}K(v)$.  Put $J=\partial_vG\,G^{-1}$.
Differentiating $\partial_uG=A_1G$ in $v$ gives $\partial_uJ=\partial_vA_1+[A_1,J]$, with $J(v,0)=A_2(0,v)$.  Then $Z:=J-A_2$ satisfies
$\partial_uZ=-F_{12}+[A_1,Z]$, $Z(v,0)=0$, and since $Y(u)=\Ad(G(v,u)G(v,u')^{-1})Y_0$ solves $\partial_uY=[A_1,Y]$, variation of constants gives
$Z(v,u)=-\int_0^u\Ad\bigl(G(v,u)G(v,u')^{-1}\bigr)F_{12}(u',v)\,du'$.  Hence $\partial_vG(v,1)=(A_2(1,v)+Z(v,1))G(v,1)$ and
\[
Q'=-G^{-1}\bigl(A_2(1,v)+Z(v,1)\bigr)K+G^{-1}A_2(1,v)K=-\Ad\bigl(G(v,1)^{-1}\bigr)Z(v,1)\;Q=W\,Q .
\]
Finally, $\partial_u\lambda_2-\partial_v\lambda_1=\lambda_1\lambda_2-\lambda_2\lambda_1=[\lambda_1,\lambda_2]$ (from $\partial\bar F=-\bar F\,\partial F\,\bar F$), so
$F_{12}=\tfrac12[\lambda_1,\lambda_2]-\tfrac14[\lambda_1,\lambda_2]=\tfrac12\lambda_1\times\lambda_2$, and $\lvert\lambda_1\times\lambda_2\rvert=\lvert\partial_uF\wedge\partial_vF\rvert$
because right multiplication by $\bar F$ is an isometry.
\end{proof}

Lemma~\ref{lem:comb} is the continuum non-abelian Stokes theorem in comb order \cite{Arefeva,KMR}, here proved by variation of
constants; it gives Corollary~\ref{cor:area} again: $Q$ has speed $\lvert W\rvert\le\int_0^1\lvert F_{12}\rvert\,du$, so
$\tfrac12\Theta(U_\gamma)=d(1,U_\gamma)\le\operatorname{length}(Q)\le\tfrac12\operatorname{area}(F)$.

\begin{proposition}[equality case]\label{prop:converse}
Let $F\colon[0,1]^2\to M$ be of class $C^3$, of rank $2$ on a dense subset of the square, with counterclockwise boundary curve $\gamma$.
If $\Theta(U_\gamma)=\operatorname{area}(F)$, then $F([0,1]^2)$ lies in a great $2$-sphere.
\end{proposition}
\begin{proof}
(1) By the chain above, $d(1,Q(1))=\operatorname{length}(Q)=\tfrac12\operatorname{area}(F)$, which is positive because $F$ has rank $2$ somewhere, and at most
$\pi$.  Since $\int_0^1\bigl(\int_0^1\lvert F_{12}\rvert\,du-\lvert W(v)\rvert\bigr)dv=0$ with a continuous non-negative integrand,
$\lvert W(v)\rvert=\int_0^1\lvert\Ad(G(v,u)^{-1})F_{12}(u,v)\rvert\,du$ for every $v$; by strict convexity of the Euclidean norm, for every $v$ with
$W(v)\ne0$ the continuous integrand is a non-negative multiple of $e(v):=W(v)/\lvert W(v)\rvert$ at every $u$.
(2) A curve whose length equals the distance between its endpoints is a minimising geodesic up to monotone reparametrisation; so
$Q(v)=\exp(\varphi(v)Z)$ for a unit $Z\in\im\HH$ and a non-decreasing $\varphi$.  The $C^1$ curve $Q$ stays on the great circle
$\{\exp(tZ)\}$, so $W=Q'Q^{-1}=c(v)Z$ with $c$ real, and $c\ge0$ because $\varphi$ is non-decreasing.  Hence $e(v)=Z$ whenever $W(v)\ne0$.  (3) Put $n:=\Ad(G)Z$, a $C^2$ unit field in $\im\HH$, and $\nu:=nF$, a $C^2$ map to $S^3$ with
$\ip\nu F=\re n=0$.  Where $F$ has rank $2$, $F_{12}\ne0$, so $W(v)\ne0$ and, by (1) and (2), $F_{12}/\lvert F_{12}\rvert=n$; by the computation in step~(3)
of the proof of Theorem~\ref{thm:disk}, $\lambda_1\times\lambda_2$ is a positive multiple of $\mu\bar F$, with $\mu$ the unit normal of
$\operatorname{span}(F,\partial_uF,\partial_vF)$ for which $\det(F,\partial_uF,\partial_vF,\mu)>0$.  So $\nu=\mu$ there, and $\ip\nu{\partial_uF}=\ip\nu{\partial_vF}=0$ on a
dense set, hence everywhere.  (4) For the adjoint action, $D_in:=\partial_in-[A_i,n]$; as in step~(2) of the proof of Theorem~\ref{thm:disk}, now
with $\nu$ varying, $n=\nu\bar F$ gives $D_in=\partial_i\nu\,\bar F+\ip\nu{\partial_iF}=\partial_i\nu\,\bar F$.  By construction $\partial_un=[A_1,n]$ on the square
and $\partial_vn=[A_2,n]$ on the left edge, so $\partial_u\nu=0$ everywhere and $\partial_v\nu(0,\cdot)=0$.  Thus $\nu$ is constant and $F$ lies in
$S^3\cap\nu^\perp$.
\end{proof}

\begin{remark}\label{rem:converse}
(i) The rank hypothesis cannot be dropped: the tripod loop, which runs out and back along three minimal arcs from $x$ whose directions span
$T_xS^3$, has $U_\gamma=1$ and a filling of rank at most $1$ and area $0$, so $\Theta(U_\gamma)=\operatorname{area}=0$, while $x$ and the three tips span $\R^4$.
(ii) A \emph{filling} of a closed curve $\gamma$ is a $C^3$ map $F$ of the square whose restriction to the boundary, traversed counterclockwise, is
a weakly monotone parametrisation of $\gamma$ (preimages of points are connected and the degree is one).  The proposition concerns a given
filling.  For the question whether $\Theta(U_\gamma)=\inf_F\operatorname{area}(F)$ forces a Jordan curve $\gamma$ into a
great sphere (Q4 of Section~\ref{sec:open}) one needs a filling attaining the infimum that is $C^3$ on the closed square with rank $2$ on a
dense set.  For a smooth Jordan curve in $S^3$ the following results, which we know from the account in \cite[\S4]{White16}
and have not checked in the original papers, would provide it.  Morrey's solution of the Plateau problem in a homogeneously regular
Riemannian manifold \cite{Morrey48} is a harmonic, almost conformal map of the closed disk of Sobolev class $W^{1,2}$, continuous up to the
boundary, minimising area among such maps whose boundary trace is a weakly monotone parametrisation of $\gamma$ (White states the
corresponding Douglas--Rad\'o theorem in $\R^N$ for maps continuous on the closed disk and locally Lipschitz inside, and attributes the
extension to Riemannian manifolds to Morrey).  A $C^3$ filling composed with the conformal map of the disk onto the square is such a
competitor (the Dirichlet energy is conformally invariant and area does not depend on the parametrisation), so $\inf_F\operatorname{area}(F)$ over
$C^3$ fillings is at least Morrey's minimum, which is positive for a Jordan curve because the minimiser is non-constant and almost
conformal.  In a Riemannian $3$-manifold the minimiser has no interior branch points
\cite{Osserman70,Gulliver73}, so it is an immersion inside; and it is as regular up to the boundary as the curve
\cite{Hildebrandt69,HH70}.  Composing with the conformal map of the square onto the disk (analytic on the closed square by the reflection
principle, a diffeomorphism inside) gives a $C^3$ map of the square of rank $2$ in the interior, of least area among $C^3$ fillings, and
Proposition~\ref{prop:converse} applies.  Boundary branch points, the main open point of the classical theory in dimension $3$
\cite{White16}, do not matter here, because the proposition needs rank $2$ only on a dense set.
\end{remark}


Theorem~\ref{thm:continuum} is a standard consequence of the Magnus expansion; Corollary~\ref{cor:area} is an inequality
of the classical type ``holonomy is bounded by curvature times area'', here with the constant obtained from the exact
discrete Corollary~\ref{cor:excess}, and Corollary~\ref{cor:convex} shows that the constant cannot be improved.  Together with Theorem~\ref{thm:stokes} and Proposition~\ref{prop:plaquette}, the face
products converge to the surface-ordered exponential of the continuum non-abelian Stokes theorem \cite{Arefeva,KMR}.

\section{Rational calculus}\label{sec:rational}

\begin{proposition}[Cayley chart; classical \cite{Cayley,Rodrigues,Shuster,SJ}]\label{prop:cayley}
Let $\Cay(r):=(1+r)(1-r)^{-1}=\bigl((1-\lvert r\rvert^2)+2r\bigr)/(1+\lvert r\rvert^2)$ for $r\in\im\HH$.
\begin{enumerate}[label=(\roman*)]
\item $\Cay$ is a bijection $\im\HH\to M\setminus\{-1\}$ with inverse $q\mapsto\im q/(1+\re q)$; it is the inverse
stereographic projection from $-1$, and $\Cay(u\tan\tfrac\theta2)=\cos\theta+u\sin\theta$ for unit $u$.
\item It restricts to a bijection $\Q^3\to M(\Q)\setminus\{-1\}$; hence the rational points $M(\Q)$ are dense.
\item $\Cay(r)\Cay(s)=\Cay(r\star s)$ with
$r\star s=\bigl((1-\lvert s\rvert^2)r+(1-\lvert r\rvert^2)s+2\,r\times s\bigr)/\bigl(1+\lvert r\rvert^2\lvert s\rvert^2-2\ip rs\bigr)$,
whenever $\Cay(r)\Cay(s)\ne-1$.
\item \emph{(two charts)} $\Cay_+=\Cay$ and $\Cay_-=-\Cay$ restricted to the closed unit ball cover $M$, with images
the hemispheres $\re q\ge0$ and $\re q\le0$; on the overlap $r_-=-r_+/\lvert r_+\rvert^2$.  On each chart
\[
\lvert r-s\rvert\le d(\Cay_\pm(r),\Cay_\pm(s))\le2\,\lvert r-s\rvert .
\]
\end{enumerate}
\end{proposition}
\begin{proof}
(i) For $r\in\im\HH$, $(1-r)(1+r)=1+\lvert r\rvert^2$, which gives the second expression and the inverse by direct
substitution.  (ii) Both maps in (i) are rational with rational coefficients.  (iii) With $A=(1+r)^2=(1-\lvert r\rvert^2)+2r$ and
$B=(1+s)^2$, the product $\Cay(r)\Cay(s)$ is $AB/D$ with $D=(1+\lvert r\rvert^2)(1+\lvert s\rvert^2)$;
expanding, $\re AB=(1-\lvert r\rvert^2)(1-\lvert s\rvert^2)-4\ip rs$ and
$\im AB=2\bigl((1-\lvert s\rvert^2)r+(1-\lvert r\rvert^2)s+2r\times s\bigr)$, and
$D+\re AB=2(1+\lvert r\rvert^2\lvert s\rvert^2-2\ip rs)$; apply the inverse in (i).  (iv) $-\Cay(r)=\Cay(-r/\lvert r\rvert^2)$ by (i).  The pull-back of the
round metric is $4\lvert dr\rvert^2/(1+\lvert r\rvert^2)^2$, with conformal factor in $[1,2]$ on the unit ball.
The segment from $r$ to $s$ stays in the ball, which gives the upper bound; a minimal geodesic between two
points of a closed hemisphere can be chosen inside it, and its preimage has Euclidean length at least
$\lvert r-s\rvert$, which gives the lower bound.
\end{proof}

\begin{theorem}[exact rational arc calculus]\label{thm:rational}
Let all vertices of a sampled curve or vertex map lie in $M(\Q)$.
\begin{enumerate}[label=(\roman*)]
\item All heap operations, comparisons $\kappa$, sagittas $v$, the recurrence of Theorem~\ref{thm:recur} and its
Chebyshev coefficients are rational.
\item The Levi-Civita transport along every minimal arc, in the right trivialisation, is
$\Ad(1+x_q\bar x_p)\in\SO(3,\Q)$, and so is the Levi-Civita holonomy of every closed polygon of minimal arcs.
The spin holonomy is $N/\sqrt{\lvert N\rvert^2}$ with $N\in\Q^4$.
\item For every triangle, $\cos E\in\Q$ is given by Theorem~\ref{thm:triangle}.
\end{enumerate}
In general the half-displacements and the excess $E$ itself are irrational: for $x=1$, $y=i$, $z=j$ one has
$h_{yx}=(1+i)/\sqrt2$ and $E=\pi/2$.
\end{theorem}
\begin{proof}
(i) is clear from the formulas.  (ii) $h_{qp}=N_{qp}/\lvert N_{qp}\rvert$ with $N_{qp}=1+x_q\bar x_p\in\Q^4$, and
$\Ad(N)\xi=N\xi\bar N/\lvert N\rvert^2$ is rational; products of rational rotations are rational, and the spin
holonomy is $\prod N_{qp}/\prod\lvert N_{qp}\rvert$.  (iii) is the second formula of Theorem~\ref{thm:triangle}.
\end{proof}

Theorem~\ref{thm:rational} is the exact counterpart of the floating-point statements of Part~\ref{part:arc}: on rational data
the arc calculus, including curvature, can be carried out in exact arithmetic, and the only irrational
quantities are the square roots that normalise half-displacements, which cancel under $\Ad$.

\begin{proposition}[rational difference quotient]\label{prop:ratdq}
For $x\in C^3((a,b),M)$ put $D_hx(t):=\tfrac2h\,\Cay^{-1}\bigl(x(t+h)\bar x(t)\bigr)$ and
$D_h^cx(t):=\tfrac1h\,\Cay^{-1}\bigl(x(t+h)\bar x(t-h)\bigr)$.  Then $D_hx=\lambda_x+O(h)$ and
$D^c_hx=\lambda_x+O(h^2)$, and conversely $x(t+h)=\Cay\bigl(\tfrac h2D_hx(t)\bigr)\,x(t)$ exactly.  For rational samples
and $h\in\Q$, $D_hx(t)$ and $D_h^cx(t)$ are rational vectors.
\end{proposition}
\begin{proof}
$G(h):=x(t+h)\bar x(t-h)$ satisfies $G(-h)=G(h)^{-1}$, so $\log G$ is odd in $h$, $\log G=2h\lambda_x(t)+O(h^3)$.
For $q=\exp\xi$, $\Cay^{-1}(q)=\tan(\lvert\xi\rvert/2)\,\xi/\lvert\xi\rvert=\xi/2+O(\lvert\xi\rvert^3)$.  The forward case is
the same with $\log(x(t+h)\bar x(t))=h\lambda_x(t)+O(h^2)$.  The exact inversion is Proposition~\ref{prop:cayley}(i).
\end{proof}

With $E=\Cay(\cdot/2)$ in Theorem~\ref{thm:ftc}(i), discrete integration of a rational logarithmic derivative is
exact over $\Q$; this is the Cayley-transform method of \cite{DLP,Iserles01}, whose orders (one for left tags,
two for midpoint tags) we confirmed numerically.

% ---- insertion point: sections of the arc calculus from modules m4 (phase synthesis, rational atlas) and m5 (arc body, tropical limit)
% ---- insertion point: Part "Arithmetic and holonomy" (module m6)
\paperpart{part:final}{Context, formalisation and open problems}

\section{Relation to known work}\label{sec:known}

\begin{center}
\small
\begin{tabular}{@{}p{0.47\textwidth}p{0.47\textwidth}@{}}
\toprule
Statement & Status\\
\midrule
Prop.~\ref{m1:algebra}, Prop.~\ref{m1:chebyshev}, Thm.~\ref{m1:trichotomy} & Elementary; Chebyshev identities classical \cite{m1:DLMF}; the scope statements (Examples~\ref{m1:algebra-boundaries}, \ref{m1:power-boundaries}) not found in the searched scope\\
Thm.~\ref{m1:chordal}, Prop.~\ref{m1:sharp} & Simplicial vertices: Dirac \cite{m1:Dirac}; exact cycle consistency classical \cite{m1:Zaslavsky,m1:LermanShi,m1:Singer}; the $\varepsilon$-bound and its sharpness not found in the searched scope\\
Thm.~\ref{m1:gauge}, Prop.~\ref{m1:clique} & Classical in substance (edge-path group \cite{m1:Spanier}, spanning-tree coordinates \cite{Levy10}); the filling estimate elementary\\
Props.~\ref{m1:gram}, \ref{m1:distance}, \ref{m1:near-gram} & Classical ingredients \cite{m1:Schoenberg,m1:JordanVonNeumann,m1:Grone}; Prop.~\ref{m1:near-gram} elementary\\
Thm.~\ref{thm:main}, classification by $(\Omega1)$--$(\Omega3)$ & New as a theorem; ingredients classical \cite{Certaine,Gleason52,MZ,Milnor,Knapp}; see Rem.~\ref{rem:dalembert}\\
Thm.~\ref{thm:noreg}, Cor.~\ref{cor:interference} & Group-level statement not found in the searched scope; Lemmas~\ref{lem:noline}, \ref{lem:circle} not claimed new (possible precursors \cite{Kuczma64,Sharkovskii73}, texts not obtained); Lemma~\ref{lem:elliptic} standard; Cor.~\ref{cor:interference} follows from \cite{Davison08} once compactness is known\\
Thm.~\ref{thm:parsimony} & Proof uses \cite{PW,Folland}\\
Thm.~\ref{thm:cc} & Two-point homogeneity; cf.\ \cite{Wang52,Tits55}, here without a metric given in advance\\
Thm.~\ref{thm:operational} & Reduction to d'Alembert via centrality not found in the searched scope; cf.\ \cite{Stetkaer,Stetkaer98,CENSZ95,Davison08}; the companion paper \cite{Companion} gives a second proof of its two-readout form\\
Props.~\ref{prop:linear}--\ref{prop:conn}, Thm.~\ref{thm:obstruction}, Prop.~\ref{prop:moufang} & Classical \cite{CS,Loos,CE,Adams}\\
Prop.~\ref{prop:moufangclass}, Lemma~\ref{lem:K3} & Uses (L) \cite{Malcev55}, (K1) \cite{Grishkov80}, (K2) proved from \cite{Sagle62,Kuzmin68,Filippov76,Loginov01}; (ii) uses \cite{Kerdman}; analyticity and di-associativity are assumed\\
Prop.~\ref{prop:mfauto} & Consequence of \cite[Thm.~1]{Kuzmin71} and Mal'cev's canonical coordinates \cite{Malcev55} (via (L)) with the identity theorem; not new\\
Prop.~\ref{prop:LCarc} & Half-angle rotor, classical; identification with Levi-Civita transport and $-ds_m$ elementary (cf.\ \cite{Helgason})\\
Thm.~\ref{thm:recur} & Classical in substance \cite{Loos,MH}\\
\bottomrule
\end{tabular}
\end{center}

\begin{center}
\small
\begin{tabular}{@{}p{0.47\textwidth}p{0.47\textwidth}@{}}
\toprule
Statement (continued) & Status\\
\midrule
Thm.~\ref{thm:action}, inequality & Elementary; not found in the searched scope\\
Thm.~\ref{thm:action}, Euler--Lagrange & Discrete variational principle \cite{Veselov88,MV,MW}\\
Thm.~\ref{thm:consist} & Elementary; not found in the searched scope\\
Prop.~\ref{prop:split}, Thm.~\ref{thm:sagitta}(a),(b) & Elementary reformulation of \cite{Follmer}\\
Thm.~\ref{thm:sagitta}(c),(d), Prop.~\ref{prop:clt} & Classical \cite{Stroock,Hsu,RY,JP12}; (d) by L\'evy's argument\\
Thm.~\ref{thm:ftc} & Classical \cite{Volterra,DF,Slavik,CG,IMNZ,Iserles01,DLP,Sharpe}\\
Thm.~\ref{thm:stokes}, Prop.~\ref{prop:grid}, Cor.~\ref{cor:poincare}, Prop.~\ref{prop:plaquette} & Classical: exact finite form of the lasso decomposition \cite{Arefeva,Bralic80,Levy03,Levy10,KMR}\\
Thm.~\ref{thm:triangle} & Classical: algebraic form of the Bargmann--Pancharatnam invariant on $\operatorname{span}(x,y,z)$ \cite{SB88,MS93}, angle formulas \cite{Todhunter,VOS,Eriksson}; new only the reading as product of Levi-Civita half-displacements\\
Cor.~\ref{cor:excess} & Direct consequence of the two preceding rows; (ii) is discrete Gauss--Bonnet\\
Lemma~\ref{lem:glue}, Thm.~\ref{thm:continuum} & Classical \cite{MT,Thomassen}; Magnus expansion \cite{Magnus}\\
Cor.~\ref{cor:area}, Cor.~\ref{cor:convex} & Classical type (holonomy $\le$ curvature $\times$ area); sharp constant\\
Lemma~\ref{lem:comb} & Continuum non-abelian Stokes theorem in comb order \cite{Arefeva,KMR}; proof by variation of constants\\
Prop.~\ref{prop:converse} & Not found in the searched scope; elementary given Lemma~\ref{lem:comb}\\
Prop.~\ref{prop:weaknot} & (A) is the Mal'cev term of the Fano Steiner quasigroup \cite{Quackenbush76}; (B) not found in the searched scope\\
Lemma~\ref{lem:hurwitz} & Classical \cite[Ch.~1]{SV00}\\
Props.~\ref{prop:onesided}, \ref{prop:onesidedpos} & Not found in the searched scope; related left-isometric operations on round spheres with discontinuous right translations in \cite{Kinyon00} (Remark~\ref{rem:onesidedlit}); uses \cite{HK15,Adams}\\
Thm.~\ref{thm:round}, Cor.~\ref{cor:round} & Not found in the searched scope; uses \cite{HK15}; nearest precursor \cite[Thm.~11]{BN09}, \cite{DN68}; Lemma~\ref{lem:roundcurv} could be replaced by the Wiedersehen theorem \cite{Green63,Besse78,Weinstein74,Yang80,McKay16}; see the paragraph on roundness below\\
Lemma~\ref{lem:normal}, Prop.~\ref{prop:dim}, Thm.~\ref{thm:rigid}, Props.~\ref{prop:S7}, \ref{prop:C} & Elementary consequences of Hurwitz's theorem \cite{Hurwitz1898,SV00,CS03}, with \cite{Baez02,Schafer66}; $(\mathrm C_o)$ for all $o$ is the property of ternary Moufang loops in \cite[Lemma~3.4, Remark~(3)]{BK12}; the statements themselves not found in the searched scope (including Sabinin's and Kikkawa's ternary operations of loops)\\
Thm.~\ref{thm:disk} & Classical: Gauss--Bonnet holonomy lifted to the spin connection \cite{doCarmo76,SB88}\\
Prop.~\ref{prop:cayley}, Prop.~\ref{prop:ratdq} & Classical \cite{Cayley,Rodrigues,Shuster,SJ,DLP,Iserles01}\\
Thm.~\ref{thm:rational} & Elementary; not found in the searched scope\\
\bottomrule
\end{tabular}
\end{center}

\emph{Heaps and symmetric spaces.}  Heaps were introduced by Pr\"ufer and Baer \cite{Prufer,Baer}; the
identities of $(\Omega1)$ and Lemma~\ref{lem:certaine} are Certaine's \cite{Certaine}.  The reflection
$s_x(w)=\hp xwx$ makes every group a symmetric space in Loos's sense \cite{Loos}, and the generation of
geodesics by iterated point reflections, Theorem~\ref{thm:recur}, is a basic fact of that theory; its Chebyshev
form on the sphere is the three-term recurrence of \cite{MH}.

\emph{Functional equations: placing Theorem~\ref{thm:noreg}, Corollary~\ref{cor:interference} and
Theorem~\ref{thm:operational}.}  The classical theory on groups concerns the two-variable equations of d'Alembert,
$f(xy)+f(xy^{-1})=2f(x)f(y)$, and Wilson, $f(xy)+f(xy^{-1})=2f(x)g(y)$ \cite{Stetkaer}.  Davison determined the continuous
d'Alembert functions on compact groups: they come from continuous homomorphisms into $\SU(2)$, and a compact group
carrying a \emph{basic} solution ($f(xy)=f(y)$ for all $y$ only if $x=e$; our faithfulness implies this) is isomorphic to a
compact subgroup of $\SU(2)$ \cite{Davison08}; see also \cite{Yang13}, and \cite{Davison09} for monoids.  (1) \emph{Corollary~\ref{cor:interference}} is the
d'Alembert case.  On a \emph{compact} group with a faithful real solution bounded by $1$ it follows quickly from
\cite{Davison08}; what is added here is that compactness need not be assumed for a connected Lie group, since
faithfulness and continuity force every one-parameter subgroup to close (Lemma~\ref{lem:noline}) and then
Lemma~\ref{lem:elliptic} applies.  (2) \emph{Theorem~\ref{thm:noreg}} uses only the one-variable duplication equation
$\psi(g^2)=2\psi(g)^2-1$, the diagonal of d'Alembert's equation, which on $\R$ has many continuous solutions; we found no
earlier classification of its faithful continuous solutions on groups.  Characterisations of the cosine on the line go back at least to
\cite{Kuczma64,Sharkovskii73}; we could not obtain their texts (zbMATH carries no review of either, and the publisher and
EuDML records give metadata only), so we cannot say whether they treat the duplication equation.  We therefore claim
nothing for the one-dimensional statements (Lemmas~\ref{lem:noline}, \ref{lem:circle} and the core F16 of
Section~\ref{sec:ready}), and claim only the classification on groups.  Its ingredients are elementary
(Lemmas~\ref{lem:noline}, \ref{lem:circle}) or standard Lie theory (Lemma~\ref{lem:elliptic}: a Lie algebra with elliptic
adjoint action and one-dimensional abelian subalgebras is $\R$ or $\su(2)$).  (3) \emph{Theorem~\ref{thm:operational}}
reduces an affine two-readout law to Wilson's equation with an additional constant term and then, because $\psi$ is
central, to d'Alembert's equation; Wilson's equation is studied in \cite{Stetkaer,EbanksStetkaer15}, and the equation
$f(xy)+f(xy^{-1})=f(x)g(y)+h(y)$ with a constant-type term was solved in \cite[Lemma~1]{CENSZ95} under the extra condition
$f(xyz)=f(xzy)$, which reduces it to an abelian quotient; our $\psi$ is central but need not satisfy that condition, and we
found no statement equivalent to $(\Omega3')$.  The search covered functional equations on groups
(duplication and cosine equations, d'Alembert, Wilson, Kannappan's condition, pre-d'Alembert functions, d'Alembert
functions on compact groups) and characterisations of $S^1$, $S^3$ among groups and homogeneous spaces; it was a
search of indexed sources and reviews, not a complete survey, and earlier occurrences may exist.  That the only spheres
carrying a continuous group structure are $S^0,S^1,S^3$ is classical \cite{Samelson40} (Theorem~\ref{thm:obstruction}), and the
compact connected Lie groups of rank one are $S^1$, $\SU(2)$, $\SO(3)$ \cite{Samelson40,Knapp}; the results above characterise $S^1$ and $S^3$
among heaps on connected manifolds by one continuous scalar comparison, without assuming a metric or smoothness.
Completeness in Theorem~\ref{thm:cc} is two-point homogeneity of $(M,\kappa)$; compact connected two-point homogeneous
metric spaces were classified by Wang \cite{Wang52} and Tits \cite{Tits55}, which gives another route to the conclusion of
Theorem~\ref{thm:cc} once $\kappa$ is known to be a function of a metric; among Lie groups this leaves $S^1$, $S^3$, $\SO(3)$.
Stetk\ae r \cite{Stetkaer98} solved a quadratic-trigonometric equation on abelian groups with involution by reduction to
d'Alembert's, Wilson's and the quadratic equation; Theorem~\ref{thm:operational} works on non-abelian groups and uses that
$\psi$ is central.

\emph{Roundness: placing Theorem~\ref{thm:round}.}  After Lemma~\ref{lem:roundgeo}, $(M,g)$ is a Wiedersehen manifold: all geodesics from any
point $w$ meet again at one point $w^\perp$ at the same distance $a$.  Such a manifold is isometric to a round sphere: Green for surfaces
\cite{Green63}, and in general by the work of Berger, Kazdan, Weinstein and Yang \cite{Besse78,Weinstein74,Yang80}; we know the statement
from the survey \cite{McKay16} and have not read the proofs.  This deep theorem could replace Lemma~\ref{lem:roundcurv}; in our homogeneous
situation the elementary index estimate suffices, and it is the estimate of Green's argument (\cite[\S13]{McKay16}).  Lemma~\ref{lem:roundgeo}(i)
says that every geodesic is an integral curve of a Killing field of constant length; by \cite[Thms.~7, 10]{BN09} $(M,g)$ is then restrictively
Clifford--Wolf homogeneous and locally symmetric.  More precisely, $\lvert X_v\rvert\equiv\lvert v\rvert$ polarises to $\langle X_u,X_v\rangle\equiv\langle u,v\rangle$,
so the fields $X_{e_i}$ for an orthonormal basis $(e_i)$ of $T_wM$ form a global orthonormal frame of Killing fields: $(M,g)$ has the
Killing property of D'Atri and Nickerson \cite{DN68}, and so does its universal cover.  Berestovski\u{\i} and Nikonorov proved that a
simply connected complete Riemannian manifold has the Killing property if and only if it is a product of a Euclidean space, compact simply
connected simple Lie groups with bi-invariant metrics, and round spheres $S^7$ \cite[Thm.~11]{BN09}.  This is the nearest geometric
precursor of Theorem~\ref{thm:round}: the step from a global Killing frame to the dimension restriction is close to that theorem, which
could replace Lemma~\ref{lem:roundcurv} and Hurwitz's theorem once the product factors are sorted out (we have not written that argument).
What is new in Theorem~\ref{thm:round} is obtaining a Riemannian structure from a merely continuous comparison, through Hofmann--Kramer
\cite{HK15}, and selecting the round factor through the duplication lemmas.  Two-point homogeneity \cite{Wang52,Tits55} is not available before
the conclusion.  The use of Hofmann--Kramer replaces Szenthe's theorem, whose proof had a gap \cite[\S1]{HK15}.  The search covered the
sources named here and was not complete.

\emph{Calculus on Lie groups.}  Replacing differences by quotients in a group is the content of
Volterra's product integral \cite{Volterra,DF,Slavik} and of Lie-group integrators \cite{CG,MK,IMNZ}; the Cayley transform as a
rational substitute for the exponential is studied in \cite{DLP,Iserles01}, and in attitude kinematics as the
modified Rodrigues parameters \cite{Rodrigues,Shuster,SJ}.  Discrete variational principles for geodesics on
spheres go back to \cite{Veselov88,MV}; see \cite{MW}.

\emph{Quadratic variation.}  Pathwise quadratic variation and the pathwise It\^o integral are due to F\"ollmer
\cite{Follmer}.  The representation of spherical Brownian motion with drift $-\tfrac{n-1}2x\,dt$ is in
\cite{Stroock,Hsu}; geodesic random walks are treated in \cite{Jorgensen}, and Wong--Zakai approximations in
\cite{WZ,IW}.  For paths of lower regularity the rough-path theory of \cite{Lyons} is the relevant framework; we
do not use it.

\emph{Stokes theorems and the triangle invariant.}  The non-abelian Stokes theorem in continuum and lattice form is
due to Aref'eva \cite{Arefeva}; the lattice computation of loop averages via decompositions into lassos is Bralić's
\cite{Bralic80}, and the decomposition of loops in graphs on surfaces into lassos is systematic in L\'evy's work
\cite{Levy03,Levy10}; a product-integral treatment is in \cite{KMR}, and Gross's lasso variables are used by Driver
\cite{Driver89}.  Theorem~\ref{thm:stokes} is the exact finite form of this decomposition along a shelling; the heap
formulation adds only the explicit order and the identification of the cochain with half-displacements.  The triangle
holonomy has the algebraic form of the Bargmann--Pancharatnam invariant \cite{SB88,MS93}, evaluated on the vertices as unit vectors of
their span (not on their Hopf images).  The angle formulas of
Theorem~\ref{thm:triangle} are classical spherical trigonometry \cite{Todhunter} and the solid-angle formula
\cite{VOS,Eriksson}.

\section{Formalisation status}\label{sec:lean}

Three units of Part~\ref{part:axiom} are machine-checked in the Lean~4 library \cite{Tru}; they were merged into branch \texttt{dev} after the
commit cited in the bibliography, through the pull requests named below (statements read at the merge commits).
\begin{enumerate}[label=(\alph*),leftmargin=2em]\sloppy
\item \emph{Rank-one classification} (pull request 31765): \texttt{D5.\allowbreak S3.\allowbreak Quantum.\allowbreak Algebra.\allowbreak Comparison.\allowbreak RoundTripClassification.\allowbreak comparison\_\allowbreak classification}.
For a Lie algebra with exponential map, an invariant comparison and its Hessian form $h$ as in Lemmas~\ref{lem:class}--\ref{lem:rankone}, it
proves that $h$ is positive definite and invariant, $\varphi(\exp tX)=\cos^2(t\sqrt{h(X,X)/2})$, abelian subalgebras are lines, the
dimension is $0$, $1$ or $3$, and in dimension $3$ there is an $h$-orthonormal basis with the $\su(2)$ brackets: the $C^2$ route of
Section~\ref{sec:class} at the Lie-algebra level.
\item \emph{Rigidity of continuous cosine duplication} (pull request 31864): \texttt{D5.\allowbreak S3.\allowbreak Analytic.\allowbreak Dilation.\allowbreak DuplicationCircleRigidity},
declarations \texttt{no\_\allowbreak faithful\_\allowbreak duplication\_\allowbreak on\_\allowbreak line} and \texttt{eq\_\allowbreak scaled\_\allowbreak cos\_\allowbreak of\_\allowbreak discrete\_\allowbreak fibre}, the items F8 and F16 of
Section~\ref{sec:ready} (Lemmas~\ref{lem:noline}, \ref{lem:circle}).
\item \emph{Two-readout reduction} (pull request 31928): \texttt{D5.\allowbreak S3.\allowbreak Quantum.\allowbreak Algebra.\allowbreak Comparison.\allowbreak TwoReadoutNormalization.\allowbreak bounded\_\allowbreak dalembert\_\allowbreak of\_\allowbreak affine\_\allowbreak law},
the item F7 of Section~\ref{sec:ready} (the group-theoretic core of Theorem~\ref{thm:operational}).
\end{enumerate}
No \texttt{sorry} occurs in the three files.  These three units are all that is machine-checked.  Not formalised are: the topological
steps of Part~\ref{part:axiom} (Lemma~\ref{lem:lie}, the Gleason--Montgomery--Zippin theorem, the regularity-free argument as a whole), the remaining results
of Sections~\ref{sec:class}--\ref{sec:geom}, including the roundness theorem (Section~\ref{sec:round}) and the one-sided results
(Section~\ref{sec:onesided}), all of Part~\ref{part:arc}, and Part~\ref{part:alg} as written here (for the library declarations
that cover parts of it, see Section~\ref{sec:leanlib}).  A preregistered formalisation of the classification core (Sections~\ref{sec:def}--\ref{sec:class}) is tracked at
\begin{center}\small\url{https://github.com/the-omega-institute/trureturing/issues/31739}.\end{center}
Earlier, the library already contained, in the file
\begin{center}\small\texttt{D5/S3/Quantum/Algebra/RankOneInvariantMetric.lean} (line 444),\end{center}
the declaration
\begin{center}\small\texttt{D5.S3.Quantum.Algebra.RankOneInvariantMetric.finrank\_eq\_one\_or\_three}:\end{center}
a finite-dimensional real Lie algebra with an invariant inner product and a nonzero element whose centraliser is the line it spans has
dimension $1$ or $3$.  It covers the Lie-algebra half of Lemma~\ref{lem:rankone}: if every abelian subalgebra is a
line, the centraliser of a nonzero element $a$ is $\R a$ (a larger centraliser would contain an abelian plane), so the
declaration gives $\dim\g\in\{1,3\}$, and with the invariant inner product $\g\cong\R$ or $\su(2)$.  Name and location
were verified at commit \texttt{153b52e} of branch \texttt{dev}.  The Gleason--Montgomery--Zippin theorem is used as a
black box; a formal statement of Theorem~\ref{thm:main} can avoid it by assuming a Lie group from the start.

All identities and limits of Part~\ref{part:arc} that are not textbook facts, and the identities of Part~\ref{part:axiom} used there, were
checked numerically on random or explicit inputs by the accompanying script \texttt{check\_arc\_calculus.py} (numpy and the
standard library; 157 checks, eight of them for Part~\ref{part:alg}), which also runs the module scripts registered in it from the subdirectory \texttt{checks/} (at present six, one of them for Part~\ref{part:alg}); everything runs
from the directory of this paper.  These include the recurrence and its Chebyshev form, the residual expansion of
Theorem~\ref{thm:consist}, the inequality of Theorem~\ref{thm:action}, the split and the Brownian limits of
Theorem~\ref{thm:sagitta} (simulated, $\tfrac32T$ to two decimals), the transport formula of
Proposition~\ref{prop:LCarc} against the explicit transport on the round sphere, all formulas of
Theorem~\ref{thm:triangle} against excesses computed from vertex angles, the shelling identity of
Theorem~\ref{thm:stokes} (with the arc condition of Example~\ref{ex:shellings}) and the order of
Proposition~\ref{prop:grid} for random cochains, Lemma~\ref{lem:sign} and Corollary~\ref{cor:excess}, the Cayley formulas
and chart bounds, Theorem~\ref{thm:rational} in exact rational arithmetic, and for Section~\ref{sec:operational} the
coefficients of Lemma~\ref{lem:opmodel}, the examples of Remark~\ref{rem:operational}, and the identities of Steps~3
and~4 of the proof of Theorem~\ref{thm:operational}, the non-faithful non-cosine solution on the line and the torus
obstruction of Section~\ref{sec:regularity}, the variance bound of Theorem~\ref{thm:sagitta}(d) along dyadic partitions, and
the $O(n^{-2})$ rate of Theorem~\ref{thm:continuum} with the bound of Corollary~\ref{cor:area}, the bound
$\varrho\le\sum\min(E_T,2\pi-E_T)$, the sharpness examples after Theorem~\ref{thm:action}, Corollary~\ref{cor:convex} for a
non-convex star-shaped region and for one of area $>2\pi$, and Proposition~\ref{prop:clt}: the moment formulas against a
simulation, and the limit variance $\tfrac32T^2$ by Monte Carlo (the estimate $1.50$ from 400 paths agrees with $1.5$ within
its standard error of about $0.1$; this is not an exact agreement), and for Section~\ref{sec:geom} the examples of
Proposition~\ref{prop:weaknot} with the explicit witnesses, Lemmas~\ref{lem:hurwitz} and \ref{lem:normal}, the family of
Proposition~\ref{prop:S7}(iii), the identity $(\mathrm C_o)$ and its failure for that family and for $T_E$, and the two normal forms of the
earlier proof of Theorem~\ref{thm:rigid} (1800 assertions), and Lemma~\ref{lem:comb} with Proposition~\ref{prop:converse} on a
Clifford-torus patch (area $0.5$, spin angle $0.4807$, a strict gap) and on a great-sphere patch (equality); for
Theorem~\ref{thm:round} the model identities of Lemmas~\ref{lem:roundgeo} and \ref{lem:roundcurv} on $S^3$ and the witnesses of
Remark~\ref{rem:roundaudit} (scripts \texttt{checks/roundness\_models.py}, \texttt{checks/roundness\_witnesses.py},
\texttt{checks/roundness\_extra.py}); for Section~\ref{sec:onesided} the counterexample of Proposition~\ref{prop:onesided} (with the
contraction bound $\pi2^{-N}$), its $\mathrm U(2)$-invariance and the limits of Remark~\ref{rem:onesidedreg} (scripts
\texttt{checks/onesided\_models.py}, \texttt{checks/onesided\_witnesses.py}).  These are sanity checks, not proofs.

\subsection{Statements ready for formalisation}\label{sec:ready}
The following are finite-dimensional algebra or elementary real analysis.  Lean~4/Mathlib names: $\HH=$
\texttt{Quaternion Real} (notation $\HH[\R]$), $\bar x=$ \texttt{star x}, $\ip xy=$ \texttt{inner x y} $=(x\bar y)$\texttt{.re}
(\texttt{Quaternion.inner\_def}), $\lVert x\rVert^2=$ \texttt{normSq x}, $\re x=$ \texttt{x.re}, $x_{\im}=$ \texttt{x.im} (a
quaternion with zero real part), unit points = \texttt{Metric.sphere (0 : }$\HH[\R]$\texttt{) 1}, $\cos=$ \texttt{Real.cos},
$U_n=$ \texttt{Polynomial.Chebyshev.U Real n} with $n:\Z$ and $U_{-1}=0$; the cross product $p\times q$ of pure quaternions is the
imaginary part of $pq$ (as $pq=-\ip pq+p\times q$), so no separate cross product is needed.  We write
$\kappa(a,b):=(1+\ip ab)/2$.
\begin{enumerate}[label=\textbf{F\arabic*.},leftmargin=2.4em]\small
\item (reflection) $\forall\,x\,w:\HH[\R]$, \ $x\bar wx=2\ip xw\,x-\lVert x\rVert^2w$.  No hypotheses.
\item (interference, round trip) $\forall\,x\,w\,y$, $\lVert x\rVert=1\to\ip y{x\bar wx}+\ip yw=2\ip wx\ip yx$; and
$\lVert x\rVert=1\to\lVert w\rVert=1\to\kappa(w,x\bar wx)=(2\kappa(w,x)-1)^2$.
\item (exact split) $\lVert x\rVert=1\to\lVert y\rVert=1\to y-x=(y\bar x)_{\im}\,x-(1-\ip xy)\,x$ and
$\ip x{(y\bar x)_{\im}\,x}=0$.
\item (Chebyshev) $x:\mathbb N\to\HH[\R]$, $\lVert x_0\rVert=\lVert x_1\rVert=1$, $\forall n$, $x_{n+2}=x_{n+1}\bar x_nx_{n+1}$; with
$c=\ip{x_0}{x_1}$: $\forall n:\mathbb N$, $x_{n+1}=U_n(c)\,x_1-U_{n-1}(c)\,x_0$, and $\lVert x_n\rVert=1$.
\item (discrete least action) $N\ge1$, $t:\mathrm{Fin}\,N\to\R$, $\forall i,\ t_i\ge0$, $\sum_it_i=\Theta$, $\Theta<\pi$
$\to\sum_i(1-\cos t_i)\ge N(1-\cos(\Theta/N))$, with equality iff $\forall i,\ t_i=\Theta/N$.
\item ($b=-1$) $G$ with \texttt{[Group G]}, $\psi:G\to\R$, $\forall g,\ \psi(g^{-1})=\psi(g)$, $\forall g\,h,\ \psi(hgh^{-1})=\psi(g)$,
$\forall g,\ (\psi(g)=1\leftrightarrow g=1)$; $a\,b\,d:\R$, $u:G$ with $\forall z,\ \psi(zu^{-1})=a\psi(z)+b\psi(zu)+d$, and $u*u\ne1$
$\to b=-1$.
\item ($(\Omega3')\Rightarrow$ d'Alembert) \texttt{[Group G]}, $\psi:G\to\R$, $\psi(1)=1$, $\psi(g^{-1})=\psi(g)$, $\psi(hgh^{-1})=\psi(g)$,
$\forall g,\ \lvert\psi(g)\rvert\le1$, $\exists g,\ \psi(g)\ne1$; $a\,d:G\to\R$ with $\forall z\,u$,
$\psi(zu)+\psi(zu^{-1})=a(u)\psi(z)+d(u)$ $\to\exists\,m<1$ such that $\psi_0:=(\psi-m)/(1-m)$ satisfies
$\forall g,\ \lvert\psi_0(g)\rvert\le1$ and $\forall z\,u$, $\psi_0(zu)+\psi_0(zu^{-1})=2\psi_0(z)\psi_0(u)$.  (Pure algebra and
boundedness; no topology.)
\item (no faithful solution on a line) $\neg\exists\,p:\R\to\R$, \texttt{Continuous p} $\wedge\ p(0)=1\wedge\forall t,\ \lvert p(t)\rvert\le1
\wedge\forall t,\ p(2t)=2p(t)^2-1\wedge\forall t>0,\ p(t)<1$.
\item (circle rigidity) $\psi:\R\to\R$, \texttt{Continuous}, \texttt{Function.Periodic }$\psi\ (2\pi)$, $\forall\theta,\ \psi(-\theta)=\psi(\theta)$,
$\lvert\psi\rvert\le1$, $\forall\theta,\ \psi(2\theta)=2\psi(\theta)^2-1$, $\forall\theta,\ (\psi(\theta)=1\leftrightarrow\exists k:\Z,\ \theta=2\pi k)$
$\to\psi=$ \texttt{Real.cos}.
\item (triangle closed form) $\lVert x\rVert=\lVert y\rVert=\lVert z\rVert=1\to(1+x\bar z)(1+z\bar y)(1+y\bar x)=
2(1+\ip xy+\ip yz+\ip zx)+2\bigl((y\bar x)_{\im}(z\bar x)_{\im}\bigr)_{\im}$, and $\lVert1+y\bar x\rVert^2=4\kappa(x,y)$.
\item (Gram identity) $\forall\,a\,b\,c:\R$, $(1+a+b+c)^2+(1+2abc-a^2-b^2-c^2)=2(1+a)(1+b)(1+c)$ (\texttt{ring}).
\item (orientation, Lemma~\ref{lem:sign}) unit $x,y,z,\nu$ with $\ip\nu x=\ip\nu y=\ip\nu z=0$ $\to(\nu\bar x)$\texttt{.re}$=0$ and
$\bigl((y\bar x)_{\im}(z\bar x)_{\im}\bigr)_{\im}=D\cdot\nu\bar x$, where $D=$ \texttt{Matrix.det} of the $4\times4$ matrix of coordinates
of $x,y,z,\nu$ (rows; coordinates by \texttt{Quaternion.equivTuple}).
\item (rationality) $x\,y:$ \texttt{Quaternion Rat}, \texttt{normSq x = 1}, \texttt{normSq y = 1}, $q=1+y\bar x\ne0$ $\to$ the matrix of
$v\mapsto qv\bar q/$\texttt{normSq q} on pure quaternions lies in \texttt{Matrix.specialOrthogonalGroup (Fin 3) Rat}.
\item (Cayley chart) for pure $r$ (\texttt{r.re = 0}), $\Cay(r):=((1-\lVert r\rVert^2)+2r)/(1+\lVert r\rVert^2)$ has norm $1$,
$\Cay(r)_{\im}/(1+\Cay(r)$\texttt{.re}$)=r$, $-\Cay(r)=\Cay(-r/\lVert r\rVert^2)$ for $r\ne0$, and the composition law of
Proposition~\ref{prop:cayley}(iii) holds whenever its denominator is nonzero.
\item (sagitta bound) $d:\mathrm{Fin}\,n\to\R$, $\forall k,\ 0\le d_k\le\pi$ $\to\sum_k(1-\cos d_k)\le\tfrac12(\max_kd_k)\sum_kd_k$.
\item (one-dimensional core of Theorem~\ref{thm:noreg}) $p:\R\to\R$, \texttt{Continuous p}, $\forall t,\ \lvert p(t)\rvert\le1$,
$p(0)=1$, $\forall t,\ p(-t)=p(t)$, $\forall t,\ p(2t)=2p(t)^2-1$; let $Z=\{t\mid p(t)=1\}$; assume $\exists\varepsilon>0,\
Z\cap(-\varepsilon,\varepsilon)=\{0\}$ and $\forall s\in Z,\forall t,\ p(t+s)=p(t)$.  Then $\exists\,T>0$ with $Z=T\Z$ and
$\forall t,\ p(t)=\cos(2\pi t/T)$.  (Proof: $Z$ is a closed subgroup by periodicity and evenness, discrete by assumption,
nonzero by F8; then F9 after rescaling.)
\item (Fano example, Proposition~\ref{prop:weaknot}(A)) on \texttt{Fin 7} encoded as the nonzero elements of
$(\Z/2)^3$ with the operation of the proposition: the Mal'cev identities, the heap identities on every line, the failure of
the third identity, and the non-di-associativity of $\hp xoy$ are finite statements decidable by \texttt{decide}.
\end{enumerate}
The Hurwitz form of Theorem~\ref{thm:rigid} (a bilinear, norm-multiplicative product on $\R^4$ with two-sided unit $1$ is
$uv$ or $vu$) is not on this list: the Mathlib version we inspected (the one pinned by \cite{Tru}) contains neither Hurwitz's
theorem nor composition algebras, so a formal proof would have to supply that classification.

\section*{Where these quantities occur in physics}
This section consists of pointers only; it adds no claim to the paper.  The companion paper \cite[Physical
examples]{Companion} collects physical instances.  Items (i), (ii) and (iv) below are exact instances of statements proved here; item
(iii) is a correspondence of algebraic form.  Items (i)--(iii) point to measurements; item (iv) is theoretical.
\begin{enumerate}[label=(\roman*),leftmargin=2em]
\item For a spin rotation $U$ by the angle $\alpha\ge0$ and every unit spinor $\psi$ one has $\re\ip\psi{U\psi}=\cos\frac\alpha2$, so
$\kappa(\psi,U\psi)=\tfrac12(1+\cos\frac\alpha2)$ for all $\alpha$, with period $4\pi$; the distance is $d(\psi,U\psi)=\alpha/2$ for $0\le\alpha\le2\pi$ and
$d=2\pi-\alpha/2$ for $2\pi\le\alpha\le4\pi$.  Read through the path$\otimes$spin
composite, which is the subject of the companion paper, $\kappa$ is the fringe function of the 1975 neutron interferometry
experiments on $4\pi$ periodicity \cite{Rauch75,WCOE75}, whose most precise early measured period is $716.8\pm3.8^\circ$
\cite{Rauch78}.
\item Theorem~\ref{thm:recur} with $c=\cos(\omega\tau/2)$ is the recurrence of the spinor $\psi_n=U^n\psi_0$,
$U=e^{-i\omega\tau\sigma_z/2}$, under stroboscopic Larmor precession; its form (d) is the Cayley--Hamilton identity
$U^2=(\operatorname{tr}U)\,U-1$.  This half-angle recurrence is visible only in phase-referenced data.  Conventional single-spin NMR
and ESR record Bloch-vector components; the transverse components obey the full-angle recurrence with $\cos(\omega\tau)$, and the
longitudinal component obeys it only if it vanishes (its defect is $(2-2\cos\omega\tau)\lvert n_z\rvert$).  The spinor sign itself was observed by
NMR interferometry with a coupled second spin \cite{SVV77}.
\item The triangle holonomy of Theorem~\ref{thm:triangle} is a statement about Levi-Civita transport on $S^3$; formula \eqref{eq:tri} has the
algebraic form of the Bargmann invariant $\operatorname{Tr}(P_1P_2P_3)$ of three spin-$\tfrac12$ states, evaluated on $x,y,z$ regarded as unit vectors of
the three-dimensional space $\operatorname{span}(x,y,z)$, with $i$ replaced by $\nu\bar x$ \cite{SB88,MS93}; it is not the Bargmann invariant of the
Hopf images of $x,y,z$, and the measured geometric phase \cite{BS88} belongs to the Hopf connection treated in the companion paper.  (Our
Hopf map $x\mapsto xi\bar x$ differs from the Bloch map of the companion paper by an orientation reversal of the $2$-sphere; signed solid
angles here refer to the companion's Bloch orientation.)  For
$x,y,z$ in one great $2$-sphere $\Sigma=S^3\cap\nu^\perp$ the identification is the same for all triangles, with $\Sigma$ in the role of the Bloch
sphere, and Corollary~\ref{cor:excess}(ii) and Theorem~\ref{thm:disk} give $\cos\frac A2+\sin\frac A2\,\nu\bar x_0$ for a region of area $A$ in $\Sigma$,
the form of the half-solid-angle law.  With the conventions stated after
Theorem~\ref{thm:triangle}, $\arg\operatorname{Tr}(P_1P_2P_3)=+\Omega/2$, while the geometric phase acquired by a spin-$\tfrac12$ state taken around
the cycle $n_1\to n_2\to n_3\to n_1$ along geodesic arcs,
\[
\arg\langle\psi(0)|\psi(T)\rangle-\operatorname{Im}\int_0^T\langle\psi|\dot\psi\rangle\,dt=-\Omega/2=-\arg\operatorname{Tr}(P_1P_2P_3)
\]
(Hermitian product, antilinear in the first slot).  The polarised-neutron measurement of Bitter and Dubbers \cite{BD87} concerns the adiabatic
Berry phase in a helical field, a continuous loop; it has the form of Theorem~\ref{thm:disk} (the half-solid-angle law) when the Bloch
sphere is identified with a great $2$-sphere, and it is not a three-point invariant.
\item (Theoretical.)  Under isotropic rotational diffusion with rotation vector of covariance $\langle\theta_i\theta_j\rangle=2Dt\,\delta_{ij}$, the spinor
performs Brownian motion on $S^3$ in the sense of Theorem~\ref{thm:sagitta} at paper time $Dt/2$: a rotation by $\theta$ is a
step of length $\lvert\theta\rvert/2$, and Brownian motion with generator $\tfrac12\Delta$ has tangential covariance $t\,\delta_{ij}$.  Hence
$\mathbb E\,c(\psi_0,\psi_t)=e^{-3Dt/4}$ by Proposition~\ref{prop:clt}, the standard decay of the $j=\tfrac12$ Casimir mode under
isotropic rotational diffusion, and sagitta sums over $[0,t]$ tend to $\tfrac34Dt$.  We cite no measurement of this decay.
\end{enumerate}

\subsection{Library declarations related to Part~\ref{part:alg}}\label{sec:leanlib}
The following declarations of the library \cite{Tru} were located by name and their statements read (local clone at commit
\texttt{9d18ab6}); no \texttt{sorry} occurs in the files of the first six.  They are results of the library, not formalisations
carried out for this paper, and several are narrower than the corresponding statements of Part~\ref{part:alg}, as indicated.
{\small\raggedright
\begin{itemize}[leftmargin=1.5em]
\item Proposition~\ref{m1:algebra}, matrix slice (skew-adjoint unitary squares to $-1$):
\texttt{D5/S3/Quantum/Algebra/Unitary/BalancedUnitaryComplexStructure.lean},
\texttt{square\_eq\_neg\_one\_of\_skewAdjoint\_unitary}; nearby \texttt{relativeComplexStructure\_square}.
\item Proposition~\ref{m1:time}: \texttt{D5/S3/Analytic/Chebyshev/PositivityThreshold.lean},
\texttt{forall\_pos\_eval\_T\_iff\_one\_le} and \texttt{eval\_T\_cos\_floor\_succ\_neg\_and\_nonneg\_before}.
\item Remark~\ref{m1:spin}, generator slice: \texttt{D5/S3/Analytic/ReflectedSpectrum/EllipticHyperbolicReflectionTrichotomy.lean},
\texttt{elliptic\_hyperbolic\_reflection\_trichotomy} (its middle generator is the zero matrix).
\item Theorem~\ref{m1:chordal}: \texttt{D5/S3/Combinatorics/Graph/Chordal/GainPotential.lean},
\texttt{SimpleGraph.IsChordal.exists\_gain\_potential} (exact case) and \texttt{GainPotentialStability.lean},
\texttt{SimpleGraph.IsChordal.exists\_approximate\_gain\_potential} (the bound $(|V|-2)\varepsilon$ with truncated subtraction, for a metric
invariant on both sides); simplicial vertex: \texttt{Dirac.lean}, \texttt{SimpleGraph.IsChordal.exists\_isSimplicial}.
\item Corollary~\ref{m1:anchored}: \texttt{D5/S3/Observer/AgencyHolonomy/AnchoredPhaseClassification.lean},
\texttt{anchored\_phase\_classification} (circle-valued links).
\item Chordal positive semidefinite completion \cite{m1:Grone}: \texttt{D5/S3/QuadraticForms/ChordalPsdCompletion.lean},
\texttt{exists\_posSemidef\_completion}; rank-bounded version \texttt{RankBoundedChordalPsdCompletion.lean},
\texttt{exists\_posSemidef\_completion\_rank\_le}.
\item Propositions~\ref{m1:gram}, \ref{m1:distance}: \texttt{D5/S3/Quantum/Measurement/AntidistinguishabilityGram.lean}, \texttt{gramPsd};
\texttt{D5/S3/Weil/Probability/FiniteGaussianSchoenberg.lean}, \texttt{gram\_distance\_conditionally\_negative};
\texttt{D5/S3/Quantum/Algebra/GramUnitaryExtension.lean}, \texttt{exists\_unitary\_mul\_eq\_of\_conjTranspose\_mul\_eq};
parallelogram law: \texttt{D5/S3/QuantumContext/ParallelogramNormLaw.lean}, \texttt{inner\_product\_norm\_parallelogram\_law}.
\item Proposition~\ref{m1:readout-identities}: \texttt{D5/S3/QuantumBounds/LagrangeGramIdentity.lean}, \texttt{lagrange\_gram\_identity};
\texttt{D5/S3/Quantum/Algebra/CommonPhasePolynomial.lean}, \texttt{complete\_common\_phase\_polynomial\_theorem}.
\item Example~\ref{m1:overlap-boundaries}: \texttt{D5/S3/Quantum/PureState/UnitGramIndistinguishability.lean},
\texttt{unit\_gram\_overlap\_characterization}; \texttt{D5/S3/Quantum/Information/BargmannImaginaryMaximizers.lean},
\texttt{Latitude.latitude\_dot}; a larger contextuality statement: \texttt{D5/S3/Quantum/Contextuality/LinearDoilyNegativeLineParity.lean},
\texttt{result}.
\item Remark~\ref{m1:phase}: \texttt{D5/S3/Quantum/Information/GroupInvariantPsi/C129QubitBargmann.lean}, \texttt{triple\_bargmann\_bloch}
and \texttt{explicit\_qubit\_orientation}; \texttt{BargmannImaginaryMaximizers.lean}, \texttt{bargmann\_path\_product}.
\item Example~\ref{m1:sequential}: \texttt{D5/S3/Quantum/Measurement/StaticEffectSequentialSeparation.lean},
\texttt{same\_effects\_different\_two\_step\_joint\_law}.
\end{itemize}}
No declaration was found for the gauge quotient of Theorem~\ref{m1:gauge}, the filling estimate of Proposition~\ref{m1:clique}, the
norm trichotomy of Theorem~\ref{m1:trichotomy} or the sharpness results of Proposition~\ref{m1:sharp}.

\section{What remains open}\label{sec:open}

\begin{enumerate}[label=\textbf{Q\arabic*.},leftmargin=2.6em]
\item \emph{Regularity of topological loops.}  Is a di-associative (or Moufang) loop whose underlying space is a connected
topological manifold, with continuous multiplication and divisions, isomorphic to an analytic loop?  For Moufang loops the differentiable case is known: a review of \cite{NS98} states that every differentiable connected
Moufang loop is analytic (we have not read the proof).  For di-associative loops the answer is negative in general: Hofmann
constructed a topological loop on $\R^3$, commutative and di-associative (any two elements lie in a subgroup) but not associative, whose
operation is analytic in no coordinate system \cite[(9.1)--(9.3)]{Hofmann58II}; it is not Moufang.  The same paper proves that a compact
connected commutative di-associative loop without small cyclic subgroups is associative \cite[(8.30)]{Hofmann58II}.  So the question
stays open for Moufang loops and for compact di-associative loops; a topological loop homeomorphic to $S^7$ whose left translations
generate a compact Lie group is one of the two seven-dimensional compact Moufang loops \cite[Thm.~6]{FS15}, which on the round $S^7$
with an invariant $\kappa$ also follows from Proposition~\ref{prop:S7}.  We did not obtain \cite{HS90} or the papers of Hudson on loops
with invariant uniformities; the global theory of analytic Moufang loops is in \cite{Kerdman,NS02}.  An affirmative answer for compact
Moufang loops removes analyticity from Proposition~\ref{prop:moufangclass} once compactness of $M$ is known; in the weakened setting
compactness is not among the hypotheses.  For the classification in this paper the question is no longer needed: under two-sided
invariance Theorem~\ref{thm:round} gives compactness and roundness without analyticity (Remark~\ref{rem:roundrel}(ii)); it remains a
question about topological loops.
\item \emph{From local to global} (answered).  An earlier version asked whether a connected analytic di-associative loop whose tangent
algebra is a Malcev algebra is a Moufang loop.  It is: Proposition~\ref{prop:mfauto}, a consequence of Mal'cev's and Kuz'min's
theorems and the identity theorem; so Proposition~\ref{prop:moufangclass}(ii) no longer assumes the Moufang identity.
\item \emph{The two-generator axiom.}  Answered except for the one-sided variants below.  The axiom does not imply di-associativity
(Proposition~\ref{prop:weaknot}).
The former parts (a) and (b) are answered: with the Mal'cev identities, $(\Omega2\mathrm a$--$\mathrm d)$ with invariance under left and
right transports, and $(\Omega3)$, $(M,\kappa)$ is the round $S^1$, $S^3$ or $S^7$ with the standard comparison, the third identity holds on
two-generated subsystems, and the operations are those of Proposition~\ref{prop:S7} (Theorem~\ref{thm:round}, Corollary~\ref{cor:round});
$(\mathrm C_o)$ for all $o$ leaves only the standard structures (Proposition~\ref{prop:C}).  The question posed in an earlier version,
whether invariance under the \emph{left} transports alone still gives a round sphere, has a negative answer as posed
(Section~\ref{sec:onesided}, Proposition~\ref{prop:onesided}); what remains true is Proposition~\ref{prop:onesidedpos}.  Open: with left invariance only,
is $(M,\kappa)$ a round sphere if in addition (a) the third identity holds on two-generated subsystems, or $B$ is power-associative; or (b)
$\kappa$ satisfies the regularity clause $(\Omega2\mathrm e)$ (twice differentiable at one point of the diagonal, or smooth); or both?  The example of
Proposition~\ref{prop:onesided} fails all three additional conditions (its comparison is continuous but not $C^2$ on the diagonal,
Remark~\ref{rem:onesidedreg}); Example~(B) of Proposition~\ref{prop:weaknot} satisfies them and lives on the round $S^3$.  By Proposition~\ref{prop:onesidedpos}(i) the question is one about a continuous section $\ell$ of $G\to G/H$ and a function $f$ on the
double cosets $H\backslash G/H$, and by part (ii) it is equivalent to the four-point Gram test.
\item \emph{Equality in Corollary~\ref{cor:area}.}  Let $\gamma$ be a Jordan curve in $S^3$ with $\Theta(U_\gamma)=\inf_F\operatorname{area}(F)$ over $C^3$
fillings $F$ of the square (Remark~\ref{rem:converse}(ii); the infimum is positive).  Must $\gamma$ lie in a great $2$-sphere?  For a given filling of rank $2$ on a dense set that attains the infimum,
yes (Proposition~\ref{prop:converse}); without a rank condition, no (the tripod loop, Remark~\ref{rem:converse}(i), which is not a Jordan curve).
For smooth Jordan curves the answer is yes provided Morrey's minimiser is $C^3$ up to the boundary and immersed inside, as the
regularity theory summarised in \cite{White16} states (Remark~\ref{rem:converse}(ii)); we have not checked the original statements
\cite{Morrey48,Osserman70,Gulliver73,Hildebrandt69,HH70}, and for piecewise $C^3$ curves the corners need a separate argument.
\item \emph{Sparse comparisons.}  (a) The optimal universal coefficient in Theorem~\ref{m1:chordal} for groups with a bi-invariant metric
is $1$ for three vertices, at most $n-3$ for $n\ge4$ vertices, and $(n-2)/n$ for the real line (Proposition~\ref{m1:sharp}); what is it in
general?  (b) The continuum programme of Remark~\ref{m1:continuum-programme}: refinement, compactness modulo gauge, and identification of
normalised face holonomies with curvature.
\end{enumerate}

\section*{Use of AI tools}
The authors used AI tools, namely Claude Fable 5.2 and Claude Opus 5.5 (via Claude Code) and Codex (gpt-6-astra and gpt-6.2-sol), for literature search, computation and double-checking, including formal verification of the results in Lean. All proofs can be checked by hand, and the computations can be reproduced from the verification package. The authors take full responsibility for their contributions.

\begin{thebibliography}{99}\small
\bibitem{Adams} J.~F.~Adams, On the non-existence of elements of Hopf invariant one, \emph{Ann. of Math.} \textbf{72} (1960), 20--104.
\bibitem{Arefeva} I.~Ya.~Aref'eva, Non-Abelian Stokes formula, \emph{Theor. Math. Phys.} \textbf{43} (1980), 353--356; Russian original \emph{Teor. Mat. Fiz.} \textbf{43} (1980), 111--116.
\bibitem{Baer} R.~Baer, Zur Einf\"uhrung des Scharbegriffs, \emph{J. Reine Angew. Math.} \textbf{160} (1929), 199--207.
\bibitem{CS} \'E.~Cartan, J.~A.~Schouten, On the geometry of the group-manifold of simple and semi-simple groups, \emph{Proc. Akad. Wetensch. Amsterdam} \textbf{29} (1926), 803--815.
% UNVERIFIED (Cayley 1846: not checked against a source)
\bibitem{Cayley} A.~Cayley, Sur quelques propri\'et\'es des d\'eterminants gauches, \emph{J. Reine Angew. Math.} \textbf{32} (1846), 119--123.
\bibitem{Certaine} J.~Certaine, The ternary operation $(abc)=ab^{-1}c$ of a group, \emph{Bull. Amer. Math. Soc.} \textbf{49} (1943), 869--877.
\bibitem{CE} C.~Chevalley, S.~Eilenberg, Cohomology theory of Lie groups and Lie algebras, \emph{Trans. Amer. Math. Soc.} \textbf{63} (1948), 85--124.
\bibitem{CG} P.~E.~Crouch, R.~Grossman, Numerical integration of ordinary differential equations on manifolds, \emph{J. Nonlinear Sci.} \textbf{3} (1993), 1--33.
\bibitem{Davison08} T.~M.~K.~Davison, D'Alembert's functional equation on topological groups, \emph{Aequationes Math.} \textbf{76} (2008), 33--53.
\bibitem{Davison09} T.~M.~K.~Davison, D'Alembert's functional equation on topological monoids, \emph{Publ. Math. Debrecen} \textbf{75} (2009), 41--66.
\bibitem{DLP} F.~Diele, L.~Lopez, R.~Peluso, The Cayley transform in the numerical solution of unitary differential systems, \emph{Adv. Comput. Math.} \textbf{8} (1998), 317--334.
\bibitem{DF} J.~D.~Dollard, C.~N.~Friedman, \emph{Product Integration with Applications to Differential Equations}, Encyclopedia Math. Appl. \textbf{10}, Addison-Wesley, 1979.
\bibitem{Eriksson} F.~Eriksson, On the measure of solid angles, \emph{Math. Mag.} \textbf{63} (1990), 184--187.
\bibitem{Follmer} H.~F\"ollmer, Calcul d'It\^o sans probabilit\'es, in \emph{S\'eminaire de Probabilit\'es XV}, Lecture Notes in Math. \textbf{850}, Springer, 1981, 143--150.
\bibitem{Folland} G.~B.~Folland, \emph{A Course in Abstract Harmonic Analysis}, 2nd ed., CRC Press, 2016.
\bibitem{Grishkov80} A.~N.~Grishkov, Structure and representations of binary-Lie algebras, \emph{Izv. Akad. Nauk SSSR Ser. Mat.} \textbf{44} (1980), no.~5, 999--1030; English transl. \emph{Math. USSR-Izv.} \textbf{17} (1981), 243--269.
\bibitem{Gleason52} A.~M.~Gleason, Groups without small subgroups, \emph{Ann. of Math.} \textbf{56} (1952), 193--212.
\bibitem{Helgason} S.~Helgason, \emph{Differential Geometry, Lie Groups, and Symmetric Spaces}, Academic Press, 1978.
\bibitem{Hsu} E.~P.~Hsu, \emph{Stochastic Analysis on Manifolds}, Grad. Stud. Math. \textbf{38}, Amer. Math. Soc., 2002.
\bibitem{IW} N.~Ikeda, S.~Watanabe, \emph{Stochastic Differential Equations and Diffusion Processes}, 2nd ed., North-Holland, 1989.
\bibitem{Iserles01} A.~Iserles, On Cayley-transform methods for the discretization of Lie-group equations, \emph{Found. Comput. Math.} \textbf{1} (2001), 129--160.
\bibitem{IMNZ} A.~Iserles, H.~Z.~Munthe-Kaas, S.~P.~N{\o}rsett, A.~Zanna, Lie-group methods, \emph{Acta Numer.} \textbf{9} (2000), 215--365.
\bibitem{Jorgensen} E.~J{\o}rgensen, The central limit problem for geodesic random walks, \emph{Z. Wahrscheinlichkeitstheorie verw. Gebiete} \textbf{32} (1975), 1--64.
\bibitem{KMR} R.~L.~Karp, F.~Mansouri, J.~S.~Rno, Product integral formalism and non-Abelian Stokes theorem, \emph{J. Math. Phys.} \textbf{40} (1999), 6033--6043.
% Content (simply connected loop unique; connected ones are quotients by discrete central normal subgroups) checked only
% through a secondary source (Loginov, hep-th/0109206), which cites the long paper.
\bibitem{Kerdman} F.~S.~Kerdman, Analytic Moufang loops in the large, \emph{Algebra i Logika} \textbf{18} (1979), 523--555; English transl. \emph{Algebra and Logic} \textbf{18} (1980), 325--347.  Announced in \emph{Soviet Math. Dokl.} \textbf{20} (1979), 1297--1300.
\bibitem{HS90} K.~H.~Hofmann, K.~Strambach, Topological and analytic loops, in \emph{Quasigroups and Loops: Theory and Applications} (O.~Chein, H.~O.~Pflugfelder, J.~D.~H.~Smith, eds.), Sigma Ser. Pure Math. \textbf{8}, Heldermann, 1990, 205--262.
\bibitem{Knapp} A.~W.~Knapp, \emph{Lie Groups Beyond an Introduction}, 2nd ed., Birkh\"auser, 2002.
% Bibliographic data verified (mathnet / Springer record).  Content checked only through secondary sources: over an
% algebraically closed field of characteristic 0 the only simple non-Lie Malcev algebra is seven-dimensional.
\bibitem{Kuzmin68} E.~N.~Kuz'min, Mal'cev algebras and their representations, \emph{Algebra i Logika} \textbf{7} (1968), no.~4, 48--69; English transl. \emph{Algebra and Logic} \textbf{7} (1968), 233--244.
\bibitem{Kuzmin71} E.~N.~Kuz'min, The connection between Mal'cev algebras and analytic Moufang loops, \emph{Algebra i Logika} \textbf{10} (1971), 3--22; English transl. \emph{Algebra and Logic} \textbf{10} (1971), 1--14.
\bibitem{Loos} O.~Loos, \emph{Symmetric Spaces I: General Theory}, Benjamin, 1969.
\bibitem{Lyons} T.~J.~Lyons, Differential equations driven by rough signals, \emph{Rev. Mat. Iberoam.} \textbf{14} (1998), 215--310.
\bibitem{Magnus} W.~Magnus, On the exponential solution of differential equations for a linear operator, \emph{Comm. Pure Appl. Math.} \textbf{7} (1954), 649--673.
% Full text read in Round 15 (mathnet.ru, Mi sm5205): Thm 1 one-parameter subgroups / canonical coordinates of local analytic
% alternative loops; Thm 2 tangent algebra binary-Lie, unique local loop, Campbell-Hausdorff formula; Thm 4 Moufang -> Malcev.
% (An earlier comment here said "verified (scan)"; the Round-12 log said "not read"; the Round-15 reading supersedes both.)
\bibitem{Malcev55} A.~I.~Mal'cev, Analytic loops, \emph{Mat. Sb. (N.S.)} \textbf{36(78)} (1955), no.~3, 569--576.
\bibitem{MW} J.~E.~Marsden, M.~West, Discrete mechanics and variational integrators, \emph{Acta Numer.} \textbf{10} (2001), 357--514.
\bibitem{MH} J.~C.~Mason, D.~C.~Handscomb, \emph{Chebyshev Polynomials}, Chapman \& Hall/CRC, 2003.
\bibitem{Milnor} J.~Milnor, Curvatures of left invariant metrics on Lie groups, \emph{Adv. Math.} \textbf{21} (1976), 293--329.
% UNVERIFIED (section number of the theta-curve lemma not checked)
\bibitem{MT} B.~Mohar, C.~Thomassen, \emph{Graphs on Surfaces}, Johns Hopkins Univ. Press, 2001.
\bibitem{NS02} P.~T.~Nagy, K.~Strambach, \emph{Loops in Group Theory and Lie Theory}, de Gruyter Exp. Math. \textbf{35}, de Gruyter, 2002.
\bibitem{MZ} D.~Montgomery, L.~Zippin, \emph{Topological Transformation Groups}, Interscience, 1955.
\bibitem{MV} J.~Moser, A.~P.~Veselov, Discrete versions of some classical integrable systems and factorization of matrix polynomials, \emph{Comm. Math. Phys.} \textbf{139} (1991), 217--243.
\bibitem{MK} H.~Munthe-Kaas, High order Runge--Kutta methods on manifolds, \emph{Appl. Numer. Math.} \textbf{29} (1999), 115--127.
\bibitem{PW} F.~Peter, H.~Weyl, Die Vollst\"andigkeit der primitiven Darstellungen einer geschlossenen kontinuierlichen Gruppe, \emph{Math. Ann.} \textbf{97} (1927), 737--755.
\bibitem{Prufer} H.~Pr\"ufer, Theorie der Abelschen Gruppen I, \emph{Math. Z.} \textbf{20} (1924), 165--187.
\bibitem{RY} D.~Revuz, M.~Yor, \emph{Continuous Martingales and Brownian Motion}, 3rd ed., Springer, 1999.
\bibitem{Rodrigues} O.~Rodrigues, Des lois g\'eom\'etriques qui r\'egissent les d\'eplacements d'un syst\`eme solide dans l'espace, \emph{J. Math. Pures Appl.} \textbf{5} (1840), 380--440.
\bibitem{SJ} H.~Schaub, J.~L.~Junkins, Stereographic orientation parameters for attitude dynamics: a generalization of the Rodrigues parameters, \emph{J. Astronaut. Sci.} \textbf{44} (1996), 1--19.
\bibitem{Sharpe} R.~W.~Sharpe, \emph{Differential Geometry: Cartan's Generalization of Klein's Erlangen Program}, Grad. Texts in Math. \textbf{166}, Springer, 1997.
\bibitem{Shuster} M.~D.~Shuster, A survey of attitude representations, \emph{J. Astronaut. Sci.} \textbf{41} (1993), 439--517.
\bibitem{Slavik} A.~Slav\'ik, \emph{Product Integration, Its History and Applications}, Matfyzpress, Prague, 2007.
\bibitem{Stetkaer} H.~Stetk\ae r, \emph{Functional Equations on Groups}, World Scientific, 2013.
\bibitem{Stroock} D.~W.~Stroock, \emph{An Introduction to the Analysis of Paths on a Riemannian Manifold}, Math. Surveys Monogr. \textbf{74}, Amer. Math. Soc., 2000.
\bibitem{Thomassen} C.~Thomassen, The Jordan--Sch\"onflies theorem and the classification of surfaces, \emph{Amer. Math. Monthly} \textbf{99} (1992), 116--130.
\bibitem{Todhunter} I.~Todhunter, \emph{Spherical Trigonometry}, 5th ed., Macmillan, London, 1886, Art.~103.
\bibitem{Tru} The Omega Institute, \emph{trureturing}, Lean~4 library, branch \texttt{dev}, commit \texttt{153b52e}, 2026, \url{https://github.com/the-omega-institute/trureturing}.
\bibitem{VOS} A.~van Oosterom, J.~Strackee, The solid angle of a plane triangle, \emph{IEEE Trans. Biomed. Eng.} \textbf{BME-30} (1983), 125--126.
% UNVERIFIED (Veselov 1988: not checked against a source)
\bibitem{Veselov88} A.~P.~Veselov, Integrable systems with discrete time, and difference operators, \emph{Funct. Anal. Appl.} \textbf{22} (1988), 83--93.
% UNVERIFIED page range (journal, series, volume and year verified; pages not found)
\bibitem{Volterra} V.~Volterra, Sui fondamenti della teoria delle equazioni differenziali lineari, \emph{Mem. Soc. Ital. Sci. (XL)} (3) \textbf{6} (1887), no.~8.
\bibitem{WZ} E.~Wong, M.~Zakai, On the convergence of ordinary integrals to stochastic integrals, \emph{Ann. Math. Statist.} \textbf{36} (1965), 1560--1564.
\bibitem{Wang52} H.-C.~Wang, Two-point homogeneous spaces, \emph{Ann. of Math.} \textbf{55} (1952), 177--191.
% Verified: doi:10.2307/1969792
\bibitem{Yamabe53} H.~Yamabe, A generalization of a theorem of Gleason, \emph{Ann. of Math.} \textbf{58} (1953), 351--365.
% End page 365 from secondary bibliographies.
% UNVERIFIED title (journal, volume, pages and doi:10.4153/CMB-2011-136-7 verified)
\bibitem{Yang13} D.~Yang, Functional equations and Fourier analysis, \emph{Canad. Math. Bull.} \textbf{56} (2013), 218--224.
% Kuczma 1964, Sharkovskii 1973: bibliographic data verified (IMPAN, EuDML); texts could not be obtained.
\bibitem{Kuczma64} M.~Kuczma, On a characterization of the cosine, \emph{Ann. Polon. Math.} \textbf{16} (1964), 53--57.
\bibitem{Sharkovskii73} A.~N.~Sharkovski\u{\i}, Characterization of the cosine (Russian), \emph{Aequationes Math.} \textbf{9} (1973), 121--128.
\bibitem{EbanksStetkaer15} B.~Ebanks, H.~Stetk\ae r, On Wilson's functional equations, \emph{Aequationes Math.} \textbf{89} (2015), 339--354.
\bibitem{CENSZ95} J.~K.~Chung, B.~R.~Ebanks, C.~T.~Ng, P.~K.~Sahoo, W.~B.~Zeng, On generalized rectangular and rhombic functional equations, \emph{Publ. Math. Debrecen} \textbf{47} (1995), 249--270.
\bibitem{Samelson40} H.~Samelson, \"Uber die Sph\"aren, die als Gruppenr\"aume auftreten, \emph{Comment. Math. Helv.} \textbf{13} (1940/41), 144--155.
% Data from a literature search; contents of Nagy-Strambach 1998 known only from a review.
\bibitem{NS98} P.~T.~Nagy, K.~Strambach, Loops, their cores and symmetric spaces, \emph{Israel J. Math.} \textbf{105} (1998), 285--322.
\bibitem{Bralic80} N.~E.~Bralić, Exact computation of loop averages in two-dimensional Yang--Mills theory, \emph{Phys. Rev. D} \textbf{22} (1980), 3090--3103.
% End page 3103 not verified.
% Content verified only through secondary sources: every partial shelling of a triangulated 2-ball extends to a shelling.
\bibitem{DK78} G.~Danaraj, V.~Klee, Which spheres are shellable?, in \emph{Algorithmic Aspects of Combinatorics}, Ann. Discrete Math. \textbf{2} (1978), 33--52.
\bibitem{Driver89} B.~K.~Driver, YM$_2$: continuum expectations, lattice convergence, and lassos, \emph{Comm. Math. Phys.} \textbf{123} (1989), 575--616.
\bibitem{JP12} J.~Jacod, P.~Protter, \emph{Discretization of Processes}, Stoch. Model. Appl. Probab. \textbf{67}, Springer, 2012.
\bibitem{Levy03} T.~L\'evy, Yang--Mills measure on compact surfaces, \emph{Mem. Amer. Math. Soc.} \textbf{166} (2003), no.~790.
\bibitem{Levy10} T.~L\'evy, Two-dimensional Markovian holonomy fields, \emph{Ast\'erisque} \textbf{329} (2010).
\bibitem{MS93} N.~Mukunda, R.~Simon, Quantum kinematic approach to the geometric phase. I. General formalism, \emph{Ann. Physics} \textbf{228} (1993), 205--268.
\bibitem{SB88} J.~Samuel, R.~Bhandari, General setting for Berry's phase, \emph{Phys. Rev. Lett.} \textbf{60} (1988), 2339--2342.
\bibitem{Stetkaer98} H.~Stetk\ae r, Functional equations on abelian groups with involution, II, \emph{Aequationes Math.} \textbf{55} (1998), 227--240.
% Bibliographic data verified; content (classification of two-point homogeneous spaces) via secondary sources.
\bibitem{Tits55} J.~Tits, \emph{Sur certaines classes d'espaces homog\`enes de groupes de Lie}, Acad. Roy. Belg. Cl. Sci. M\'em. Coll. in-8$^\circ$ \textbf{29} (1955), no.~3.
\bibitem{Loginov01} E.~K.~Loginov, Analytic loops and gauge fields, \emph{Nuclear Phys. B} \textbf{606} (2001), 636--646; arXiv:hep-th/0109206.
% Round 15: title verified in zbMATH (Zbl 0346.17019, "Zentrale einfache Mal'cevsche Algebren", Algebra Logika 15, 235-242 (1976));
% content (central simple non-Lie Malcev algebras are traceless octonions) from Loginov 2001, ref. [18].
\bibitem{Filippov76} V.~T.~Filippov, Central simple Mal'tsev algebras, \emph{Algebra i Logika} \textbf{15} (1976), 235--242.
\bibitem{Sagle62} A.~A.~Sagle, Simple Malcev algebras over fields of characteristic zero, \emph{Pacific J. Math.} \textbf{12} (1962), 1057--1078.
\bibitem{Schafer66} R.~D.~Schafer, \emph{An Introduction to Nonassociative Algebras}, Academic Press, 1966.
\bibitem{Companion} Companion paper, \emph{Axiom $\Omega$ and the origin of quantum mechanics: an origin-free comparison, its clock, and its space of events}, in preparation, 2026.
\bibitem{Rauch75} H.~Rauch, A.~Zeilinger, G.~Badurek, A.~Wilfing, W.~Bauspiess, U.~Bonse, Verification of coherent spinor rotation of fermions, \emph{Phys. Lett. A} \textbf{54} (1975), 425--427.
\bibitem{WCOE75} S.~A.~Werner, R.~Colella, A.~W.~Overhauser, C.~F.~Eagen, Observation of the phase shift of a neutron due to precession in a magnetic field, \emph{Phys. Rev. Lett.} \textbf{35} (1975), 1053--1055.
% UNVERIFIED section number (the holonomy form of Gauss--Bonnet is in Chapter 4)
\bibitem{doCarmo76} M.~P.~do~Carmo, \emph{Differential Geometry of Curves and Surfaces}, Prentice-Hall, 1976.
\bibitem{Rauch78} H.~Rauch, A.~Wilfing, W.~Bauspiess, U.~Bonse, Precise determination of the $4\pi$-periodicity factor of a spinor wave function, \emph{Z. Phys. B} \textbf{29} (1978), 281--284.
\bibitem{BS88} R.~Bhandari, J.~Samuel, Observation of topological phase by use of a laser interferometer, \emph{Phys. Rev. Lett.} \textbf{60} (1988), 1211--1213.
\bibitem{BD87} T.~Bitter, D.~Dubbers, Manifestation of Berry's topological phase in neutron spin rotation, \emph{Phys. Rev. Lett.} \textbf{59} (1987), 251--254.
% Round 10: data of the following seven items checked against publisher or index records (2026-11-20); contents of
% SV00 Ch. 1 and Baez02 Sec. 4.1 checked only via tables of contents and excerpts.
\bibitem{Quackenbush76} R.~W.~Quackenbush, Varieties of Steiner loops and Steiner quasigroups, \emph{Canad. J. Math.} \textbf{28} (1976), 1187--1198.
\bibitem{Hurwitz1898} A.~Hurwitz, \"Uber die Composition der quadratischen Formen von beliebig vielen Variablen, \emph{Nachr. Ges. Wiss. G\"ottingen, Math.-Phys. Kl.} (1898), 309--316.
\bibitem{SV00} T.~A.~Springer, F.~D.~Veldkamp, \emph{Octonions, Jordan Algebras and Exceptional Groups}, Springer Monogr. Math., Springer, Berlin, 2000.
\bibitem{CS03} J.~H.~Conway, D.~A.~Smith, \emph{On Quaternions and Octonions: Their Geometry, Arithmetic, and Symmetry}, A~K~Peters, Natick, MA, 2003.
\bibitem{Baez02} J.~C.~Baez, The octonions, \emph{Bull. Amer. Math. Soc. (N.S.)} \textbf{39} (2002), 145--205.
\bibitem{SVV77} M.~E.~Stoll, A.~J.~Vega, R.~W.~Vaughan, Explicit demonstration of spinor character for a spin-$\tfrac12$ nucleus via NMR interferometry, \emph{Phys. Rev. A} \textbf{16} (1977), 1521--1524.
\bibitem{Potts51} D.~H.~Potts, A note on Green's theorem, \emph{J. London Math. Soc.} \textbf{26} (1951), 302--304.
% Round 13: the six items below: data from the reference list of White16 (full text read); contents only via White16.
\bibitem{White16} B.~White, Lectures on minimal surface theory, IAS/Park City lectures 2013, arXiv:1308.3325v4, 2016.
\bibitem{Morrey48} C.~B.~Morrey, Jr., The problem of Plateau on a Riemannian manifold, \emph{Ann. of Math. (2)} \textbf{49} (1948), 807--851.
\bibitem{Osserman70} R.~Osserman, A proof of the regularity everywhere of the classical solution to Plateau's problem, \emph{Ann. of Math. (2)} \textbf{91} (1970), 550--569.
\bibitem{Gulliver73} R.~D.~Gulliver~II, Regularity of minimizing surfaces of prescribed mean curvature, \emph{Ann. of Math. (2)} \textbf{97} (1973), 275--305.
\bibitem{Hildebrandt69} S.~Hildebrandt, Boundary behavior of minimal surfaces, \emph{Arch. Rational Mech. Anal.} \textbf{35} (1969), 47--82.
\bibitem{HH70} E.~Heinz, S.~Hildebrandt, Some remarks on minimal surfaces in Riemannian manifolds, \emph{Comm. Pure Appl. Math.} \textbf{23} (1970), 371--377.
% Round 14/15 additions: Hofmann58II full text (GDZ scan, OCR) read for (8.30), (9.1)-(9.3); FS15 arXiv full text read (Thm 6);
% BK12 arXiv full text read (Lemma 3.4, Remark (3)).
\bibitem{Hofmann58II} K.~H.~Hofmann, Topologische Loops mit schwachen Assoziativit\"atsforderungen, \emph{Math. Z.} \textbf{70} (1958), 125--155.
\bibitem{FS15} \'A.~Figula, K.~Strambach, Loops as sections in compact Lie groups, arXiv:1502.06911, 2015.
\bibitem{BK12} W.~Bertram, M.~Kinyon, Torsors and ternary Moufang loops arising in projective geometry, arXiv:1206.2222, 2012.
% Round 18: HK15/HK15e: arXiv v4 full text read (Thm 5.4 and erratum), journal data via Crossref; Hatcher Thm 2B.3 read;
% BN09 arXiv full text read (Prop. 1, Thms. 1, 5, 7, 10); Lee13/Lee18 chapter titles and pages via Crossref only (textbook facts);
% Green63, Besse78, Weinstein74, Yang80 data and the Blaschke/Wiedersehen statement from McKay16 (full text read), originals not read.
\bibitem{Hatcher} A.~Hatcher, \emph{Algebraic Topology}, Cambridge Univ. Press, 2002.
\bibitem{HK15} K.~H.~Hofmann, L.~Kramer, Transitive actions of locally compact groups on locally contractible spaces, \emph{J. Reine Angew. Math.} \textbf{702} (2015), 227--243; arXiv:1301.5114.
\bibitem{HK15e} K.~H.~Hofmann, L.~Kramer, Erratum to: Transitive actions of locally compact groups on locally contractible spaces, \emph{J. Reine Angew. Math.} \textbf{702} (2015), 245--246.
\bibitem{Lee13} J.~M.~Lee, \emph{Introduction to Smooth Manifolds}, 2nd ed., Grad. Texts in Math. \textbf{218}, Springer, 2013.
\bibitem{Lee18} J.~M.~Lee, \emph{Introduction to Riemannian Manifolds}, 2nd ed., Grad. Texts in Math. \textbf{176}, Springer, 2018.
\bibitem{BN09} V.~N.~Berestovski\u{\i}, Yu.~G.~Nikonorov, Clifford--Wolf homogeneous Riemannian manifolds, \emph{J. Differential Geom.} \textbf{82} (2009), 467--500; arXiv:0803.4368v1 (theorem numbers cited here are those of the arXiv version).
% DN68: data via Crossref (doi 10.4310/jdg/1214428656) and BN09's reference list; not read.
% Kinyon00: arXiv math/9910111v2 read (abstract, Section 3, Definition 3.1, (3.6)); journal data via search (EuDML/dml.cz).
\bibitem{Kinyon00} M.~K.~Kinyon, Global left loop structures on spheres, \emph{Comment. Math. Univ. Carolin.} \textbf{41} (2000), 325--346; arXiv:math/9910111.
\bibitem{DN68} J.~E.~D'Atri, H.~K.~Nickerson, The existence of special orthonormal frames, \emph{J. Differential Geom.} \textbf{2} (1968), 393--409.
\bibitem{Green63} L.~W.~Green, Auf Wiedersehensfl\"achen, \emph{Ann. of Math. (2)} \textbf{78} (1963), 289--299.
\bibitem{Besse78} A.~L.~Besse, \emph{Manifolds All of Whose Geodesics Are Closed}, Ergeb. Math. Grenzgeb. \textbf{93}, Springer, 1978.
\bibitem{Weinstein74} A.~Weinstein, On the volume of manifolds all of whose geodesics are closed, \emph{J. Differential Geom.} \textbf{9} (1974), 513--517.
\bibitem{Yang80} C.~T.~Yang, Odd-dimensional Wiedersehen manifolds are spheres, \emph{J. Differential Geom.} \textbf{15} (1980), 91--96.
\bibitem{McKay16} B.~McKay, Summary of progress on the Blaschke conjecture, arXiv:1309.1326v5, 2016.
% Part 0 references: data of all eleven entries checked via Crossref (journal items), arXiv API (Duncan-Kileel) and the DLMF page
% (2026-11-20); Spanier is a standard textbook.
\bibitem{m1:Dirac} G.~A.~Dirac, On rigid circuit graphs, \emph{Abh. Math. Sem. Univ. Hamburg} \textbf{25} (1961), 71--76.
\bibitem{m1:DLMF} F.~W.~J.~Olver et al.\ (eds.), \emph{NIST Digital Library of Mathematical Functions}, \S18.5(i), equations 18.5.1--18.5.2.
\bibitem{m1:Grone} R.~Grone, C.~R.~Johnson, E.~M.~S\'a, H.~Wolkowicz, Positive definite completions of partial Hermitian matrices, \emph{Linear Algebra Appl.} \textbf{58} (1984), 109--124.
\bibitem{m1:JordanVonNeumann} P.~Jordan, J.~von Neumann, On inner products in linear, metric spaces, \emph{Ann. of Math. (2)} \textbf{36} (1935), 719--723.
\bibitem{m1:LermanShi} G.~Lerman, Y.~Shi, Robust group synchronization via cycle-edge message passing, \emph{Found. Comput. Math.} \textbf{22} (2022), 1665--1741, arXiv:1912.11347.
\bibitem{m1:Schoenberg} I.~J.~Schoenberg, Remarks to Maurice Fr\'echet's article {Sur la d\'efinition axiomatique d'une classe d'espace distanci\'es vectoriellement applicable sur l'espace de Hilbert}, \emph{Ann. of Math. (2)} \textbf{36} (1935), 724--732.
\bibitem{m1:Singer} A.~Singer, Angular synchronization by eigenvectors and semidefinite programming, \emph{Appl. Comput. Harmon. Anal.} \textbf{30} (2011), 20--36, arXiv:0905.3174.
\bibitem{m1:Spanier} E.~H.~Spanier, \emph{Algebraic Topology}, McGraw--Hill, New York, 1966, Chapter 3.
\bibitem{m1:Zaslavsky} T.~Zaslavsky, Biased graphs. I. Bias, balance, and gains, \emph{J. Combin. Theory Ser. B} \textbf{47} (1989), 32--52.
\bibitem{m1:DuncanKileel} A.~L.~Duncan, J.~Kileel, Higher-order group synchronization, arXiv:2505.21932v2 (2025).
\bibitem{m1:Mermin} N.~D.~Mermin, Simple unified form for the major no-hidden-variables theorems, \emph{Phys. Rev. Lett.} \textbf{65} (1990), 3373--3376.
\end{thebibliography}

\end{document}
